\chapter{Detailed Per-Run Training Graphics}
\label{chap:training}

Extracting summarized means inherently compresses the behavioral nuances observed during multi-million step rollouts. To contextualize the summary metrics, the per-run diagnostics were generated, capturing learning curves, discrete point-evaluations, and application behaviors over global steps.

\section{Diagnostic Panels and Evaluation Trajectories}

Every execution in the 42-run study dumped an immensely rich set of diagnostic panels mapping the core training flow: the primary episode length scaling, dynamic primary metric convergence, and normalized training rewards.

\begin{figure}[H]
    \centering
    \includegraphics[width=0.9\textwidth]{figures/001_p1_ppo_adaptive_fixed_weather_seed0_ent0__diagnostics_panel.png}
    \caption{Typical diagnostics panel extracted from a Point 1 Entropy ablation run (Seed 0, ent=0.0). Notice the aggressive early convergence enabled by clipping entropy.}
    \label{fig:per_run_diag}
\end{figure}

\begin{figure}[H]
    \centering
    \includegraphics[width=0.9\textwidth]{figures/013_p2_a2c_fixed_weather_seed0_blockpen0__checkpoint_eval_curves.png}
    \caption{Checkpoint Evaluation Curves from Point 2 (A2C). This illustrates the extreme volatility of A2C's non-clipped updates recovering dynamically over a strict evaluation cycle.}
    \label{fig:per_run_eval}
\end{figure}

\section{Behavioral Observations in Fertilization}

Specifically, in runs targeting blocked penalties (Point 2) or fertilization targets, detailed application curves mapped exactly \textit{when} the agent elected to deploy compounds. A hallmark of sophisticated policy convergence was the agent learning not only the quantity but the precise temporal distribution mapped directly to crop-stage logic.

\begin{figure}[H]
    \centering
    \includegraphics[width=0.9\textwidth]{figures/013_p2_a2c_fixed_weather_seed0_blockpen0__weekly_npk_behavior.png}
    \caption{Detailed Weekly NPK Application Behavior traced against specific calendar milestones across the episode.}
    \label{fig:per_run_npk}
\end{figure}

\begin{figure}[H]
    \centering
    \includegraphics[width=0.9\textwidth]{figures/013_p2_a2c_fixed_weather_seed0_blockpen0__crop_decision_timeline.png}
    \caption{The hierarchical crop decision timeline demonstrating sequential transitions deployed over the evaluated episode.}
    \label{fig:per_run_crops}
\end{figure}

\section{Key Takeaways from the Diagnostics}
The per-run artifacts fundamentally validate the grouped summary metrics:
\begin{enumerate}
    \item PPO generated smoother, monotonically increasing check-point evaluation curves than A2C. 
    \item Total episode lengths routinely saturated early in fixed-weather sequences, denoting successful structural survival policies. Random weather introduced mid-training "crashes", visible as drops in the \texttt{episode\_length\_vs\_global\_step} plots, corresponding precisely with injected deep drought states.
    \item Agents operating with \texttt{cost\_weight=0.8} (Point 3) visibly sprayed NPK aggressively across timelines, resulting in dense activation graphs on the weekly NPK timelines but lower final episodic yields. 
\end{enumerate}

\section{Complete Per-Run Graphs Portfolio}
The following sections present the comprehensive per-run metrics and graphics generated across both the 113-experiment benchmark and the 42-run ablation study. Due to the volume of runs, each run is allocated its own subsection detailing the respective plots.

\subsection{Run: 001 Fertilization Ppo Adaptive Fixed Weather Years 1000 Seed 0}
\begin{figure}[H]
    \centering
    \includegraphics[width=0.8\textwidth]{figures/per_run/001_fertilization_ppo_adaptive_fixed_weather_years_1000_seed_0__checkpoint_eval_curves.png}
    \caption{001 Fertilization Ppo Adaptive Fixed Weather Years 1000 Seed 0 - Checkpoint Eval Curves}
\end{figure}

\begin{figure}[H]
    \centering
    \includegraphics[width=0.8\textwidth]{figures/per_run/001_fertilization_ppo_adaptive_fixed_weather_years_1000_seed_0__diagnostics_panel.png}
    \caption{001 Fertilization Ppo Adaptive Fixed Weather Years 1000 Seed 0 - Diagnostics Panel}
\end{figure}

\begin{figure}[H]
    \centering
    \includegraphics[width=0.8\textwidth]{figures/per_run/001_fertilization_ppo_adaptive_fixed_weather_years_1000_seed_0__episode_length_vs_global_step.png}
    \caption{001 Fertilization Ppo Adaptive Fixed Weather Years 1000 Seed 0 - Episode Length Vs Global Step}
\end{figure}

\begin{figure}[H]
    \centering
    \includegraphics[width=0.8\textwidth]{figures/per_run/001_fertilization_ppo_adaptive_fixed_weather_years_1000_seed_0__primary_metric_vs_global_step.png}
    \caption{001 Fertilization Ppo Adaptive Fixed Weather Years 1000 Seed 0 - Primary Metric Vs Global Step}
\end{figure}

\begin{figure}[H]
    \centering
    \includegraphics[width=0.8\textwidth]{figures/per_run/001_fertilization_ppo_adaptive_fixed_weather_years_1000_seed_0__training_reward_vs_global_step.png}
    \caption{001 Fertilization Ppo Adaptive Fixed Weather Years 1000 Seed 0 - Training Reward Vs Global Step}
\end{figure}

\clearpage
\subsection{Run: 001 P1 Ppo Adaptive Fixed Weather Seed0 Ent0}
\begin{figure}[H]
    \centering
    \includegraphics[width=0.8\textwidth]{figures/per_run/001_p1_ppo_adaptive_fixed_weather_seed0_ent0__checkpoint_eval_curves.png}
    \caption{001 P1 Ppo Adaptive Fixed Weather Seed0 Ent0 - Checkpoint Eval Curves}
\end{figure}

\begin{figure}[H]
    \centering
    \includegraphics[width=0.8\textwidth]{figures/per_run/001_p1_ppo_adaptive_fixed_weather_seed0_ent0__diagnostics_panel.png}
    \caption{001 P1 Ppo Adaptive Fixed Weather Seed0 Ent0 - Diagnostics Panel}
\end{figure}

\begin{figure}[H]
    \centering
    \includegraphics[width=0.8\textwidth]{figures/per_run/001_p1_ppo_adaptive_fixed_weather_seed0_ent0__episode_length_vs_global_step.png}
    \caption{001 P1 Ppo Adaptive Fixed Weather Seed0 Ent0 - Episode Length Vs Global Step}
\end{figure}

\begin{figure}[H]
    \centering
    \includegraphics[width=0.8\textwidth]{figures/per_run/001_p1_ppo_adaptive_fixed_weather_seed0_ent0__primary_metric_vs_global_step.png}
    \caption{001 P1 Ppo Adaptive Fixed Weather Seed0 Ent0 - Primary Metric Vs Global Step}
\end{figure}

\begin{figure}[H]
    \centering
    \includegraphics[width=0.8\textwidth]{figures/per_run/001_p1_ppo_adaptive_fixed_weather_seed0_ent0__training_reward_vs_global_step.png}
    \caption{001 P1 Ppo Adaptive Fixed Weather Seed0 Ent0 - Training Reward Vs Global Step}
\end{figure}

\clearpage
\subsection{Run: 002 Fertilization Ppo Nonadaptive Fixed Weather Years 1000 Seed 0}
\begin{figure}[H]
    \centering
    \includegraphics[width=0.8\textwidth]{figures/per_run/002_fertilization_ppo_nonadaptive_fixed_weather_years_1000_seed_0__checkpoint_eval_curves.png}
    \caption{002 Fertilization Ppo Nonadaptive Fixed Weather Years 1000 Seed 0 - Checkpoint Eval Curves}
\end{figure}

\begin{figure}[H]
    \centering
    \includegraphics[width=0.8\textwidth]{figures/per_run/002_fertilization_ppo_nonadaptive_fixed_weather_years_1000_seed_0__diagnostics_panel.png}
    \caption{002 Fertilization Ppo Nonadaptive Fixed Weather Years 1000 Seed 0 - Diagnostics Panel}
\end{figure}

\begin{figure}[H]
    \centering
    \includegraphics[width=0.8\textwidth]{figures/per_run/002_fertilization_ppo_nonadaptive_fixed_weather_years_1000_seed_0__episode_length_vs_global_step.png}
    \caption{002 Fertilization Ppo Nonadaptive Fixed Weather Years 1000 Seed 0 - Episode Length Vs Global Step}
\end{figure}

\begin{figure}[H]
    \centering
    \includegraphics[width=0.8\textwidth]{figures/per_run/002_fertilization_ppo_nonadaptive_fixed_weather_years_1000_seed_0__primary_metric_vs_global_step.png}
    \caption{002 Fertilization Ppo Nonadaptive Fixed Weather Years 1000 Seed 0 - Primary Metric Vs Global Step}
\end{figure}

\begin{figure}[H]
    \centering
    \includegraphics[width=0.8\textwidth]{figures/per_run/002_fertilization_ppo_nonadaptive_fixed_weather_years_1000_seed_0__training_reward_vs_global_step.png}
    \caption{002 Fertilization Ppo Nonadaptive Fixed Weather Years 1000 Seed 0 - Training Reward Vs Global Step}
\end{figure}

\clearpage
\subsection{Run: 002 P1 Ppo Adaptive Fixed Weather Seed0 Ent0 01}
\begin{figure}[H]
    \centering
    \includegraphics[width=0.8\textwidth]{figures/per_run/002_p1_ppo_adaptive_fixed_weather_seed0_ent0_01__checkpoint_eval_curves.png}
    \caption{002 P1 Ppo Adaptive Fixed Weather Seed0 Ent0 01 - Checkpoint Eval Curves}
\end{figure}

\begin{figure}[H]
    \centering
    \includegraphics[width=0.8\textwidth]{figures/per_run/002_p1_ppo_adaptive_fixed_weather_seed0_ent0_01__diagnostics_panel.png}
    \caption{002 P1 Ppo Adaptive Fixed Weather Seed0 Ent0 01 - Diagnostics Panel}
\end{figure}

\begin{figure}[H]
    \centering
    \includegraphics[width=0.8\textwidth]{figures/per_run/002_p1_ppo_adaptive_fixed_weather_seed0_ent0_01__episode_length_vs_global_step.png}
    \caption{002 P1 Ppo Adaptive Fixed Weather Seed0 Ent0 01 - Episode Length Vs Global Step}
\end{figure}

\begin{figure}[H]
    \centering
    \includegraphics[width=0.8\textwidth]{figures/per_run/002_p1_ppo_adaptive_fixed_weather_seed0_ent0_01__primary_metric_vs_global_step.png}
    \caption{002 P1 Ppo Adaptive Fixed Weather Seed0 Ent0 01 - Primary Metric Vs Global Step}
\end{figure}

\begin{figure}[H]
    \centering
    \includegraphics[width=0.8\textwidth]{figures/per_run/002_p1_ppo_adaptive_fixed_weather_seed0_ent0_01__training_reward_vs_global_step.png}
    \caption{002 P1 Ppo Adaptive Fixed Weather Seed0 Ent0 01 - Training Reward Vs Global Step}
\end{figure}

\clearpage
\subsection{Run: 003 Fertilization Ppo Adaptive Random Weather Years 1000 Seed 0}
\begin{figure}[H]
    \centering
    \includegraphics[width=0.8\textwidth]{figures/per_run/003_fertilization_ppo_adaptive_random_weather_years_1000_seed_0__checkpoint_eval_curves.png}
    \caption{003 Fertilization Ppo Adaptive Random Weather Years 1000 Seed 0 - Checkpoint Eval Curves}
\end{figure}

\begin{figure}[H]
    \centering
    \includegraphics[width=0.8\textwidth]{figures/per_run/003_fertilization_ppo_adaptive_random_weather_years_1000_seed_0__diagnostics_panel.png}
    \caption{003 Fertilization Ppo Adaptive Random Weather Years 1000 Seed 0 - Diagnostics Panel}
\end{figure}

\begin{figure}[H]
    \centering
    \includegraphics[width=0.8\textwidth]{figures/per_run/003_fertilization_ppo_adaptive_random_weather_years_1000_seed_0__episode_length_vs_global_step.png}
    \caption{003 Fertilization Ppo Adaptive Random Weather Years 1000 Seed 0 - Episode Length Vs Global Step}
\end{figure}

\begin{figure}[H]
    \centering
    \includegraphics[width=0.8\textwidth]{figures/per_run/003_fertilization_ppo_adaptive_random_weather_years_1000_seed_0__primary_metric_vs_global_step.png}
    \caption{003 Fertilization Ppo Adaptive Random Weather Years 1000 Seed 0 - Primary Metric Vs Global Step}
\end{figure}

\begin{figure}[H]
    \centering
    \includegraphics[width=0.8\textwidth]{figures/per_run/003_fertilization_ppo_adaptive_random_weather_years_1000_seed_0__training_reward_vs_global_step.png}
    \caption{003 Fertilization Ppo Adaptive Random Weather Years 1000 Seed 0 - Training Reward Vs Global Step}
\end{figure}

\clearpage
\subsection{Run: 003 P1 Ppo Adaptive Random Weather Seed0 Ent0}
\begin{figure}[H]
    \centering
    \includegraphics[width=0.8\textwidth]{figures/per_run/003_p1_ppo_adaptive_random_weather_seed0_ent0__checkpoint_eval_curves.png}
    \caption{003 P1 Ppo Adaptive Random Weather Seed0 Ent0 - Checkpoint Eval Curves}
\end{figure}

\begin{figure}[H]
    \centering
    \includegraphics[width=0.8\textwidth]{figures/per_run/003_p1_ppo_adaptive_random_weather_seed0_ent0__diagnostics_panel.png}
    \caption{003 P1 Ppo Adaptive Random Weather Seed0 Ent0 - Diagnostics Panel}
\end{figure}

\begin{figure}[H]
    \centering
    \includegraphics[width=0.8\textwidth]{figures/per_run/003_p1_ppo_adaptive_random_weather_seed0_ent0__episode_length_vs_global_step.png}
    \caption{003 P1 Ppo Adaptive Random Weather Seed0 Ent0 - Episode Length Vs Global Step}
\end{figure}

\begin{figure}[H]
    \centering
    \includegraphics[width=0.8\textwidth]{figures/per_run/003_p1_ppo_adaptive_random_weather_seed0_ent0__primary_metric_vs_global_step.png}
    \caption{003 P1 Ppo Adaptive Random Weather Seed0 Ent0 - Primary Metric Vs Global Step}
\end{figure}

\begin{figure}[H]
    \centering
    \includegraphics[width=0.8\textwidth]{figures/per_run/003_p1_ppo_adaptive_random_weather_seed0_ent0__training_reward_vs_global_step.png}
    \caption{003 P1 Ppo Adaptive Random Weather Seed0 Ent0 - Training Reward Vs Global Step}
\end{figure}

\clearpage
\subsection{Run: 004 Fertilization Ppo Nonadaptive Random Weather Years 1000 Seed 0}
\begin{figure}[H]
    \centering
    \includegraphics[width=0.8\textwidth]{figures/per_run/004_fertilization_ppo_nonadaptive_random_weather_years_1000_seed_0__checkpoint_eval_curves.png}
    \caption{004 Fertilization Ppo Nonadaptive Random Weather Years 1000 Seed 0 - Checkpoint Eval Curves}
\end{figure}

\begin{figure}[H]
    \centering
    \includegraphics[width=0.8\textwidth]{figures/per_run/004_fertilization_ppo_nonadaptive_random_weather_years_1000_seed_0__diagnostics_panel.png}
    \caption{004 Fertilization Ppo Nonadaptive Random Weather Years 1000 Seed 0 - Diagnostics Panel}
\end{figure}

\begin{figure}[H]
    \centering
    \includegraphics[width=0.8\textwidth]{figures/per_run/004_fertilization_ppo_nonadaptive_random_weather_years_1000_seed_0__episode_length_vs_global_step.png}
    \caption{004 Fertilization Ppo Nonadaptive Random Weather Years 1000 Seed 0 - Episode Length Vs Global Step}
\end{figure}

\begin{figure}[H]
    \centering
    \includegraphics[width=0.8\textwidth]{figures/per_run/004_fertilization_ppo_nonadaptive_random_weather_years_1000_seed_0__primary_metric_vs_global_step.png}
    \caption{004 Fertilization Ppo Nonadaptive Random Weather Years 1000 Seed 0 - Primary Metric Vs Global Step}
\end{figure}

\begin{figure}[H]
    \centering
    \includegraphics[width=0.8\textwidth]{figures/per_run/004_fertilization_ppo_nonadaptive_random_weather_years_1000_seed_0__training_reward_vs_global_step.png}
    \caption{004 Fertilization Ppo Nonadaptive Random Weather Years 1000 Seed 0 - Training Reward Vs Global Step}
\end{figure}

\clearpage
\subsection{Run: 004 P1 Ppo Adaptive Random Weather Seed0 Ent0 01}
\begin{figure}[H]
    \centering
    \includegraphics[width=0.8\textwidth]{figures/per_run/004_p1_ppo_adaptive_random_weather_seed0_ent0_01__checkpoint_eval_curves.png}
    \caption{004 P1 Ppo Adaptive Random Weather Seed0 Ent0 01 - Checkpoint Eval Curves}
\end{figure}

\begin{figure}[H]
    \centering
    \includegraphics[width=0.8\textwidth]{figures/per_run/004_p1_ppo_adaptive_random_weather_seed0_ent0_01__diagnostics_panel.png}
    \caption{004 P1 Ppo Adaptive Random Weather Seed0 Ent0 01 - Diagnostics Panel}
\end{figure}

\begin{figure}[H]
    \centering
    \includegraphics[width=0.8\textwidth]{figures/per_run/004_p1_ppo_adaptive_random_weather_seed0_ent0_01__episode_length_vs_global_step.png}
    \caption{004 P1 Ppo Adaptive Random Weather Seed0 Ent0 01 - Episode Length Vs Global Step}
\end{figure}

\begin{figure}[H]
    \centering
    \includegraphics[width=0.8\textwidth]{figures/per_run/004_p1_ppo_adaptive_random_weather_seed0_ent0_01__primary_metric_vs_global_step.png}
    \caption{004 P1 Ppo Adaptive Random Weather Seed0 Ent0 01 - Primary Metric Vs Global Step}
\end{figure}

\begin{figure}[H]
    \centering
    \includegraphics[width=0.8\textwidth]{figures/per_run/004_p1_ppo_adaptive_random_weather_seed0_ent0_01__training_reward_vs_global_step.png}
    \caption{004 P1 Ppo Adaptive Random Weather Seed0 Ent0 01 - Training Reward Vs Global Step}
\end{figure}

\clearpage
\subsection{Run: 005 Fertilization Ppo Adaptive Fixed Weather Years 1000 Seed 1}
\begin{figure}[H]
    \centering
    \includegraphics[width=0.8\textwidth]{figures/per_run/005_fertilization_ppo_adaptive_fixed_weather_years_1000_seed_1__checkpoint_eval_curves.png}
    \caption{005 Fertilization Ppo Adaptive Fixed Weather Years 1000 Seed 1 - Checkpoint Eval Curves}
\end{figure}

\begin{figure}[H]
    \centering
    \includegraphics[width=0.8\textwidth]{figures/per_run/005_fertilization_ppo_adaptive_fixed_weather_years_1000_seed_1__diagnostics_panel.png}
    \caption{005 Fertilization Ppo Adaptive Fixed Weather Years 1000 Seed 1 - Diagnostics Panel}
\end{figure}

\begin{figure}[H]
    \centering
    \includegraphics[width=0.8\textwidth]{figures/per_run/005_fertilization_ppo_adaptive_fixed_weather_years_1000_seed_1__episode_length_vs_global_step.png}
    \caption{005 Fertilization Ppo Adaptive Fixed Weather Years 1000 Seed 1 - Episode Length Vs Global Step}
\end{figure}

\begin{figure}[H]
    \centering
    \includegraphics[width=0.8\textwidth]{figures/per_run/005_fertilization_ppo_adaptive_fixed_weather_years_1000_seed_1__primary_metric_vs_global_step.png}
    \caption{005 Fertilization Ppo Adaptive Fixed Weather Years 1000 Seed 1 - Primary Metric Vs Global Step}
\end{figure}

\begin{figure}[H]
    \centering
    \includegraphics[width=0.8\textwidth]{figures/per_run/005_fertilization_ppo_adaptive_fixed_weather_years_1000_seed_1__training_reward_vs_global_step.png}
    \caption{005 Fertilization Ppo Adaptive Fixed Weather Years 1000 Seed 1 - Training Reward Vs Global Step}
\end{figure}

\clearpage
\subsection{Run: 005 P1 Ppo Adaptive Fixed Weather Seed1 Ent0}
\begin{figure}[H]
    \centering
    \includegraphics[width=0.8\textwidth]{figures/per_run/005_p1_ppo_adaptive_fixed_weather_seed1_ent0__checkpoint_eval_curves.png}
    \caption{005 P1 Ppo Adaptive Fixed Weather Seed1 Ent0 - Checkpoint Eval Curves}
\end{figure}

\begin{figure}[H]
    \centering
    \includegraphics[width=0.8\textwidth]{figures/per_run/005_p1_ppo_adaptive_fixed_weather_seed1_ent0__diagnostics_panel.png}
    \caption{005 P1 Ppo Adaptive Fixed Weather Seed1 Ent0 - Diagnostics Panel}
\end{figure}

\begin{figure}[H]
    \centering
    \includegraphics[width=0.8\textwidth]{figures/per_run/005_p1_ppo_adaptive_fixed_weather_seed1_ent0__episode_length_vs_global_step.png}
    \caption{005 P1 Ppo Adaptive Fixed Weather Seed1 Ent0 - Episode Length Vs Global Step}
\end{figure}

\begin{figure}[H]
    \centering
    \includegraphics[width=0.8\textwidth]{figures/per_run/005_p1_ppo_adaptive_fixed_weather_seed1_ent0__primary_metric_vs_global_step.png}
    \caption{005 P1 Ppo Adaptive Fixed Weather Seed1 Ent0 - Primary Metric Vs Global Step}
\end{figure}

\begin{figure}[H]
    \centering
    \includegraphics[width=0.8\textwidth]{figures/per_run/005_p1_ppo_adaptive_fixed_weather_seed1_ent0__training_reward_vs_global_step.png}
    \caption{005 P1 Ppo Adaptive Fixed Weather Seed1 Ent0 - Training Reward Vs Global Step}
\end{figure}

\clearpage
\subsection{Run: 006 Fertilization Ppo Nonadaptive Fixed Weather Years 1000 Seed 1}
\begin{figure}[H]
    \centering
    \includegraphics[width=0.8\textwidth]{figures/per_run/006_fertilization_ppo_nonadaptive_fixed_weather_years_1000_seed_1__checkpoint_eval_curves.png}
    \caption{006 Fertilization Ppo Nonadaptive Fixed Weather Years 1000 Seed 1 - Checkpoint Eval Curves}
\end{figure}

\begin{figure}[H]
    \centering
    \includegraphics[width=0.8\textwidth]{figures/per_run/006_fertilization_ppo_nonadaptive_fixed_weather_years_1000_seed_1__diagnostics_panel.png}
    \caption{006 Fertilization Ppo Nonadaptive Fixed Weather Years 1000 Seed 1 - Diagnostics Panel}
\end{figure}

\begin{figure}[H]
    \centering
    \includegraphics[width=0.8\textwidth]{figures/per_run/006_fertilization_ppo_nonadaptive_fixed_weather_years_1000_seed_1__episode_length_vs_global_step.png}
    \caption{006 Fertilization Ppo Nonadaptive Fixed Weather Years 1000 Seed 1 - Episode Length Vs Global Step}
\end{figure}

\begin{figure}[H]
    \centering
    \includegraphics[width=0.8\textwidth]{figures/per_run/006_fertilization_ppo_nonadaptive_fixed_weather_years_1000_seed_1__primary_metric_vs_global_step.png}
    \caption{006 Fertilization Ppo Nonadaptive Fixed Weather Years 1000 Seed 1 - Primary Metric Vs Global Step}
\end{figure}

\begin{figure}[H]
    \centering
    \includegraphics[width=0.8\textwidth]{figures/per_run/006_fertilization_ppo_nonadaptive_fixed_weather_years_1000_seed_1__training_reward_vs_global_step.png}
    \caption{006 Fertilization Ppo Nonadaptive Fixed Weather Years 1000 Seed 1 - Training Reward Vs Global Step}
\end{figure}

\clearpage
\subsection{Run: 006 P1 Ppo Adaptive Fixed Weather Seed1 Ent0 01}
\begin{figure}[H]
    \centering
    \includegraphics[width=0.8\textwidth]{figures/per_run/006_p1_ppo_adaptive_fixed_weather_seed1_ent0_01__checkpoint_eval_curves.png}
    \caption{006 P1 Ppo Adaptive Fixed Weather Seed1 Ent0 01 - Checkpoint Eval Curves}
\end{figure}

\begin{figure}[H]
    \centering
    \includegraphics[width=0.8\textwidth]{figures/per_run/006_p1_ppo_adaptive_fixed_weather_seed1_ent0_01__diagnostics_panel.png}
    \caption{006 P1 Ppo Adaptive Fixed Weather Seed1 Ent0 01 - Diagnostics Panel}
\end{figure}

\begin{figure}[H]
    \centering
    \includegraphics[width=0.8\textwidth]{figures/per_run/006_p1_ppo_adaptive_fixed_weather_seed1_ent0_01__episode_length_vs_global_step.png}
    \caption{006 P1 Ppo Adaptive Fixed Weather Seed1 Ent0 01 - Episode Length Vs Global Step}
\end{figure}

\begin{figure}[H]
    \centering
    \includegraphics[width=0.8\textwidth]{figures/per_run/006_p1_ppo_adaptive_fixed_weather_seed1_ent0_01__primary_metric_vs_global_step.png}
    \caption{006 P1 Ppo Adaptive Fixed Weather Seed1 Ent0 01 - Primary Metric Vs Global Step}
\end{figure}

\begin{figure}[H]
    \centering
    \includegraphics[width=0.8\textwidth]{figures/per_run/006_p1_ppo_adaptive_fixed_weather_seed1_ent0_01__training_reward_vs_global_step.png}
    \caption{006 P1 Ppo Adaptive Fixed Weather Seed1 Ent0 01 - Training Reward Vs Global Step}
\end{figure}

\clearpage
\subsection{Run: 007 Fertilization Ppo Adaptive Random Weather Years 1000 Seed 1}
\begin{figure}[H]
    \centering
    \includegraphics[width=0.8\textwidth]{figures/per_run/007_fertilization_ppo_adaptive_random_weather_years_1000_seed_1__checkpoint_eval_curves.png}
    \caption{007 Fertilization Ppo Adaptive Random Weather Years 1000 Seed 1 - Checkpoint Eval Curves}
\end{figure}

\begin{figure}[H]
    \centering
    \includegraphics[width=0.8\textwidth]{figures/per_run/007_fertilization_ppo_adaptive_random_weather_years_1000_seed_1__diagnostics_panel.png}
    \caption{007 Fertilization Ppo Adaptive Random Weather Years 1000 Seed 1 - Diagnostics Panel}
\end{figure}

\begin{figure}[H]
    \centering
    \includegraphics[width=0.8\textwidth]{figures/per_run/007_fertilization_ppo_adaptive_random_weather_years_1000_seed_1__episode_length_vs_global_step.png}
    \caption{007 Fertilization Ppo Adaptive Random Weather Years 1000 Seed 1 - Episode Length Vs Global Step}
\end{figure}

\begin{figure}[H]
    \centering
    \includegraphics[width=0.8\textwidth]{figures/per_run/007_fertilization_ppo_adaptive_random_weather_years_1000_seed_1__primary_metric_vs_global_step.png}
    \caption{007 Fertilization Ppo Adaptive Random Weather Years 1000 Seed 1 - Primary Metric Vs Global Step}
\end{figure}

\begin{figure}[H]
    \centering
    \includegraphics[width=0.8\textwidth]{figures/per_run/007_fertilization_ppo_adaptive_random_weather_years_1000_seed_1__training_reward_vs_global_step.png}
    \caption{007 Fertilization Ppo Adaptive Random Weather Years 1000 Seed 1 - Training Reward Vs Global Step}
\end{figure}

\clearpage
\subsection{Run: 007 P1 Ppo Adaptive Fixed Weather Seed2 Ent0}
\begin{figure}[H]
    \centering
    \includegraphics[width=0.8\textwidth]{figures/per_run/007_p1_ppo_adaptive_fixed_weather_seed2_ent0__checkpoint_eval_curves.png}
    \caption{007 P1 Ppo Adaptive Fixed Weather Seed2 Ent0 - Checkpoint Eval Curves}
\end{figure}

\begin{figure}[H]
    \centering
    \includegraphics[width=0.8\textwidth]{figures/per_run/007_p1_ppo_adaptive_fixed_weather_seed2_ent0__diagnostics_panel.png}
    \caption{007 P1 Ppo Adaptive Fixed Weather Seed2 Ent0 - Diagnostics Panel}
\end{figure}

\begin{figure}[H]
    \centering
    \includegraphics[width=0.8\textwidth]{figures/per_run/007_p1_ppo_adaptive_fixed_weather_seed2_ent0__episode_length_vs_global_step.png}
    \caption{007 P1 Ppo Adaptive Fixed Weather Seed2 Ent0 - Episode Length Vs Global Step}
\end{figure}

\begin{figure}[H]
    \centering
    \includegraphics[width=0.8\textwidth]{figures/per_run/007_p1_ppo_adaptive_fixed_weather_seed2_ent0__primary_metric_vs_global_step.png}
    \caption{007 P1 Ppo Adaptive Fixed Weather Seed2 Ent0 - Primary Metric Vs Global Step}
\end{figure}

\begin{figure}[H]
    \centering
    \includegraphics[width=0.8\textwidth]{figures/per_run/007_p1_ppo_adaptive_fixed_weather_seed2_ent0__training_reward_vs_global_step.png}
    \caption{007 P1 Ppo Adaptive Fixed Weather Seed2 Ent0 - Training Reward Vs Global Step}
\end{figure}

\clearpage
\subsection{Run: 008 Fertilization Ppo Nonadaptive Random Weather Years 1000 Seed 1}
\begin{figure}[H]
    \centering
    \includegraphics[width=0.8\textwidth]{figures/per_run/008_fertilization_ppo_nonadaptive_random_weather_years_1000_seed_1__checkpoint_eval_curves.png}
    \caption{008 Fertilization Ppo Nonadaptive Random Weather Years 1000 Seed 1 - Checkpoint Eval Curves}
\end{figure}

\begin{figure}[H]
    \centering
    \includegraphics[width=0.8\textwidth]{figures/per_run/008_fertilization_ppo_nonadaptive_random_weather_years_1000_seed_1__diagnostics_panel.png}
    \caption{008 Fertilization Ppo Nonadaptive Random Weather Years 1000 Seed 1 - Diagnostics Panel}
\end{figure}

\begin{figure}[H]
    \centering
    \includegraphics[width=0.8\textwidth]{figures/per_run/008_fertilization_ppo_nonadaptive_random_weather_years_1000_seed_1__episode_length_vs_global_step.png}
    \caption{008 Fertilization Ppo Nonadaptive Random Weather Years 1000 Seed 1 - Episode Length Vs Global Step}
\end{figure}

\begin{figure}[H]
    \centering
    \includegraphics[width=0.8\textwidth]{figures/per_run/008_fertilization_ppo_nonadaptive_random_weather_years_1000_seed_1__primary_metric_vs_global_step.png}
    \caption{008 Fertilization Ppo Nonadaptive Random Weather Years 1000 Seed 1 - Primary Metric Vs Global Step}
\end{figure}

\begin{figure}[H]
    \centering
    \includegraphics[width=0.8\textwidth]{figures/per_run/008_fertilization_ppo_nonadaptive_random_weather_years_1000_seed_1__training_reward_vs_global_step.png}
    \caption{008 Fertilization Ppo Nonadaptive Random Weather Years 1000 Seed 1 - Training Reward Vs Global Step}
\end{figure}

\clearpage
\subsection{Run: 008 P1 Ppo Adaptive Fixed Weather Seed2 Ent0 01}
\begin{figure}[H]
    \centering
    \includegraphics[width=0.8\textwidth]{figures/per_run/008_p1_ppo_adaptive_fixed_weather_seed2_ent0_01__checkpoint_eval_curves.png}
    \caption{008 P1 Ppo Adaptive Fixed Weather Seed2 Ent0 01 - Checkpoint Eval Curves}
\end{figure}

\begin{figure}[H]
    \centering
    \includegraphics[width=0.8\textwidth]{figures/per_run/008_p1_ppo_adaptive_fixed_weather_seed2_ent0_01__diagnostics_panel.png}
    \caption{008 P1 Ppo Adaptive Fixed Weather Seed2 Ent0 01 - Diagnostics Panel}
\end{figure}

\begin{figure}[H]
    \centering
    \includegraphics[width=0.8\textwidth]{figures/per_run/008_p1_ppo_adaptive_fixed_weather_seed2_ent0_01__episode_length_vs_global_step.png}
    \caption{008 P1 Ppo Adaptive Fixed Weather Seed2 Ent0 01 - Episode Length Vs Global Step}
\end{figure}

\begin{figure}[H]
    \centering
    \includegraphics[width=0.8\textwidth]{figures/per_run/008_p1_ppo_adaptive_fixed_weather_seed2_ent0_01__primary_metric_vs_global_step.png}
    \caption{008 P1 Ppo Adaptive Fixed Weather Seed2 Ent0 01 - Primary Metric Vs Global Step}
\end{figure}

\begin{figure}[H]
    \centering
    \includegraphics[width=0.8\textwidth]{figures/per_run/008_p1_ppo_adaptive_fixed_weather_seed2_ent0_01__training_reward_vs_global_step.png}
    \caption{008 P1 Ppo Adaptive Fixed Weather Seed2 Ent0 01 - Training Reward Vs Global Step}
\end{figure}

\clearpage
\subsection{Run: 009 Fertilization Ppo Adaptive Fixed Weather Years 1000 Seed 2}
\begin{figure}[H]
    \centering
    \includegraphics[width=0.8\textwidth]{figures/per_run/009_fertilization_ppo_adaptive_fixed_weather_years_1000_seed_2__checkpoint_eval_curves.png}
    \caption{009 Fertilization Ppo Adaptive Fixed Weather Years 1000 Seed 2 - Checkpoint Eval Curves}
\end{figure}

\begin{figure}[H]
    \centering
    \includegraphics[width=0.8\textwidth]{figures/per_run/009_fertilization_ppo_adaptive_fixed_weather_years_1000_seed_2__diagnostics_panel.png}
    \caption{009 Fertilization Ppo Adaptive Fixed Weather Years 1000 Seed 2 - Diagnostics Panel}
\end{figure}

\begin{figure}[H]
    \centering
    \includegraphics[width=0.8\textwidth]{figures/per_run/009_fertilization_ppo_adaptive_fixed_weather_years_1000_seed_2__episode_length_vs_global_step.png}
    \caption{009 Fertilization Ppo Adaptive Fixed Weather Years 1000 Seed 2 - Episode Length Vs Global Step}
\end{figure}

\begin{figure}[H]
    \centering
    \includegraphics[width=0.8\textwidth]{figures/per_run/009_fertilization_ppo_adaptive_fixed_weather_years_1000_seed_2__primary_metric_vs_global_step.png}
    \caption{009 Fertilization Ppo Adaptive Fixed Weather Years 1000 Seed 2 - Primary Metric Vs Global Step}
\end{figure}

\begin{figure}[H]
    \centering
    \includegraphics[width=0.8\textwidth]{figures/per_run/009_fertilization_ppo_adaptive_fixed_weather_years_1000_seed_2__training_reward_vs_global_step.png}
    \caption{009 Fertilization Ppo Adaptive Fixed Weather Years 1000 Seed 2 - Training Reward Vs Global Step}
\end{figure}

\clearpage
\subsection{Run: 009 P1 Ppo Adaptive Random Weather Seed1 Ent0}
\begin{figure}[H]
    \centering
    \includegraphics[width=0.8\textwidth]{figures/per_run/009_p1_ppo_adaptive_random_weather_seed1_ent0__checkpoint_eval_curves.png}
    \caption{009 P1 Ppo Adaptive Random Weather Seed1 Ent0 - Checkpoint Eval Curves}
\end{figure}

\begin{figure}[H]
    \centering
    \includegraphics[width=0.8\textwidth]{figures/per_run/009_p1_ppo_adaptive_random_weather_seed1_ent0__diagnostics_panel.png}
    \caption{009 P1 Ppo Adaptive Random Weather Seed1 Ent0 - Diagnostics Panel}
\end{figure}

\begin{figure}[H]
    \centering
    \includegraphics[width=0.8\textwidth]{figures/per_run/009_p1_ppo_adaptive_random_weather_seed1_ent0__episode_length_vs_global_step.png}
    \caption{009 P1 Ppo Adaptive Random Weather Seed1 Ent0 - Episode Length Vs Global Step}
\end{figure}

\begin{figure}[H]
    \centering
    \includegraphics[width=0.8\textwidth]{figures/per_run/009_p1_ppo_adaptive_random_weather_seed1_ent0__primary_metric_vs_global_step.png}
    \caption{009 P1 Ppo Adaptive Random Weather Seed1 Ent0 - Primary Metric Vs Global Step}
\end{figure}

\begin{figure}[H]
    \centering
    \includegraphics[width=0.8\textwidth]{figures/per_run/009_p1_ppo_adaptive_random_weather_seed1_ent0__training_reward_vs_global_step.png}
    \caption{009 P1 Ppo Adaptive Random Weather Seed1 Ent0 - Training Reward Vs Global Step}
\end{figure}

\clearpage
\subsection{Run: 010 Fertilization Ppo Nonadaptive Fixed Weather Years 1000 Seed 2}
\begin{figure}[H]
    \centering
    \includegraphics[width=0.8\textwidth]{figures/per_run/010_fertilization_ppo_nonadaptive_fixed_weather_years_1000_seed_2__checkpoint_eval_curves.png}
    \caption{010 Fertilization Ppo Nonadaptive Fixed Weather Years 1000 Seed 2 - Checkpoint Eval Curves}
\end{figure}

\begin{figure}[H]
    \centering
    \includegraphics[width=0.8\textwidth]{figures/per_run/010_fertilization_ppo_nonadaptive_fixed_weather_years_1000_seed_2__diagnostics_panel.png}
    \caption{010 Fertilization Ppo Nonadaptive Fixed Weather Years 1000 Seed 2 - Diagnostics Panel}
\end{figure}

\begin{figure}[H]
    \centering
    \includegraphics[width=0.8\textwidth]{figures/per_run/010_fertilization_ppo_nonadaptive_fixed_weather_years_1000_seed_2__episode_length_vs_global_step.png}
    \caption{010 Fertilization Ppo Nonadaptive Fixed Weather Years 1000 Seed 2 - Episode Length Vs Global Step}
\end{figure}

\begin{figure}[H]
    \centering
    \includegraphics[width=0.8\textwidth]{figures/per_run/010_fertilization_ppo_nonadaptive_fixed_weather_years_1000_seed_2__primary_metric_vs_global_step.png}
    \caption{010 Fertilization Ppo Nonadaptive Fixed Weather Years 1000 Seed 2 - Primary Metric Vs Global Step}
\end{figure}

\begin{figure}[H]
    \centering
    \includegraphics[width=0.8\textwidth]{figures/per_run/010_fertilization_ppo_nonadaptive_fixed_weather_years_1000_seed_2__training_reward_vs_global_step.png}
    \caption{010 Fertilization Ppo Nonadaptive Fixed Weather Years 1000 Seed 2 - Training Reward Vs Global Step}
\end{figure}

\clearpage
\subsection{Run: 010 P1 Ppo Adaptive Random Weather Seed1 Ent0 01}
\begin{figure}[H]
    \centering
    \includegraphics[width=0.8\textwidth]{figures/per_run/010_p1_ppo_adaptive_random_weather_seed1_ent0_01__checkpoint_eval_curves.png}
    \caption{010 P1 Ppo Adaptive Random Weather Seed1 Ent0 01 - Checkpoint Eval Curves}
\end{figure}

\begin{figure}[H]
    \centering
    \includegraphics[width=0.8\textwidth]{figures/per_run/010_p1_ppo_adaptive_random_weather_seed1_ent0_01__diagnostics_panel.png}
    \caption{010 P1 Ppo Adaptive Random Weather Seed1 Ent0 01 - Diagnostics Panel}
\end{figure}

\begin{figure}[H]
    \centering
    \includegraphics[width=0.8\textwidth]{figures/per_run/010_p1_ppo_adaptive_random_weather_seed1_ent0_01__episode_length_vs_global_step.png}
    \caption{010 P1 Ppo Adaptive Random Weather Seed1 Ent0 01 - Episode Length Vs Global Step}
\end{figure}

\begin{figure}[H]
    \centering
    \includegraphics[width=0.8\textwidth]{figures/per_run/010_p1_ppo_adaptive_random_weather_seed1_ent0_01__primary_metric_vs_global_step.png}
    \caption{010 P1 Ppo Adaptive Random Weather Seed1 Ent0 01 - Primary Metric Vs Global Step}
\end{figure}

\begin{figure}[H]
    \centering
    \includegraphics[width=0.8\textwidth]{figures/per_run/010_p1_ppo_adaptive_random_weather_seed1_ent0_01__training_reward_vs_global_step.png}
    \caption{010 P1 Ppo Adaptive Random Weather Seed1 Ent0 01 - Training Reward Vs Global Step}
\end{figure}

\clearpage
\subsection{Run: 011 Fertilization Ppo Adaptive Random Weather Years 1000 Seed 2}
\begin{figure}[H]
    \centering
    \includegraphics[width=0.8\textwidth]{figures/per_run/011_fertilization_ppo_adaptive_random_weather_years_1000_seed_2__checkpoint_eval_curves.png}
    \caption{011 Fertilization Ppo Adaptive Random Weather Years 1000 Seed 2 - Checkpoint Eval Curves}
\end{figure}

\begin{figure}[H]
    \centering
    \includegraphics[width=0.8\textwidth]{figures/per_run/011_fertilization_ppo_adaptive_random_weather_years_1000_seed_2__diagnostics_panel.png}
    \caption{011 Fertilization Ppo Adaptive Random Weather Years 1000 Seed 2 - Diagnostics Panel}
\end{figure}

\begin{figure}[H]
    \centering
    \includegraphics[width=0.8\textwidth]{figures/per_run/011_fertilization_ppo_adaptive_random_weather_years_1000_seed_2__episode_length_vs_global_step.png}
    \caption{011 Fertilization Ppo Adaptive Random Weather Years 1000 Seed 2 - Episode Length Vs Global Step}
\end{figure}

\begin{figure}[H]
    \centering
    \includegraphics[width=0.8\textwidth]{figures/per_run/011_fertilization_ppo_adaptive_random_weather_years_1000_seed_2__primary_metric_vs_global_step.png}
    \caption{011 Fertilization Ppo Adaptive Random Weather Years 1000 Seed 2 - Primary Metric Vs Global Step}
\end{figure}

\begin{figure}[H]
    \centering
    \includegraphics[width=0.8\textwidth]{figures/per_run/011_fertilization_ppo_adaptive_random_weather_years_1000_seed_2__training_reward_vs_global_step.png}
    \caption{011 Fertilization Ppo Adaptive Random Weather Years 1000 Seed 2 - Training Reward Vs Global Step}
\end{figure}

\clearpage
\subsection{Run: 011 P1 Ppo Adaptive Random Weather Seed2 Ent0}
\begin{figure}[H]
    \centering
    \includegraphics[width=0.8\textwidth]{figures/per_run/011_p1_ppo_adaptive_random_weather_seed2_ent0__checkpoint_eval_curves.png}
    \caption{011 P1 Ppo Adaptive Random Weather Seed2 Ent0 - Checkpoint Eval Curves}
\end{figure}

\begin{figure}[H]
    \centering
    \includegraphics[width=0.8\textwidth]{figures/per_run/011_p1_ppo_adaptive_random_weather_seed2_ent0__diagnostics_panel.png}
    \caption{011 P1 Ppo Adaptive Random Weather Seed2 Ent0 - Diagnostics Panel}
\end{figure}

\begin{figure}[H]
    \centering
    \includegraphics[width=0.8\textwidth]{figures/per_run/011_p1_ppo_adaptive_random_weather_seed2_ent0__episode_length_vs_global_step.png}
    \caption{011 P1 Ppo Adaptive Random Weather Seed2 Ent0 - Episode Length Vs Global Step}
\end{figure}

\begin{figure}[H]
    \centering
    \includegraphics[width=0.8\textwidth]{figures/per_run/011_p1_ppo_adaptive_random_weather_seed2_ent0__primary_metric_vs_global_step.png}
    \caption{011 P1 Ppo Adaptive Random Weather Seed2 Ent0 - Primary Metric Vs Global Step}
\end{figure}

\begin{figure}[H]
    \centering
    \includegraphics[width=0.8\textwidth]{figures/per_run/011_p1_ppo_adaptive_random_weather_seed2_ent0__training_reward_vs_global_step.png}
    \caption{011 P1 Ppo Adaptive Random Weather Seed2 Ent0 - Training Reward Vs Global Step}
\end{figure}

\clearpage
\subsection{Run: 012 Fertilization Ppo Nonadaptive Random Weather Years 1000 Seed 2}
\begin{figure}[H]
    \centering
    \includegraphics[width=0.8\textwidth]{figures/per_run/012_fertilization_ppo_nonadaptive_random_weather_years_1000_seed_2__checkpoint_eval_curves.png}
    \caption{012 Fertilization Ppo Nonadaptive Random Weather Years 1000 Seed 2 - Checkpoint Eval Curves}
\end{figure}

\begin{figure}[H]
    \centering
    \includegraphics[width=0.8\textwidth]{figures/per_run/012_fertilization_ppo_nonadaptive_random_weather_years_1000_seed_2__diagnostics_panel.png}
    \caption{012 Fertilization Ppo Nonadaptive Random Weather Years 1000 Seed 2 - Diagnostics Panel}
\end{figure}

\begin{figure}[H]
    \centering
    \includegraphics[width=0.8\textwidth]{figures/per_run/012_fertilization_ppo_nonadaptive_random_weather_years_1000_seed_2__episode_length_vs_global_step.png}
    \caption{012 Fertilization Ppo Nonadaptive Random Weather Years 1000 Seed 2 - Episode Length Vs Global Step}
\end{figure}

\begin{figure}[H]
    \centering
    \includegraphics[width=0.8\textwidth]{figures/per_run/012_fertilization_ppo_nonadaptive_random_weather_years_1000_seed_2__primary_metric_vs_global_step.png}
    \caption{012 Fertilization Ppo Nonadaptive Random Weather Years 1000 Seed 2 - Primary Metric Vs Global Step}
\end{figure}

\begin{figure}[H]
    \centering
    \includegraphics[width=0.8\textwidth]{figures/per_run/012_fertilization_ppo_nonadaptive_random_weather_years_1000_seed_2__training_reward_vs_global_step.png}
    \caption{012 Fertilization Ppo Nonadaptive Random Weather Years 1000 Seed 2 - Training Reward Vs Global Step}
\end{figure}

\clearpage
\subsection{Run: 012 P1 Ppo Adaptive Random Weather Seed2 Ent0 01}
\begin{figure}[H]
    \centering
    \includegraphics[width=0.8\textwidth]{figures/per_run/012_p1_ppo_adaptive_random_weather_seed2_ent0_01__checkpoint_eval_curves.png}
    \caption{012 P1 Ppo Adaptive Random Weather Seed2 Ent0 01 - Checkpoint Eval Curves}
\end{figure}

\begin{figure}[H]
    \centering
    \includegraphics[width=0.8\textwidth]{figures/per_run/012_p1_ppo_adaptive_random_weather_seed2_ent0_01__diagnostics_panel.png}
    \caption{012 P1 Ppo Adaptive Random Weather Seed2 Ent0 01 - Diagnostics Panel}
\end{figure}

\begin{figure}[H]
    \centering
    \includegraphics[width=0.8\textwidth]{figures/per_run/012_p1_ppo_adaptive_random_weather_seed2_ent0_01__episode_length_vs_global_step.png}
    \caption{012 P1 Ppo Adaptive Random Weather Seed2 Ent0 01 - Episode Length Vs Global Step}
\end{figure}

\begin{figure}[H]
    \centering
    \includegraphics[width=0.8\textwidth]{figures/per_run/012_p1_ppo_adaptive_random_weather_seed2_ent0_01__primary_metric_vs_global_step.png}
    \caption{012 P1 Ppo Adaptive Random Weather Seed2 Ent0 01 - Primary Metric Vs Global Step}
\end{figure}

\begin{figure}[H]
    \centering
    \includegraphics[width=0.8\textwidth]{figures/per_run/012_p1_ppo_adaptive_random_weather_seed2_ent0_01__training_reward_vs_global_step.png}
    \caption{012 P1 Ppo Adaptive Random Weather Seed2 Ent0 01 - Training Reward Vs Global Step}
\end{figure}

\clearpage
\subsection{Run: 013 Fertilization Ppo Adaptive Fixed Weather Years 3000 Seed 0}
\begin{figure}[H]
    \centering
    \includegraphics[width=0.8\textwidth]{figures/per_run/013_fertilization_ppo_adaptive_fixed_weather_years_3000_seed_0__checkpoint_eval_curves.png}
    \caption{013 Fertilization Ppo Adaptive Fixed Weather Years 3000 Seed 0 - Checkpoint Eval Curves}
\end{figure}

\begin{figure}[H]
    \centering
    \includegraphics[width=0.8\textwidth]{figures/per_run/013_fertilization_ppo_adaptive_fixed_weather_years_3000_seed_0__diagnostics_panel.png}
    \caption{013 Fertilization Ppo Adaptive Fixed Weather Years 3000 Seed 0 - Diagnostics Panel}
\end{figure}

\begin{figure}[H]
    \centering
    \includegraphics[width=0.8\textwidth]{figures/per_run/013_fertilization_ppo_adaptive_fixed_weather_years_3000_seed_0__episode_length_vs_global_step.png}
    \caption{013 Fertilization Ppo Adaptive Fixed Weather Years 3000 Seed 0 - Episode Length Vs Global Step}
\end{figure}

\begin{figure}[H]
    \centering
    \includegraphics[width=0.8\textwidth]{figures/per_run/013_fertilization_ppo_adaptive_fixed_weather_years_3000_seed_0__primary_metric_vs_global_step.png}
    \caption{013 Fertilization Ppo Adaptive Fixed Weather Years 3000 Seed 0 - Primary Metric Vs Global Step}
\end{figure}

\begin{figure}[H]
    \centering
    \includegraphics[width=0.8\textwidth]{figures/per_run/013_fertilization_ppo_adaptive_fixed_weather_years_3000_seed_0__training_reward_vs_global_step.png}
    \caption{013 Fertilization Ppo Adaptive Fixed Weather Years 3000 Seed 0 - Training Reward Vs Global Step}
\end{figure}

\clearpage
\subsection{Run: 013 P2 A2C Fixed Weather Seed0 Blockpen0}
\begin{figure}[H]
    \centering
    \includegraphics[width=0.8\textwidth]{figures/per_run/013_p2_a2c_fixed_weather_seed0_blockpen0__blocked_cost_summary.png}
    \caption{013 P2 A2C Fixed Weather Seed0 Blockpen0 - Blocked Cost Summary}
\end{figure}

\begin{figure}[H]
    \centering
    \includegraphics[width=0.8\textwidth]{figures/per_run/013_p2_a2c_fixed_weather_seed0_blockpen0__checkpoint_eval_curves.png}
    \caption{013 P2 A2C Fixed Weather Seed0 Blockpen0 - Checkpoint Eval Curves}
\end{figure}

\begin{figure}[H]
    \centering
    \includegraphics[width=0.8\textwidth]{figures/per_run/013_p2_a2c_fixed_weather_seed0_blockpen0__compliance_summary.png}
    \caption{013 P2 A2C Fixed Weather Seed0 Blockpen0 - Compliance Summary}
\end{figure}

\begin{figure}[H]
    \centering
    \includegraphics[width=0.8\textwidth]{figures/per_run/013_p2_a2c_fixed_weather_seed0_blockpen0__crop_decision_timeline.png}
    \caption{013 P2 A2C Fixed Weather Seed0 Blockpen0 - Crop Decision Timeline}
\end{figure}

\begin{figure}[H]
    \centering
    \includegraphics[width=0.8\textwidth]{figures/per_run/013_p2_a2c_fixed_weather_seed0_blockpen0__diagnostics_panel.png}
    \caption{013 P2 A2C Fixed Weather Seed0 Blockpen0 - Diagnostics Panel}
\end{figure}

\begin{figure}[H]
    \centering
    \includegraphics[width=0.8\textwidth]{figures/per_run/013_p2_a2c_fixed_weather_seed0_blockpen0__episode_length_vs_global_step.png}
    \caption{013 P2 A2C Fixed Weather Seed0 Blockpen0 - Episode Length Vs Global Step}
\end{figure}

\begin{figure}[H]
    \centering
    \includegraphics[width=0.8\textwidth]{figures/per_run/013_p2_a2c_fixed_weather_seed0_blockpen0__primary_metric_vs_global_step.png}
    \caption{013 P2 A2C Fixed Weather Seed0 Blockpen0 - Primary Metric Vs Global Step}
\end{figure}

\begin{figure}[H]
    \centering
    \includegraphics[width=0.8\textwidth]{figures/per_run/013_p2_a2c_fixed_weather_seed0_blockpen0__training_reward_vs_global_step.png}
    \caption{013 P2 A2C Fixed Weather Seed0 Blockpen0 - Training Reward Vs Global Step}
\end{figure}

\begin{figure}[H]
    \centering
    \includegraphics[width=0.8\textwidth]{figures/per_run/013_p2_a2c_fixed_weather_seed0_blockpen0__weekly_npk_behavior.png}
    \caption{013 P2 A2C Fixed Weather Seed0 Blockpen0 - Weekly Npk Behavior}
\end{figure}

\clearpage
\subsection{Run: 014 Fertilization Ppo Nonadaptive Fixed Weather Years 3000 Seed 0}
\begin{figure}[H]
    \centering
    \includegraphics[width=0.8\textwidth]{figures/per_run/014_fertilization_ppo_nonadaptive_fixed_weather_years_3000_seed_0__checkpoint_eval_curves.png}
    \caption{014 Fertilization Ppo Nonadaptive Fixed Weather Years 3000 Seed 0 - Checkpoint Eval Curves}
\end{figure}

\begin{figure}[H]
    \centering
    \includegraphics[width=0.8\textwidth]{figures/per_run/014_fertilization_ppo_nonadaptive_fixed_weather_years_3000_seed_0__diagnostics_panel.png}
    \caption{014 Fertilization Ppo Nonadaptive Fixed Weather Years 3000 Seed 0 - Diagnostics Panel}
\end{figure}

\begin{figure}[H]
    \centering
    \includegraphics[width=0.8\textwidth]{figures/per_run/014_fertilization_ppo_nonadaptive_fixed_weather_years_3000_seed_0__episode_length_vs_global_step.png}
    \caption{014 Fertilization Ppo Nonadaptive Fixed Weather Years 3000 Seed 0 - Episode Length Vs Global Step}
\end{figure}

\begin{figure}[H]
    \centering
    \includegraphics[width=0.8\textwidth]{figures/per_run/014_fertilization_ppo_nonadaptive_fixed_weather_years_3000_seed_0__primary_metric_vs_global_step.png}
    \caption{014 Fertilization Ppo Nonadaptive Fixed Weather Years 3000 Seed 0 - Primary Metric Vs Global Step}
\end{figure}

\begin{figure}[H]
    \centering
    \includegraphics[width=0.8\textwidth]{figures/per_run/014_fertilization_ppo_nonadaptive_fixed_weather_years_3000_seed_0__training_reward_vs_global_step.png}
    \caption{014 Fertilization Ppo Nonadaptive Fixed Weather Years 3000 Seed 0 - Training Reward Vs Global Step}
\end{figure}

\clearpage
\subsection{Run: 014 P2 A2C Fixed Weather Seed0 Blockpen0 02}
\begin{figure}[H]
    \centering
    \includegraphics[width=0.8\textwidth]{figures/per_run/014_p2_a2c_fixed_weather_seed0_blockpen0_02__blocked_cost_summary.png}
    \caption{014 P2 A2C Fixed Weather Seed0 Blockpen0 02 - Blocked Cost Summary}
\end{figure}

\begin{figure}[H]
    \centering
    \includegraphics[width=0.8\textwidth]{figures/per_run/014_p2_a2c_fixed_weather_seed0_blockpen0_02__checkpoint_eval_curves.png}
    \caption{014 P2 A2C Fixed Weather Seed0 Blockpen0 02 - Checkpoint Eval Curves}
\end{figure}

\begin{figure}[H]
    \centering
    \includegraphics[width=0.8\textwidth]{figures/per_run/014_p2_a2c_fixed_weather_seed0_blockpen0_02__compliance_summary.png}
    \caption{014 P2 A2C Fixed Weather Seed0 Blockpen0 02 - Compliance Summary}
\end{figure}

\begin{figure}[H]
    \centering
    \includegraphics[width=0.8\textwidth]{figures/per_run/014_p2_a2c_fixed_weather_seed0_blockpen0_02__crop_decision_timeline.png}
    \caption{014 P2 A2C Fixed Weather Seed0 Blockpen0 02 - Crop Decision Timeline}
\end{figure}

\begin{figure}[H]
    \centering
    \includegraphics[width=0.8\textwidth]{figures/per_run/014_p2_a2c_fixed_weather_seed0_blockpen0_02__diagnostics_panel.png}
    \caption{014 P2 A2C Fixed Weather Seed0 Blockpen0 02 - Diagnostics Panel}
\end{figure}

\begin{figure}[H]
    \centering
    \includegraphics[width=0.8\textwidth]{figures/per_run/014_p2_a2c_fixed_weather_seed0_blockpen0_02__episode_length_vs_global_step.png}
    \caption{014 P2 A2C Fixed Weather Seed0 Blockpen0 02 - Episode Length Vs Global Step}
\end{figure}

\begin{figure}[H]
    \centering
    \includegraphics[width=0.8\textwidth]{figures/per_run/014_p2_a2c_fixed_weather_seed0_blockpen0_02__primary_metric_vs_global_step.png}
    \caption{014 P2 A2C Fixed Weather Seed0 Blockpen0 02 - Primary Metric Vs Global Step}
\end{figure}

\begin{figure}[H]
    \centering
    \includegraphics[width=0.8\textwidth]{figures/per_run/014_p2_a2c_fixed_weather_seed0_blockpen0_02__training_reward_vs_global_step.png}
    \caption{014 P2 A2C Fixed Weather Seed0 Blockpen0 02 - Training Reward Vs Global Step}
\end{figure}

\begin{figure}[H]
    \centering
    \includegraphics[width=0.8\textwidth]{figures/per_run/014_p2_a2c_fixed_weather_seed0_blockpen0_02__weekly_npk_behavior.png}
    \caption{014 P2 A2C Fixed Weather Seed0 Blockpen0 02 - Weekly Npk Behavior}
\end{figure}

\clearpage
\subsection{Run: 015 Fertilization Ppo Adaptive Random Weather Years 3000 Seed 0}
\begin{figure}[H]
    \centering
    \includegraphics[width=0.8\textwidth]{figures/per_run/015_fertilization_ppo_adaptive_random_weather_years_3000_seed_0__checkpoint_eval_curves.png}
    \caption{015 Fertilization Ppo Adaptive Random Weather Years 3000 Seed 0 - Checkpoint Eval Curves}
\end{figure}

\begin{figure}[H]
    \centering
    \includegraphics[width=0.8\textwidth]{figures/per_run/015_fertilization_ppo_adaptive_random_weather_years_3000_seed_0__diagnostics_panel.png}
    \caption{015 Fertilization Ppo Adaptive Random Weather Years 3000 Seed 0 - Diagnostics Panel}
\end{figure}

\begin{figure}[H]
    \centering
    \includegraphics[width=0.8\textwidth]{figures/per_run/015_fertilization_ppo_adaptive_random_weather_years_3000_seed_0__episode_length_vs_global_step.png}
    \caption{015 Fertilization Ppo Adaptive Random Weather Years 3000 Seed 0 - Episode Length Vs Global Step}
\end{figure}

\begin{figure}[H]
    \centering
    \includegraphics[width=0.8\textwidth]{figures/per_run/015_fertilization_ppo_adaptive_random_weather_years_3000_seed_0__primary_metric_vs_global_step.png}
    \caption{015 Fertilization Ppo Adaptive Random Weather Years 3000 Seed 0 - Primary Metric Vs Global Step}
\end{figure}

\begin{figure}[H]
    \centering
    \includegraphics[width=0.8\textwidth]{figures/per_run/015_fertilization_ppo_adaptive_random_weather_years_3000_seed_0__training_reward_vs_global_step.png}
    \caption{015 Fertilization Ppo Adaptive Random Weather Years 3000 Seed 0 - Training Reward Vs Global Step}
\end{figure}

\clearpage
\subsection{Run: 015 P2 A2C Fixed Weather Seed0 Blockpen0 05}
\begin{figure}[H]
    \centering
    \includegraphics[width=0.8\textwidth]{figures/per_run/015_p2_a2c_fixed_weather_seed0_blockpen0_05__blocked_cost_summary.png}
    \caption{015 P2 A2C Fixed Weather Seed0 Blockpen0 05 - Blocked Cost Summary}
\end{figure}

\begin{figure}[H]
    \centering
    \includegraphics[width=0.8\textwidth]{figures/per_run/015_p2_a2c_fixed_weather_seed0_blockpen0_05__checkpoint_eval_curves.png}
    \caption{015 P2 A2C Fixed Weather Seed0 Blockpen0 05 - Checkpoint Eval Curves}
\end{figure}

\begin{figure}[H]
    \centering
    \includegraphics[width=0.8\textwidth]{figures/per_run/015_p2_a2c_fixed_weather_seed0_blockpen0_05__compliance_summary.png}
    \caption{015 P2 A2C Fixed Weather Seed0 Blockpen0 05 - Compliance Summary}
\end{figure}

\begin{figure}[H]
    \centering
    \includegraphics[width=0.8\textwidth]{figures/per_run/015_p2_a2c_fixed_weather_seed0_blockpen0_05__crop_decision_timeline.png}
    \caption{015 P2 A2C Fixed Weather Seed0 Blockpen0 05 - Crop Decision Timeline}
\end{figure}

\begin{figure}[H]
    \centering
    \includegraphics[width=0.8\textwidth]{figures/per_run/015_p2_a2c_fixed_weather_seed0_blockpen0_05__diagnostics_panel.png}
    \caption{015 P2 A2C Fixed Weather Seed0 Blockpen0 05 - Diagnostics Panel}
\end{figure}

\begin{figure}[H]
    \centering
    \includegraphics[width=0.8\textwidth]{figures/per_run/015_p2_a2c_fixed_weather_seed0_blockpen0_05__episode_length_vs_global_step.png}
    \caption{015 P2 A2C Fixed Weather Seed0 Blockpen0 05 - Episode Length Vs Global Step}
\end{figure}

\begin{figure}[H]
    \centering
    \includegraphics[width=0.8\textwidth]{figures/per_run/015_p2_a2c_fixed_weather_seed0_blockpen0_05__primary_metric_vs_global_step.png}
    \caption{015 P2 A2C Fixed Weather Seed0 Blockpen0 05 - Primary Metric Vs Global Step}
\end{figure}

\begin{figure}[H]
    \centering
    \includegraphics[width=0.8\textwidth]{figures/per_run/015_p2_a2c_fixed_weather_seed0_blockpen0_05__training_reward_vs_global_step.png}
    \caption{015 P2 A2C Fixed Weather Seed0 Blockpen0 05 - Training Reward Vs Global Step}
\end{figure}

\begin{figure}[H]
    \centering
    \includegraphics[width=0.8\textwidth]{figures/per_run/015_p2_a2c_fixed_weather_seed0_blockpen0_05__weekly_npk_behavior.png}
    \caption{015 P2 A2C Fixed Weather Seed0 Blockpen0 05 - Weekly Npk Behavior}
\end{figure}

\clearpage
\subsection{Run: 016 Fertilization Ppo Nonadaptive Random Weather Years 3000 Seed 0}
\begin{figure}[H]
    \centering
    \includegraphics[width=0.8\textwidth]{figures/per_run/016_fertilization_ppo_nonadaptive_random_weather_years_3000_seed_0__checkpoint_eval_curves.png}
    \caption{016 Fertilization Ppo Nonadaptive Random Weather Years 3000 Seed 0 - Checkpoint Eval Curves}
\end{figure}

\begin{figure}[H]
    \centering
    \includegraphics[width=0.8\textwidth]{figures/per_run/016_fertilization_ppo_nonadaptive_random_weather_years_3000_seed_0__diagnostics_panel.png}
    \caption{016 Fertilization Ppo Nonadaptive Random Weather Years 3000 Seed 0 - Diagnostics Panel}
\end{figure}

\begin{figure}[H]
    \centering
    \includegraphics[width=0.8\textwidth]{figures/per_run/016_fertilization_ppo_nonadaptive_random_weather_years_3000_seed_0__episode_length_vs_global_step.png}
    \caption{016 Fertilization Ppo Nonadaptive Random Weather Years 3000 Seed 0 - Episode Length Vs Global Step}
\end{figure}

\begin{figure}[H]
    \centering
    \includegraphics[width=0.8\textwidth]{figures/per_run/016_fertilization_ppo_nonadaptive_random_weather_years_3000_seed_0__primary_metric_vs_global_step.png}
    \caption{016 Fertilization Ppo Nonadaptive Random Weather Years 3000 Seed 0 - Primary Metric Vs Global Step}
\end{figure}

\begin{figure}[H]
    \centering
    \includegraphics[width=0.8\textwidth]{figures/per_run/016_fertilization_ppo_nonadaptive_random_weather_years_3000_seed_0__training_reward_vs_global_step.png}
    \caption{016 Fertilization Ppo Nonadaptive Random Weather Years 3000 Seed 0 - Training Reward Vs Global Step}
\end{figure}

\clearpage
\subsection{Run: 016 P2 A2C Random Weather Seed0 Blockpen0}
\begin{figure}[H]
    \centering
    \includegraphics[width=0.8\textwidth]{figures/per_run/016_p2_a2c_random_weather_seed0_blockpen0__blocked_cost_summary.png}
    \caption{016 P2 A2C Random Weather Seed0 Blockpen0 - Blocked Cost Summary}
\end{figure}

\begin{figure}[H]
    \centering
    \includegraphics[width=0.8\textwidth]{figures/per_run/016_p2_a2c_random_weather_seed0_blockpen0__checkpoint_eval_curves.png}
    \caption{016 P2 A2C Random Weather Seed0 Blockpen0 - Checkpoint Eval Curves}
\end{figure}

\begin{figure}[H]
    \centering
    \includegraphics[width=0.8\textwidth]{figures/per_run/016_p2_a2c_random_weather_seed0_blockpen0__compliance_summary.png}
    \caption{016 P2 A2C Random Weather Seed0 Blockpen0 - Compliance Summary}
\end{figure}

\begin{figure}[H]
    \centering
    \includegraphics[width=0.8\textwidth]{figures/per_run/016_p2_a2c_random_weather_seed0_blockpen0__crop_decision_timeline.png}
    \caption{016 P2 A2C Random Weather Seed0 Blockpen0 - Crop Decision Timeline}
\end{figure}

\begin{figure}[H]
    \centering
    \includegraphics[width=0.8\textwidth]{figures/per_run/016_p2_a2c_random_weather_seed0_blockpen0__diagnostics_panel.png}
    \caption{016 P2 A2C Random Weather Seed0 Blockpen0 - Diagnostics Panel}
\end{figure}

\begin{figure}[H]
    \centering
    \includegraphics[width=0.8\textwidth]{figures/per_run/016_p2_a2c_random_weather_seed0_blockpen0__episode_length_vs_global_step.png}
    \caption{016 P2 A2C Random Weather Seed0 Blockpen0 - Episode Length Vs Global Step}
\end{figure}

\begin{figure}[H]
    \centering
    \includegraphics[width=0.8\textwidth]{figures/per_run/016_p2_a2c_random_weather_seed0_blockpen0__primary_metric_vs_global_step.png}
    \caption{016 P2 A2C Random Weather Seed0 Blockpen0 - Primary Metric Vs Global Step}
\end{figure}

\begin{figure}[H]
    \centering
    \includegraphics[width=0.8\textwidth]{figures/per_run/016_p2_a2c_random_weather_seed0_blockpen0__training_reward_vs_global_step.png}
    \caption{016 P2 A2C Random Weather Seed0 Blockpen0 - Training Reward Vs Global Step}
\end{figure}

\begin{figure}[H]
    \centering
    \includegraphics[width=0.8\textwidth]{figures/per_run/016_p2_a2c_random_weather_seed0_blockpen0__weekly_npk_behavior.png}
    \caption{016 P2 A2C Random Weather Seed0 Blockpen0 - Weekly Npk Behavior}
\end{figure}

\clearpage
\subsection{Run: 017 Fertilization Ppo Adaptive Fixed Weather Years 3000 Seed 1}
\begin{figure}[H]
    \centering
    \includegraphics[width=0.8\textwidth]{figures/per_run/017_fertilization_ppo_adaptive_fixed_weather_years_3000_seed_1__checkpoint_eval_curves.png}
    \caption{017 Fertilization Ppo Adaptive Fixed Weather Years 3000 Seed 1 - Checkpoint Eval Curves}
\end{figure}

\begin{figure}[H]
    \centering
    \includegraphics[width=0.8\textwidth]{figures/per_run/017_fertilization_ppo_adaptive_fixed_weather_years_3000_seed_1__diagnostics_panel.png}
    \caption{017 Fertilization Ppo Adaptive Fixed Weather Years 3000 Seed 1 - Diagnostics Panel}
\end{figure}

\begin{figure}[H]
    \centering
    \includegraphics[width=0.8\textwidth]{figures/per_run/017_fertilization_ppo_adaptive_fixed_weather_years_3000_seed_1__episode_length_vs_global_step.png}
    \caption{017 Fertilization Ppo Adaptive Fixed Weather Years 3000 Seed 1 - Episode Length Vs Global Step}
\end{figure}

\begin{figure}[H]
    \centering
    \includegraphics[width=0.8\textwidth]{figures/per_run/017_fertilization_ppo_adaptive_fixed_weather_years_3000_seed_1__primary_metric_vs_global_step.png}
    \caption{017 Fertilization Ppo Adaptive Fixed Weather Years 3000 Seed 1 - Primary Metric Vs Global Step}
\end{figure}

\begin{figure}[H]
    \centering
    \includegraphics[width=0.8\textwidth]{figures/per_run/017_fertilization_ppo_adaptive_fixed_weather_years_3000_seed_1__training_reward_vs_global_step.png}
    \caption{017 Fertilization Ppo Adaptive Fixed Weather Years 3000 Seed 1 - Training Reward Vs Global Step}
\end{figure}

\clearpage
\subsection{Run: 017 P2 A2C Random Weather Seed0 Blockpen0 02}
\begin{figure}[H]
    \centering
    \includegraphics[width=0.8\textwidth]{figures/per_run/017_p2_a2c_random_weather_seed0_blockpen0_02__blocked_cost_summary.png}
    \caption{017 P2 A2C Random Weather Seed0 Blockpen0 02 - Blocked Cost Summary}
\end{figure}

\begin{figure}[H]
    \centering
    \includegraphics[width=0.8\textwidth]{figures/per_run/017_p2_a2c_random_weather_seed0_blockpen0_02__checkpoint_eval_curves.png}
    \caption{017 P2 A2C Random Weather Seed0 Blockpen0 02 - Checkpoint Eval Curves}
\end{figure}

\begin{figure}[H]
    \centering
    \includegraphics[width=0.8\textwidth]{figures/per_run/017_p2_a2c_random_weather_seed0_blockpen0_02__compliance_summary.png}
    \caption{017 P2 A2C Random Weather Seed0 Blockpen0 02 - Compliance Summary}
\end{figure}

\begin{figure}[H]
    \centering
    \includegraphics[width=0.8\textwidth]{figures/per_run/017_p2_a2c_random_weather_seed0_blockpen0_02__crop_decision_timeline.png}
    \caption{017 P2 A2C Random Weather Seed0 Blockpen0 02 - Crop Decision Timeline}
\end{figure}

\begin{figure}[H]
    \centering
    \includegraphics[width=0.8\textwidth]{figures/per_run/017_p2_a2c_random_weather_seed0_blockpen0_02__diagnostics_panel.png}
    \caption{017 P2 A2C Random Weather Seed0 Blockpen0 02 - Diagnostics Panel}
\end{figure}

\begin{figure}[H]
    \centering
    \includegraphics[width=0.8\textwidth]{figures/per_run/017_p2_a2c_random_weather_seed0_blockpen0_02__episode_length_vs_global_step.png}
    \caption{017 P2 A2C Random Weather Seed0 Blockpen0 02 - Episode Length Vs Global Step}
\end{figure}

\begin{figure}[H]
    \centering
    \includegraphics[width=0.8\textwidth]{figures/per_run/017_p2_a2c_random_weather_seed0_blockpen0_02__primary_metric_vs_global_step.png}
    \caption{017 P2 A2C Random Weather Seed0 Blockpen0 02 - Primary Metric Vs Global Step}
\end{figure}

\begin{figure}[H]
    \centering
    \includegraphics[width=0.8\textwidth]{figures/per_run/017_p2_a2c_random_weather_seed0_blockpen0_02__training_reward_vs_global_step.png}
    \caption{017 P2 A2C Random Weather Seed0 Blockpen0 02 - Training Reward Vs Global Step}
\end{figure}

\begin{figure}[H]
    \centering
    \includegraphics[width=0.8\textwidth]{figures/per_run/017_p2_a2c_random_weather_seed0_blockpen0_02__weekly_npk_behavior.png}
    \caption{017 P2 A2C Random Weather Seed0 Blockpen0 02 - Weekly Npk Behavior}
\end{figure}

\clearpage
\subsection{Run: 018 Fertilization Ppo Nonadaptive Fixed Weather Years 3000 Seed 1}
\begin{figure}[H]
    \centering
    \includegraphics[width=0.8\textwidth]{figures/per_run/018_fertilization_ppo_nonadaptive_fixed_weather_years_3000_seed_1__checkpoint_eval_curves.png}
    \caption{018 Fertilization Ppo Nonadaptive Fixed Weather Years 3000 Seed 1 - Checkpoint Eval Curves}
\end{figure}

\begin{figure}[H]
    \centering
    \includegraphics[width=0.8\textwidth]{figures/per_run/018_fertilization_ppo_nonadaptive_fixed_weather_years_3000_seed_1__diagnostics_panel.png}
    \caption{018 Fertilization Ppo Nonadaptive Fixed Weather Years 3000 Seed 1 - Diagnostics Panel}
\end{figure}

\begin{figure}[H]
    \centering
    \includegraphics[width=0.8\textwidth]{figures/per_run/018_fertilization_ppo_nonadaptive_fixed_weather_years_3000_seed_1__episode_length_vs_global_step.png}
    \caption{018 Fertilization Ppo Nonadaptive Fixed Weather Years 3000 Seed 1 - Episode Length Vs Global Step}
\end{figure}

\begin{figure}[H]
    \centering
    \includegraphics[width=0.8\textwidth]{figures/per_run/018_fertilization_ppo_nonadaptive_fixed_weather_years_3000_seed_1__primary_metric_vs_global_step.png}
    \caption{018 Fertilization Ppo Nonadaptive Fixed Weather Years 3000 Seed 1 - Primary Metric Vs Global Step}
\end{figure}

\begin{figure}[H]
    \centering
    \includegraphics[width=0.8\textwidth]{figures/per_run/018_fertilization_ppo_nonadaptive_fixed_weather_years_3000_seed_1__training_reward_vs_global_step.png}
    \caption{018 Fertilization Ppo Nonadaptive Fixed Weather Years 3000 Seed 1 - Training Reward Vs Global Step}
\end{figure}

\clearpage
\subsection{Run: 018 P2 A2C Random Weather Seed0 Blockpen0 05}
\begin{figure}[H]
    \centering
    \includegraphics[width=0.8\textwidth]{figures/per_run/018_p2_a2c_random_weather_seed0_blockpen0_05__blocked_cost_summary.png}
    \caption{018 P2 A2C Random Weather Seed0 Blockpen0 05 - Blocked Cost Summary}
\end{figure}

\begin{figure}[H]
    \centering
    \includegraphics[width=0.8\textwidth]{figures/per_run/018_p2_a2c_random_weather_seed0_blockpen0_05__checkpoint_eval_curves.png}
    \caption{018 P2 A2C Random Weather Seed0 Blockpen0 05 - Checkpoint Eval Curves}
\end{figure}

\begin{figure}[H]
    \centering
    \includegraphics[width=0.8\textwidth]{figures/per_run/018_p2_a2c_random_weather_seed0_blockpen0_05__compliance_summary.png}
    \caption{018 P2 A2C Random Weather Seed0 Blockpen0 05 - Compliance Summary}
\end{figure}

\begin{figure}[H]
    \centering
    \includegraphics[width=0.8\textwidth]{figures/per_run/018_p2_a2c_random_weather_seed0_blockpen0_05__crop_decision_timeline.png}
    \caption{018 P2 A2C Random Weather Seed0 Blockpen0 05 - Crop Decision Timeline}
\end{figure}

\begin{figure}[H]
    \centering
    \includegraphics[width=0.8\textwidth]{figures/per_run/018_p2_a2c_random_weather_seed0_blockpen0_05__diagnostics_panel.png}
    \caption{018 P2 A2C Random Weather Seed0 Blockpen0 05 - Diagnostics Panel}
\end{figure}

\begin{figure}[H]
    \centering
    \includegraphics[width=0.8\textwidth]{figures/per_run/018_p2_a2c_random_weather_seed0_blockpen0_05__episode_length_vs_global_step.png}
    \caption{018 P2 A2C Random Weather Seed0 Blockpen0 05 - Episode Length Vs Global Step}
\end{figure}

\begin{figure}[H]
    \centering
    \includegraphics[width=0.8\textwidth]{figures/per_run/018_p2_a2c_random_weather_seed0_blockpen0_05__primary_metric_vs_global_step.png}
    \caption{018 P2 A2C Random Weather Seed0 Blockpen0 05 - Primary Metric Vs Global Step}
\end{figure}

\begin{figure}[H]
    \centering
    \includegraphics[width=0.8\textwidth]{figures/per_run/018_p2_a2c_random_weather_seed0_blockpen0_05__training_reward_vs_global_step.png}
    \caption{018 P2 A2C Random Weather Seed0 Blockpen0 05 - Training Reward Vs Global Step}
\end{figure}

\begin{figure}[H]
    \centering
    \includegraphics[width=0.8\textwidth]{figures/per_run/018_p2_a2c_random_weather_seed0_blockpen0_05__weekly_npk_behavior.png}
    \caption{018 P2 A2C Random Weather Seed0 Blockpen0 05 - Weekly Npk Behavior}
\end{figure}

\clearpage
\subsection{Run: 019 Fertilization Ppo Adaptive Random Weather Years 3000 Seed 1}
\begin{figure}[H]
    \centering
    \includegraphics[width=0.8\textwidth]{figures/per_run/019_fertilization_ppo_adaptive_random_weather_years_3000_seed_1__checkpoint_eval_curves.png}
    \caption{019 Fertilization Ppo Adaptive Random Weather Years 3000 Seed 1 - Checkpoint Eval Curves}
\end{figure}

\begin{figure}[H]
    \centering
    \includegraphics[width=0.8\textwidth]{figures/per_run/019_fertilization_ppo_adaptive_random_weather_years_3000_seed_1__diagnostics_panel.png}
    \caption{019 Fertilization Ppo Adaptive Random Weather Years 3000 Seed 1 - Diagnostics Panel}
\end{figure}

\begin{figure}[H]
    \centering
    \includegraphics[width=0.8\textwidth]{figures/per_run/019_fertilization_ppo_adaptive_random_weather_years_3000_seed_1__episode_length_vs_global_step.png}
    \caption{019 Fertilization Ppo Adaptive Random Weather Years 3000 Seed 1 - Episode Length Vs Global Step}
\end{figure}

\begin{figure}[H]
    \centering
    \includegraphics[width=0.8\textwidth]{figures/per_run/019_fertilization_ppo_adaptive_random_weather_years_3000_seed_1__primary_metric_vs_global_step.png}
    \caption{019 Fertilization Ppo Adaptive Random Weather Years 3000 Seed 1 - Primary Metric Vs Global Step}
\end{figure}

\begin{figure}[H]
    \centering
    \includegraphics[width=0.8\textwidth]{figures/per_run/019_fertilization_ppo_adaptive_random_weather_years_3000_seed_1__training_reward_vs_global_step.png}
    \caption{019 Fertilization Ppo Adaptive Random Weather Years 3000 Seed 1 - Training Reward Vs Global Step}
\end{figure}

\clearpage
\subsection{Run: 019 P2 Ppo Fixed Weather Seed0 Blockpen0}
\begin{figure}[H]
    \centering
    \includegraphics[width=0.8\textwidth]{figures/per_run/019_p2_ppo_fixed_weather_seed0_blockpen0__blocked_cost_summary.png}
    \caption{019 P2 Ppo Fixed Weather Seed0 Blockpen0 - Blocked Cost Summary}
\end{figure}

\begin{figure}[H]
    \centering
    \includegraphics[width=0.8\textwidth]{figures/per_run/019_p2_ppo_fixed_weather_seed0_blockpen0__checkpoint_eval_curves.png}
    \caption{019 P2 Ppo Fixed Weather Seed0 Blockpen0 - Checkpoint Eval Curves}
\end{figure}

\begin{figure}[H]
    \centering
    \includegraphics[width=0.8\textwidth]{figures/per_run/019_p2_ppo_fixed_weather_seed0_blockpen0__compliance_summary.png}
    \caption{019 P2 Ppo Fixed Weather Seed0 Blockpen0 - Compliance Summary}
\end{figure}

\begin{figure}[H]
    \centering
    \includegraphics[width=0.8\textwidth]{figures/per_run/019_p2_ppo_fixed_weather_seed0_blockpen0__crop_decision_timeline.png}
    \caption{019 P2 Ppo Fixed Weather Seed0 Blockpen0 - Crop Decision Timeline}
\end{figure}

\begin{figure}[H]
    \centering
    \includegraphics[width=0.8\textwidth]{figures/per_run/019_p2_ppo_fixed_weather_seed0_blockpen0__diagnostics_panel.png}
    \caption{019 P2 Ppo Fixed Weather Seed0 Blockpen0 - Diagnostics Panel}
\end{figure}

\begin{figure}[H]
    \centering
    \includegraphics[width=0.8\textwidth]{figures/per_run/019_p2_ppo_fixed_weather_seed0_blockpen0__episode_length_vs_global_step.png}
    \caption{019 P2 Ppo Fixed Weather Seed0 Blockpen0 - Episode Length Vs Global Step}
\end{figure}

\begin{figure}[H]
    \centering
    \includegraphics[width=0.8\textwidth]{figures/per_run/019_p2_ppo_fixed_weather_seed0_blockpen0__primary_metric_vs_global_step.png}
    \caption{019 P2 Ppo Fixed Weather Seed0 Blockpen0 - Primary Metric Vs Global Step}
\end{figure}

\begin{figure}[H]
    \centering
    \includegraphics[width=0.8\textwidth]{figures/per_run/019_p2_ppo_fixed_weather_seed0_blockpen0__training_reward_vs_global_step.png}
    \caption{019 P2 Ppo Fixed Weather Seed0 Blockpen0 - Training Reward Vs Global Step}
\end{figure}

\begin{figure}[H]
    \centering
    \includegraphics[width=0.8\textwidth]{figures/per_run/019_p2_ppo_fixed_weather_seed0_blockpen0__weekly_npk_behavior.png}
    \caption{019 P2 Ppo Fixed Weather Seed0 Blockpen0 - Weekly Npk Behavior}
\end{figure}

\clearpage
\subsection{Run: 020 Fertilization Ppo Nonadaptive Random Weather Years 3000 Seed 1}
\begin{figure}[H]
    \centering
    \includegraphics[width=0.8\textwidth]{figures/per_run/020_fertilization_ppo_nonadaptive_random_weather_years_3000_seed_1__checkpoint_eval_curves.png}
    \caption{020 Fertilization Ppo Nonadaptive Random Weather Years 3000 Seed 1 - Checkpoint Eval Curves}
\end{figure}

\begin{figure}[H]
    \centering
    \includegraphics[width=0.8\textwidth]{figures/per_run/020_fertilization_ppo_nonadaptive_random_weather_years_3000_seed_1__diagnostics_panel.png}
    \caption{020 Fertilization Ppo Nonadaptive Random Weather Years 3000 Seed 1 - Diagnostics Panel}
\end{figure}

\begin{figure}[H]
    \centering
    \includegraphics[width=0.8\textwidth]{figures/per_run/020_fertilization_ppo_nonadaptive_random_weather_years_3000_seed_1__episode_length_vs_global_step.png}
    \caption{020 Fertilization Ppo Nonadaptive Random Weather Years 3000 Seed 1 - Episode Length Vs Global Step}
\end{figure}

\begin{figure}[H]
    \centering
    \includegraphics[width=0.8\textwidth]{figures/per_run/020_fertilization_ppo_nonadaptive_random_weather_years_3000_seed_1__primary_metric_vs_global_step.png}
    \caption{020 Fertilization Ppo Nonadaptive Random Weather Years 3000 Seed 1 - Primary Metric Vs Global Step}
\end{figure}

\begin{figure}[H]
    \centering
    \includegraphics[width=0.8\textwidth]{figures/per_run/020_fertilization_ppo_nonadaptive_random_weather_years_3000_seed_1__training_reward_vs_global_step.png}
    \caption{020 Fertilization Ppo Nonadaptive Random Weather Years 3000 Seed 1 - Training Reward Vs Global Step}
\end{figure}

\clearpage
\subsection{Run: 020 P2 Ppo Fixed Weather Seed0 Blockpen0 02}
\begin{figure}[H]
    \centering
    \includegraphics[width=0.8\textwidth]{figures/per_run/020_p2_ppo_fixed_weather_seed0_blockpen0_02__blocked_cost_summary.png}
    \caption{020 P2 Ppo Fixed Weather Seed0 Blockpen0 02 - Blocked Cost Summary}
\end{figure}

\begin{figure}[H]
    \centering
    \includegraphics[width=0.8\textwidth]{figures/per_run/020_p2_ppo_fixed_weather_seed0_blockpen0_02__checkpoint_eval_curves.png}
    \caption{020 P2 Ppo Fixed Weather Seed0 Blockpen0 02 - Checkpoint Eval Curves}
\end{figure}

\begin{figure}[H]
    \centering
    \includegraphics[width=0.8\textwidth]{figures/per_run/020_p2_ppo_fixed_weather_seed0_blockpen0_02__compliance_summary.png}
    \caption{020 P2 Ppo Fixed Weather Seed0 Blockpen0 02 - Compliance Summary}
\end{figure}

\begin{figure}[H]
    \centering
    \includegraphics[width=0.8\textwidth]{figures/per_run/020_p2_ppo_fixed_weather_seed0_blockpen0_02__crop_decision_timeline.png}
    \caption{020 P2 Ppo Fixed Weather Seed0 Blockpen0 02 - Crop Decision Timeline}
\end{figure}

\begin{figure}[H]
    \centering
    \includegraphics[width=0.8\textwidth]{figures/per_run/020_p2_ppo_fixed_weather_seed0_blockpen0_02__diagnostics_panel.png}
    \caption{020 P2 Ppo Fixed Weather Seed0 Blockpen0 02 - Diagnostics Panel}
\end{figure}

\begin{figure}[H]
    \centering
    \includegraphics[width=0.8\textwidth]{figures/per_run/020_p2_ppo_fixed_weather_seed0_blockpen0_02__episode_length_vs_global_step.png}
    \caption{020 P2 Ppo Fixed Weather Seed0 Blockpen0 02 - Episode Length Vs Global Step}
\end{figure}

\begin{figure}[H]
    \centering
    \includegraphics[width=0.8\textwidth]{figures/per_run/020_p2_ppo_fixed_weather_seed0_blockpen0_02__primary_metric_vs_global_step.png}
    \caption{020 P2 Ppo Fixed Weather Seed0 Blockpen0 02 - Primary Metric Vs Global Step}
\end{figure}

\begin{figure}[H]
    \centering
    \includegraphics[width=0.8\textwidth]{figures/per_run/020_p2_ppo_fixed_weather_seed0_blockpen0_02__training_reward_vs_global_step.png}
    \caption{020 P2 Ppo Fixed Weather Seed0 Blockpen0 02 - Training Reward Vs Global Step}
\end{figure}

\begin{figure}[H]
    \centering
    \includegraphics[width=0.8\textwidth]{figures/per_run/020_p2_ppo_fixed_weather_seed0_blockpen0_02__weekly_npk_behavior.png}
    \caption{020 P2 Ppo Fixed Weather Seed0 Blockpen0 02 - Weekly Npk Behavior}
\end{figure}

\clearpage
\subsection{Run: 021 Fertilization Ppo Adaptive Fixed Weather Years 3000 Seed 2}
\begin{figure}[H]
    \centering
    \includegraphics[width=0.8\textwidth]{figures/per_run/021_fertilization_ppo_adaptive_fixed_weather_years_3000_seed_2__checkpoint_eval_curves.png}
    \caption{021 Fertilization Ppo Adaptive Fixed Weather Years 3000 Seed 2 - Checkpoint Eval Curves}
\end{figure}

\begin{figure}[H]
    \centering
    \includegraphics[width=0.8\textwidth]{figures/per_run/021_fertilization_ppo_adaptive_fixed_weather_years_3000_seed_2__diagnostics_panel.png}
    \caption{021 Fertilization Ppo Adaptive Fixed Weather Years 3000 Seed 2 - Diagnostics Panel}
\end{figure}

\begin{figure}[H]
    \centering
    \includegraphics[width=0.8\textwidth]{figures/per_run/021_fertilization_ppo_adaptive_fixed_weather_years_3000_seed_2__episode_length_vs_global_step.png}
    \caption{021 Fertilization Ppo Adaptive Fixed Weather Years 3000 Seed 2 - Episode Length Vs Global Step}
\end{figure}

\begin{figure}[H]
    \centering
    \includegraphics[width=0.8\textwidth]{figures/per_run/021_fertilization_ppo_adaptive_fixed_weather_years_3000_seed_2__primary_metric_vs_global_step.png}
    \caption{021 Fertilization Ppo Adaptive Fixed Weather Years 3000 Seed 2 - Primary Metric Vs Global Step}
\end{figure}

\begin{figure}[H]
    \centering
    \includegraphics[width=0.8\textwidth]{figures/per_run/021_fertilization_ppo_adaptive_fixed_weather_years_3000_seed_2__training_reward_vs_global_step.png}
    \caption{021 Fertilization Ppo Adaptive Fixed Weather Years 3000 Seed 2 - Training Reward Vs Global Step}
\end{figure}

\clearpage
\subsection{Run: 021 P2 Ppo Fixed Weather Seed0 Blockpen0 05}
\begin{figure}[H]
    \centering
    \includegraphics[width=0.8\textwidth]{figures/per_run/021_p2_ppo_fixed_weather_seed0_blockpen0_05__blocked_cost_summary.png}
    \caption{021 P2 Ppo Fixed Weather Seed0 Blockpen0 05 - Blocked Cost Summary}
\end{figure}

\begin{figure}[H]
    \centering
    \includegraphics[width=0.8\textwidth]{figures/per_run/021_p2_ppo_fixed_weather_seed0_blockpen0_05__checkpoint_eval_curves.png}
    \caption{021 P2 Ppo Fixed Weather Seed0 Blockpen0 05 - Checkpoint Eval Curves}
\end{figure}

\begin{figure}[H]
    \centering
    \includegraphics[width=0.8\textwidth]{figures/per_run/021_p2_ppo_fixed_weather_seed0_blockpen0_05__compliance_summary.png}
    \caption{021 P2 Ppo Fixed Weather Seed0 Blockpen0 05 - Compliance Summary}
\end{figure}

\begin{figure}[H]
    \centering
    \includegraphics[width=0.8\textwidth]{figures/per_run/021_p2_ppo_fixed_weather_seed0_blockpen0_05__crop_decision_timeline.png}
    \caption{021 P2 Ppo Fixed Weather Seed0 Blockpen0 05 - Crop Decision Timeline}
\end{figure}

\begin{figure}[H]
    \centering
    \includegraphics[width=0.8\textwidth]{figures/per_run/021_p2_ppo_fixed_weather_seed0_blockpen0_05__diagnostics_panel.png}
    \caption{021 P2 Ppo Fixed Weather Seed0 Blockpen0 05 - Diagnostics Panel}
\end{figure}

\begin{figure}[H]
    \centering
    \includegraphics[width=0.8\textwidth]{figures/per_run/021_p2_ppo_fixed_weather_seed0_blockpen0_05__episode_length_vs_global_step.png}
    \caption{021 P2 Ppo Fixed Weather Seed0 Blockpen0 05 - Episode Length Vs Global Step}
\end{figure}

\begin{figure}[H]
    \centering
    \includegraphics[width=0.8\textwidth]{figures/per_run/021_p2_ppo_fixed_weather_seed0_blockpen0_05__primary_metric_vs_global_step.png}
    \caption{021 P2 Ppo Fixed Weather Seed0 Blockpen0 05 - Primary Metric Vs Global Step}
\end{figure}

\begin{figure}[H]
    \centering
    \includegraphics[width=0.8\textwidth]{figures/per_run/021_p2_ppo_fixed_weather_seed0_blockpen0_05__training_reward_vs_global_step.png}
    \caption{021 P2 Ppo Fixed Weather Seed0 Blockpen0 05 - Training Reward Vs Global Step}
\end{figure}

\begin{figure}[H]
    \centering
    \includegraphics[width=0.8\textwidth]{figures/per_run/021_p2_ppo_fixed_weather_seed0_blockpen0_05__weekly_npk_behavior.png}
    \caption{021 P2 Ppo Fixed Weather Seed0 Blockpen0 05 - Weekly Npk Behavior}
\end{figure}

\clearpage
\subsection{Run: 022 Fertilization Ppo Nonadaptive Fixed Weather Years 3000 Seed 2}
\begin{figure}[H]
    \centering
    \includegraphics[width=0.8\textwidth]{figures/per_run/022_fertilization_ppo_nonadaptive_fixed_weather_years_3000_seed_2__checkpoint_eval_curves.png}
    \caption{022 Fertilization Ppo Nonadaptive Fixed Weather Years 3000 Seed 2 - Checkpoint Eval Curves}
\end{figure}

\begin{figure}[H]
    \centering
    \includegraphics[width=0.8\textwidth]{figures/per_run/022_fertilization_ppo_nonadaptive_fixed_weather_years_3000_seed_2__diagnostics_panel.png}
    \caption{022 Fertilization Ppo Nonadaptive Fixed Weather Years 3000 Seed 2 - Diagnostics Panel}
\end{figure}

\begin{figure}[H]
    \centering
    \includegraphics[width=0.8\textwidth]{figures/per_run/022_fertilization_ppo_nonadaptive_fixed_weather_years_3000_seed_2__episode_length_vs_global_step.png}
    \caption{022 Fertilization Ppo Nonadaptive Fixed Weather Years 3000 Seed 2 - Episode Length Vs Global Step}
\end{figure}

\begin{figure}[H]
    \centering
    \includegraphics[width=0.8\textwidth]{figures/per_run/022_fertilization_ppo_nonadaptive_fixed_weather_years_3000_seed_2__primary_metric_vs_global_step.png}
    \caption{022 Fertilization Ppo Nonadaptive Fixed Weather Years 3000 Seed 2 - Primary Metric Vs Global Step}
\end{figure}

\begin{figure}[H]
    \centering
    \includegraphics[width=0.8\textwidth]{figures/per_run/022_fertilization_ppo_nonadaptive_fixed_weather_years_3000_seed_2__training_reward_vs_global_step.png}
    \caption{022 Fertilization Ppo Nonadaptive Fixed Weather Years 3000 Seed 2 - Training Reward Vs Global Step}
\end{figure}

\clearpage
\subsection{Run: 022 P2 Ppo Random Weather Seed0 Blockpen0}
\begin{figure}[H]
    \centering
    \includegraphics[width=0.8\textwidth]{figures/per_run/022_p2_ppo_random_weather_seed0_blockpen0__blocked_cost_summary.png}
    \caption{022 P2 Ppo Random Weather Seed0 Blockpen0 - Blocked Cost Summary}
\end{figure}

\begin{figure}[H]
    \centering
    \includegraphics[width=0.8\textwidth]{figures/per_run/022_p2_ppo_random_weather_seed0_blockpen0__checkpoint_eval_curves.png}
    \caption{022 P2 Ppo Random Weather Seed0 Blockpen0 - Checkpoint Eval Curves}
\end{figure}

\begin{figure}[H]
    \centering
    \includegraphics[width=0.8\textwidth]{figures/per_run/022_p2_ppo_random_weather_seed0_blockpen0__compliance_summary.png}
    \caption{022 P2 Ppo Random Weather Seed0 Blockpen0 - Compliance Summary}
\end{figure}

\begin{figure}[H]
    \centering
    \includegraphics[width=0.8\textwidth]{figures/per_run/022_p2_ppo_random_weather_seed0_blockpen0__crop_decision_timeline.png}
    \caption{022 P2 Ppo Random Weather Seed0 Blockpen0 - Crop Decision Timeline}
\end{figure}

\begin{figure}[H]
    \centering
    \includegraphics[width=0.8\textwidth]{figures/per_run/022_p2_ppo_random_weather_seed0_blockpen0__diagnostics_panel.png}
    \caption{022 P2 Ppo Random Weather Seed0 Blockpen0 - Diagnostics Panel}
\end{figure}

\begin{figure}[H]
    \centering
    \includegraphics[width=0.8\textwidth]{figures/per_run/022_p2_ppo_random_weather_seed0_blockpen0__episode_length_vs_global_step.png}
    \caption{022 P2 Ppo Random Weather Seed0 Blockpen0 - Episode Length Vs Global Step}
\end{figure}

\begin{figure}[H]
    \centering
    \includegraphics[width=0.8\textwidth]{figures/per_run/022_p2_ppo_random_weather_seed0_blockpen0__primary_metric_vs_global_step.png}
    \caption{022 P2 Ppo Random Weather Seed0 Blockpen0 - Primary Metric Vs Global Step}
\end{figure}

\begin{figure}[H]
    \centering
    \includegraphics[width=0.8\textwidth]{figures/per_run/022_p2_ppo_random_weather_seed0_blockpen0__training_reward_vs_global_step.png}
    \caption{022 P2 Ppo Random Weather Seed0 Blockpen0 - Training Reward Vs Global Step}
\end{figure}

\begin{figure}[H]
    \centering
    \includegraphics[width=0.8\textwidth]{figures/per_run/022_p2_ppo_random_weather_seed0_blockpen0__weekly_npk_behavior.png}
    \caption{022 P2 Ppo Random Weather Seed0 Blockpen0 - Weekly Npk Behavior}
\end{figure}

\clearpage
\subsection{Run: 023 Fertilization Ppo Adaptive Random Weather Years 3000 Seed 2}
\begin{figure}[H]
    \centering
    \includegraphics[width=0.8\textwidth]{figures/per_run/023_fertilization_ppo_adaptive_random_weather_years_3000_seed_2__checkpoint_eval_curves.png}
    \caption{023 Fertilization Ppo Adaptive Random Weather Years 3000 Seed 2 - Checkpoint Eval Curves}
\end{figure}

\begin{figure}[H]
    \centering
    \includegraphics[width=0.8\textwidth]{figures/per_run/023_fertilization_ppo_adaptive_random_weather_years_3000_seed_2__diagnostics_panel.png}
    \caption{023 Fertilization Ppo Adaptive Random Weather Years 3000 Seed 2 - Diagnostics Panel}
\end{figure}

\begin{figure}[H]
    \centering
    \includegraphics[width=0.8\textwidth]{figures/per_run/023_fertilization_ppo_adaptive_random_weather_years_3000_seed_2__episode_length_vs_global_step.png}
    \caption{023 Fertilization Ppo Adaptive Random Weather Years 3000 Seed 2 - Episode Length Vs Global Step}
\end{figure}

\begin{figure}[H]
    \centering
    \includegraphics[width=0.8\textwidth]{figures/per_run/023_fertilization_ppo_adaptive_random_weather_years_3000_seed_2__primary_metric_vs_global_step.png}
    \caption{023 Fertilization Ppo Adaptive Random Weather Years 3000 Seed 2 - Primary Metric Vs Global Step}
\end{figure}

\begin{figure}[H]
    \centering
    \includegraphics[width=0.8\textwidth]{figures/per_run/023_fertilization_ppo_adaptive_random_weather_years_3000_seed_2__training_reward_vs_global_step.png}
    \caption{023 Fertilization Ppo Adaptive Random Weather Years 3000 Seed 2 - Training Reward Vs Global Step}
\end{figure}

\clearpage
\subsection{Run: 023 P2 Ppo Random Weather Seed0 Blockpen0 02}
\begin{figure}[H]
    \centering
    \includegraphics[width=0.8\textwidth]{figures/per_run/023_p2_ppo_random_weather_seed0_blockpen0_02__blocked_cost_summary.png}
    \caption{023 P2 Ppo Random Weather Seed0 Blockpen0 02 - Blocked Cost Summary}
\end{figure}

\begin{figure}[H]
    \centering
    \includegraphics[width=0.8\textwidth]{figures/per_run/023_p2_ppo_random_weather_seed0_blockpen0_02__checkpoint_eval_curves.png}
    \caption{023 P2 Ppo Random Weather Seed0 Blockpen0 02 - Checkpoint Eval Curves}
\end{figure}

\begin{figure}[H]
    \centering
    \includegraphics[width=0.8\textwidth]{figures/per_run/023_p2_ppo_random_weather_seed0_blockpen0_02__compliance_summary.png}
    \caption{023 P2 Ppo Random Weather Seed0 Blockpen0 02 - Compliance Summary}
\end{figure}

\begin{figure}[H]
    \centering
    \includegraphics[width=0.8\textwidth]{figures/per_run/023_p2_ppo_random_weather_seed0_blockpen0_02__crop_decision_timeline.png}
    \caption{023 P2 Ppo Random Weather Seed0 Blockpen0 02 - Crop Decision Timeline}
\end{figure}

\begin{figure}[H]
    \centering
    \includegraphics[width=0.8\textwidth]{figures/per_run/023_p2_ppo_random_weather_seed0_blockpen0_02__diagnostics_panel.png}
    \caption{023 P2 Ppo Random Weather Seed0 Blockpen0 02 - Diagnostics Panel}
\end{figure}

\begin{figure}[H]
    \centering
    \includegraphics[width=0.8\textwidth]{figures/per_run/023_p2_ppo_random_weather_seed0_blockpen0_02__episode_length_vs_global_step.png}
    \caption{023 P2 Ppo Random Weather Seed0 Blockpen0 02 - Episode Length Vs Global Step}
\end{figure}

\begin{figure}[H]
    \centering
    \includegraphics[width=0.8\textwidth]{figures/per_run/023_p2_ppo_random_weather_seed0_blockpen0_02__primary_metric_vs_global_step.png}
    \caption{023 P2 Ppo Random Weather Seed0 Blockpen0 02 - Primary Metric Vs Global Step}
\end{figure}

\begin{figure}[H]
    \centering
    \includegraphics[width=0.8\textwidth]{figures/per_run/023_p2_ppo_random_weather_seed0_blockpen0_02__training_reward_vs_global_step.png}
    \caption{023 P2 Ppo Random Weather Seed0 Blockpen0 02 - Training Reward Vs Global Step}
\end{figure}

\begin{figure}[H]
    \centering
    \includegraphics[width=0.8\textwidth]{figures/per_run/023_p2_ppo_random_weather_seed0_blockpen0_02__weekly_npk_behavior.png}
    \caption{023 P2 Ppo Random Weather Seed0 Blockpen0 02 - Weekly Npk Behavior}
\end{figure}

\clearpage
\subsection{Run: 024 Fertilization Ppo Nonadaptive Random Weather Years 3000 Seed 2}
\begin{figure}[H]
    \centering
    \includegraphics[width=0.8\textwidth]{figures/per_run/024_fertilization_ppo_nonadaptive_random_weather_years_3000_seed_2__checkpoint_eval_curves.png}
    \caption{024 Fertilization Ppo Nonadaptive Random Weather Years 3000 Seed 2 - Checkpoint Eval Curves}
\end{figure}

\begin{figure}[H]
    \centering
    \includegraphics[width=0.8\textwidth]{figures/per_run/024_fertilization_ppo_nonadaptive_random_weather_years_3000_seed_2__diagnostics_panel.png}
    \caption{024 Fertilization Ppo Nonadaptive Random Weather Years 3000 Seed 2 - Diagnostics Panel}
\end{figure}

\begin{figure}[H]
    \centering
    \includegraphics[width=0.8\textwidth]{figures/per_run/024_fertilization_ppo_nonadaptive_random_weather_years_3000_seed_2__episode_length_vs_global_step.png}
    \caption{024 Fertilization Ppo Nonadaptive Random Weather Years 3000 Seed 2 - Episode Length Vs Global Step}
\end{figure}

\begin{figure}[H]
    \centering
    \includegraphics[width=0.8\textwidth]{figures/per_run/024_fertilization_ppo_nonadaptive_random_weather_years_3000_seed_2__primary_metric_vs_global_step.png}
    \caption{024 Fertilization Ppo Nonadaptive Random Weather Years 3000 Seed 2 - Primary Metric Vs Global Step}
\end{figure}

\begin{figure}[H]
    \centering
    \includegraphics[width=0.8\textwidth]{figures/per_run/024_fertilization_ppo_nonadaptive_random_weather_years_3000_seed_2__training_reward_vs_global_step.png}
    \caption{024 Fertilization Ppo Nonadaptive Random Weather Years 3000 Seed 2 - Training Reward Vs Global Step}
\end{figure}

\clearpage
\subsection{Run: 024 P2 Ppo Random Weather Seed0 Blockpen0 05}
\begin{figure}[H]
    \centering
    \includegraphics[width=0.8\textwidth]{figures/per_run/024_p2_ppo_random_weather_seed0_blockpen0_05__blocked_cost_summary.png}
    \caption{024 P2 Ppo Random Weather Seed0 Blockpen0 05 - Blocked Cost Summary}
\end{figure}

\begin{figure}[H]
    \centering
    \includegraphics[width=0.8\textwidth]{figures/per_run/024_p2_ppo_random_weather_seed0_blockpen0_05__checkpoint_eval_curves.png}
    \caption{024 P2 Ppo Random Weather Seed0 Blockpen0 05 - Checkpoint Eval Curves}
\end{figure}

\begin{figure}[H]
    \centering
    \includegraphics[width=0.8\textwidth]{figures/per_run/024_p2_ppo_random_weather_seed0_blockpen0_05__compliance_summary.png}
    \caption{024 P2 Ppo Random Weather Seed0 Blockpen0 05 - Compliance Summary}
\end{figure}

\begin{figure}[H]
    \centering
    \includegraphics[width=0.8\textwidth]{figures/per_run/024_p2_ppo_random_weather_seed0_blockpen0_05__crop_decision_timeline.png}
    \caption{024 P2 Ppo Random Weather Seed0 Blockpen0 05 - Crop Decision Timeline}
\end{figure}

\begin{figure}[H]
    \centering
    \includegraphics[width=0.8\textwidth]{figures/per_run/024_p2_ppo_random_weather_seed0_blockpen0_05__diagnostics_panel.png}
    \caption{024 P2 Ppo Random Weather Seed0 Blockpen0 05 - Diagnostics Panel}
\end{figure}

\begin{figure}[H]
    \centering
    \includegraphics[width=0.8\textwidth]{figures/per_run/024_p2_ppo_random_weather_seed0_blockpen0_05__episode_length_vs_global_step.png}
    \caption{024 P2 Ppo Random Weather Seed0 Blockpen0 05 - Episode Length Vs Global Step}
\end{figure}

\begin{figure}[H]
    \centering
    \includegraphics[width=0.8\textwidth]{figures/per_run/024_p2_ppo_random_weather_seed0_blockpen0_05__primary_metric_vs_global_step.png}
    \caption{024 P2 Ppo Random Weather Seed0 Blockpen0 05 - Primary Metric Vs Global Step}
\end{figure}

\begin{figure}[H]
    \centering
    \includegraphics[width=0.8\textwidth]{figures/per_run/024_p2_ppo_random_weather_seed0_blockpen0_05__training_reward_vs_global_step.png}
    \caption{024 P2 Ppo Random Weather Seed0 Blockpen0 05 - Training Reward Vs Global Step}
\end{figure}

\begin{figure}[H]
    \centering
    \includegraphics[width=0.8\textwidth]{figures/per_run/024_p2_ppo_random_weather_seed0_blockpen0_05__weekly_npk_behavior.png}
    \caption{024 P2 Ppo Random Weather Seed0 Blockpen0 05 - Weekly Npk Behavior}
\end{figure}

\clearpage
\subsection{Run: 025 Fertilization Ppo Adaptive Fixed Weather Years 5000 Seed 0}
\begin{figure}[H]
    \centering
    \includegraphics[width=0.8\textwidth]{figures/per_run/025_fertilization_ppo_adaptive_fixed_weather_years_5000_seed_0__checkpoint_eval_curves.png}
    \caption{025 Fertilization Ppo Adaptive Fixed Weather Years 5000 Seed 0 - Checkpoint Eval Curves}
\end{figure}

\begin{figure}[H]
    \centering
    \includegraphics[width=0.8\textwidth]{figures/per_run/025_fertilization_ppo_adaptive_fixed_weather_years_5000_seed_0__diagnostics_panel.png}
    \caption{025 Fertilization Ppo Adaptive Fixed Weather Years 5000 Seed 0 - Diagnostics Panel}
\end{figure}

\begin{figure}[H]
    \centering
    \includegraphics[width=0.8\textwidth]{figures/per_run/025_fertilization_ppo_adaptive_fixed_weather_years_5000_seed_0__episode_length_vs_global_step.png}
    \caption{025 Fertilization Ppo Adaptive Fixed Weather Years 5000 Seed 0 - Episode Length Vs Global Step}
\end{figure}

\begin{figure}[H]
    \centering
    \includegraphics[width=0.8\textwidth]{figures/per_run/025_fertilization_ppo_adaptive_fixed_weather_years_5000_seed_0__primary_metric_vs_global_step.png}
    \caption{025 Fertilization Ppo Adaptive Fixed Weather Years 5000 Seed 0 - Primary Metric Vs Global Step}
\end{figure}

\begin{figure}[H]
    \centering
    \includegraphics[width=0.8\textwidth]{figures/per_run/025_fertilization_ppo_adaptive_fixed_weather_years_5000_seed_0__training_reward_vs_global_step.png}
    \caption{025 Fertilization Ppo Adaptive Fixed Weather Years 5000 Seed 0 - Training Reward Vs Global Step}
\end{figure}

\clearpage
\subsection{Run: 025 P3 Ppo Adaptive Fixed Weather Seed0 Costw0 8}
\begin{figure}[H]
    \centering
    \includegraphics[width=0.8\textwidth]{figures/per_run/025_p3_ppo_adaptive_fixed_weather_seed0_costw0_8__checkpoint_eval_curves.png}
    \caption{025 P3 Ppo Adaptive Fixed Weather Seed0 Costw0 8 - Checkpoint Eval Curves}
\end{figure}

\begin{figure}[H]
    \centering
    \includegraphics[width=0.8\textwidth]{figures/per_run/025_p3_ppo_adaptive_fixed_weather_seed0_costw0_8__diagnostics_panel.png}
    \caption{025 P3 Ppo Adaptive Fixed Weather Seed0 Costw0 8 - Diagnostics Panel}
\end{figure}

\begin{figure}[H]
    \centering
    \includegraphics[width=0.8\textwidth]{figures/per_run/025_p3_ppo_adaptive_fixed_weather_seed0_costw0_8__episode_length_vs_global_step.png}
    \caption{025 P3 Ppo Adaptive Fixed Weather Seed0 Costw0 8 - Episode Length Vs Global Step}
\end{figure}

\begin{figure}[H]
    \centering
    \includegraphics[width=0.8\textwidth]{figures/per_run/025_p3_ppo_adaptive_fixed_weather_seed0_costw0_8__primary_metric_vs_global_step.png}
    \caption{025 P3 Ppo Adaptive Fixed Weather Seed0 Costw0 8 - Primary Metric Vs Global Step}
\end{figure}

\begin{figure}[H]
    \centering
    \includegraphics[width=0.8\textwidth]{figures/per_run/025_p3_ppo_adaptive_fixed_weather_seed0_costw0_8__training_reward_vs_global_step.png}
    \caption{025 P3 Ppo Adaptive Fixed Weather Seed0 Costw0 8 - Training Reward Vs Global Step}
\end{figure}

\clearpage
\subsection{Run: 026 Fertilization Ppo Nonadaptive Fixed Weather Years 5000 Seed 0}
\begin{figure}[H]
    \centering
    \includegraphics[width=0.8\textwidth]{figures/per_run/026_fertilization_ppo_nonadaptive_fixed_weather_years_5000_seed_0__checkpoint_eval_curves.png}
    \caption{026 Fertilization Ppo Nonadaptive Fixed Weather Years 5000 Seed 0 - Checkpoint Eval Curves}
\end{figure}

\begin{figure}[H]
    \centering
    \includegraphics[width=0.8\textwidth]{figures/per_run/026_fertilization_ppo_nonadaptive_fixed_weather_years_5000_seed_0__diagnostics_panel.png}
    \caption{026 Fertilization Ppo Nonadaptive Fixed Weather Years 5000 Seed 0 - Diagnostics Panel}
\end{figure}

\begin{figure}[H]
    \centering
    \includegraphics[width=0.8\textwidth]{figures/per_run/026_fertilization_ppo_nonadaptive_fixed_weather_years_5000_seed_0__episode_length_vs_global_step.png}
    \caption{026 Fertilization Ppo Nonadaptive Fixed Weather Years 5000 Seed 0 - Episode Length Vs Global Step}
\end{figure}

\begin{figure}[H]
    \centering
    \includegraphics[width=0.8\textwidth]{figures/per_run/026_fertilization_ppo_nonadaptive_fixed_weather_years_5000_seed_0__primary_metric_vs_global_step.png}
    \caption{026 Fertilization Ppo Nonadaptive Fixed Weather Years 5000 Seed 0 - Primary Metric Vs Global Step}
\end{figure}

\begin{figure}[H]
    \centering
    \includegraphics[width=0.8\textwidth]{figures/per_run/026_fertilization_ppo_nonadaptive_fixed_weather_years_5000_seed_0__training_reward_vs_global_step.png}
    \caption{026 Fertilization Ppo Nonadaptive Fixed Weather Years 5000 Seed 0 - Training Reward Vs Global Step}
\end{figure}

\clearpage
\subsection{Run: 026 P3 Ppo Adaptive Fixed Weather Seed0 Costw1}
\begin{figure}[H]
    \centering
    \includegraphics[width=0.8\textwidth]{figures/per_run/026_p3_ppo_adaptive_fixed_weather_seed0_costw1__checkpoint_eval_curves.png}
    \caption{026 P3 Ppo Adaptive Fixed Weather Seed0 Costw1 - Checkpoint Eval Curves}
\end{figure}

\begin{figure}[H]
    \centering
    \includegraphics[width=0.8\textwidth]{figures/per_run/026_p3_ppo_adaptive_fixed_weather_seed0_costw1__diagnostics_panel.png}
    \caption{026 P3 Ppo Adaptive Fixed Weather Seed0 Costw1 - Diagnostics Panel}
\end{figure}

\begin{figure}[H]
    \centering
    \includegraphics[width=0.8\textwidth]{figures/per_run/026_p3_ppo_adaptive_fixed_weather_seed0_costw1__episode_length_vs_global_step.png}
    \caption{026 P3 Ppo Adaptive Fixed Weather Seed0 Costw1 - Episode Length Vs Global Step}
\end{figure}

\begin{figure}[H]
    \centering
    \includegraphics[width=0.8\textwidth]{figures/per_run/026_p3_ppo_adaptive_fixed_weather_seed0_costw1__primary_metric_vs_global_step.png}
    \caption{026 P3 Ppo Adaptive Fixed Weather Seed0 Costw1 - Primary Metric Vs Global Step}
\end{figure}

\begin{figure}[H]
    \centering
    \includegraphics[width=0.8\textwidth]{figures/per_run/026_p3_ppo_adaptive_fixed_weather_seed0_costw1__training_reward_vs_global_step.png}
    \caption{026 P3 Ppo Adaptive Fixed Weather Seed0 Costw1 - Training Reward Vs Global Step}
\end{figure}

\clearpage
\subsection{Run: 027 Fertilization Ppo Adaptive Random Weather Years 5000 Seed 0}
\begin{figure}[H]
    \centering
    \includegraphics[width=0.8\textwidth]{figures/per_run/027_fertilization_ppo_adaptive_random_weather_years_5000_seed_0__checkpoint_eval_curves.png}
    \caption{027 Fertilization Ppo Adaptive Random Weather Years 5000 Seed 0 - Checkpoint Eval Curves}
\end{figure}

\begin{figure}[H]
    \centering
    \includegraphics[width=0.8\textwidth]{figures/per_run/027_fertilization_ppo_adaptive_random_weather_years_5000_seed_0__diagnostics_panel.png}
    \caption{027 Fertilization Ppo Adaptive Random Weather Years 5000 Seed 0 - Diagnostics Panel}
\end{figure}

\begin{figure}[H]
    \centering
    \includegraphics[width=0.8\textwidth]{figures/per_run/027_fertilization_ppo_adaptive_random_weather_years_5000_seed_0__episode_length_vs_global_step.png}
    \caption{027 Fertilization Ppo Adaptive Random Weather Years 5000 Seed 0 - Episode Length Vs Global Step}
\end{figure}

\begin{figure}[H]
    \centering
    \includegraphics[width=0.8\textwidth]{figures/per_run/027_fertilization_ppo_adaptive_random_weather_years_5000_seed_0__primary_metric_vs_global_step.png}
    \caption{027 Fertilization Ppo Adaptive Random Weather Years 5000 Seed 0 - Primary Metric Vs Global Step}
\end{figure}

\begin{figure}[H]
    \centering
    \includegraphics[width=0.8\textwidth]{figures/per_run/027_fertilization_ppo_adaptive_random_weather_years_5000_seed_0__training_reward_vs_global_step.png}
    \caption{027 Fertilization Ppo Adaptive Random Weather Years 5000 Seed 0 - Training Reward Vs Global Step}
\end{figure}

\clearpage
\subsection{Run: 027 P3 Ppo Adaptive Fixed Weather Seed0 Costw1 2}
\begin{figure}[H]
    \centering
    \includegraphics[width=0.8\textwidth]{figures/per_run/027_p3_ppo_adaptive_fixed_weather_seed0_costw1_2__checkpoint_eval_curves.png}
    \caption{027 P3 Ppo Adaptive Fixed Weather Seed0 Costw1 2 - Checkpoint Eval Curves}
\end{figure}

\begin{figure}[H]
    \centering
    \includegraphics[width=0.8\textwidth]{figures/per_run/027_p3_ppo_adaptive_fixed_weather_seed0_costw1_2__diagnostics_panel.png}
    \caption{027 P3 Ppo Adaptive Fixed Weather Seed0 Costw1 2 - Diagnostics Panel}
\end{figure}

\begin{figure}[H]
    \centering
    \includegraphics[width=0.8\textwidth]{figures/per_run/027_p3_ppo_adaptive_fixed_weather_seed0_costw1_2__episode_length_vs_global_step.png}
    \caption{027 P3 Ppo Adaptive Fixed Weather Seed0 Costw1 2 - Episode Length Vs Global Step}
\end{figure}

\begin{figure}[H]
    \centering
    \includegraphics[width=0.8\textwidth]{figures/per_run/027_p3_ppo_adaptive_fixed_weather_seed0_costw1_2__primary_metric_vs_global_step.png}
    \caption{027 P3 Ppo Adaptive Fixed Weather Seed0 Costw1 2 - Primary Metric Vs Global Step}
\end{figure}

\begin{figure}[H]
    \centering
    \includegraphics[width=0.8\textwidth]{figures/per_run/027_p3_ppo_adaptive_fixed_weather_seed0_costw1_2__training_reward_vs_global_step.png}
    \caption{027 P3 Ppo Adaptive Fixed Weather Seed0 Costw1 2 - Training Reward Vs Global Step}
\end{figure}

\clearpage
\subsection{Run: 028 Fertilization Ppo Nonadaptive Random Weather Years 5000 Seed 0}
\begin{figure}[H]
    \centering
    \includegraphics[width=0.8\textwidth]{figures/per_run/028_fertilization_ppo_nonadaptive_random_weather_years_5000_seed_0__checkpoint_eval_curves.png}
    \caption{028 Fertilization Ppo Nonadaptive Random Weather Years 5000 Seed 0 - Checkpoint Eval Curves}
\end{figure}

\begin{figure}[H]
    \centering
    \includegraphics[width=0.8\textwidth]{figures/per_run/028_fertilization_ppo_nonadaptive_random_weather_years_5000_seed_0__diagnostics_panel.png}
    \caption{028 Fertilization Ppo Nonadaptive Random Weather Years 5000 Seed 0 - Diagnostics Panel}
\end{figure}

\begin{figure}[H]
    \centering
    \includegraphics[width=0.8\textwidth]{figures/per_run/028_fertilization_ppo_nonadaptive_random_weather_years_5000_seed_0__episode_length_vs_global_step.png}
    \caption{028 Fertilization Ppo Nonadaptive Random Weather Years 5000 Seed 0 - Episode Length Vs Global Step}
\end{figure}

\begin{figure}[H]
    \centering
    \includegraphics[width=0.8\textwidth]{figures/per_run/028_fertilization_ppo_nonadaptive_random_weather_years_5000_seed_0__primary_metric_vs_global_step.png}
    \caption{028 Fertilization Ppo Nonadaptive Random Weather Years 5000 Seed 0 - Primary Metric Vs Global Step}
\end{figure}

\begin{figure}[H]
    \centering
    \includegraphics[width=0.8\textwidth]{figures/per_run/028_fertilization_ppo_nonadaptive_random_weather_years_5000_seed_0__training_reward_vs_global_step.png}
    \caption{028 Fertilization Ppo Nonadaptive Random Weather Years 5000 Seed 0 - Training Reward Vs Global Step}
\end{figure}

\clearpage
\subsection{Run: 028 P3 Ppo Adaptive Random Weather Seed0 Costw0 8}
\begin{figure}[H]
    \centering
    \includegraphics[width=0.8\textwidth]{figures/per_run/028_p3_ppo_adaptive_random_weather_seed0_costw0_8__checkpoint_eval_curves.png}
    \caption{028 P3 Ppo Adaptive Random Weather Seed0 Costw0 8 - Checkpoint Eval Curves}
\end{figure}

\begin{figure}[H]
    \centering
    \includegraphics[width=0.8\textwidth]{figures/per_run/028_p3_ppo_adaptive_random_weather_seed0_costw0_8__diagnostics_panel.png}
    \caption{028 P3 Ppo Adaptive Random Weather Seed0 Costw0 8 - Diagnostics Panel}
\end{figure}

\begin{figure}[H]
    \centering
    \includegraphics[width=0.8\textwidth]{figures/per_run/028_p3_ppo_adaptive_random_weather_seed0_costw0_8__episode_length_vs_global_step.png}
    \caption{028 P3 Ppo Adaptive Random Weather Seed0 Costw0 8 - Episode Length Vs Global Step}
\end{figure}

\begin{figure}[H]
    \centering
    \includegraphics[width=0.8\textwidth]{figures/per_run/028_p3_ppo_adaptive_random_weather_seed0_costw0_8__primary_metric_vs_global_step.png}
    \caption{028 P3 Ppo Adaptive Random Weather Seed0 Costw0 8 - Primary Metric Vs Global Step}
\end{figure}

\begin{figure}[H]
    \centering
    \includegraphics[width=0.8\textwidth]{figures/per_run/028_p3_ppo_adaptive_random_weather_seed0_costw0_8__training_reward_vs_global_step.png}
    \caption{028 P3 Ppo Adaptive Random Weather Seed0 Costw0 8 - Training Reward Vs Global Step}
\end{figure}

\clearpage
\subsection{Run: 029 Fertilization Ppo Adaptive Fixed Weather Years 5000 Seed 1}
\begin{figure}[H]
    \centering
    \includegraphics[width=0.8\textwidth]{figures/per_run/029_fertilization_ppo_adaptive_fixed_weather_years_5000_seed_1__checkpoint_eval_curves.png}
    \caption{029 Fertilization Ppo Adaptive Fixed Weather Years 5000 Seed 1 - Checkpoint Eval Curves}
\end{figure}

\begin{figure}[H]
    \centering
    \includegraphics[width=0.8\textwidth]{figures/per_run/029_fertilization_ppo_adaptive_fixed_weather_years_5000_seed_1__diagnostics_panel.png}
    \caption{029 Fertilization Ppo Adaptive Fixed Weather Years 5000 Seed 1 - Diagnostics Panel}
\end{figure}

\begin{figure}[H]
    \centering
    \includegraphics[width=0.8\textwidth]{figures/per_run/029_fertilization_ppo_adaptive_fixed_weather_years_5000_seed_1__episode_length_vs_global_step.png}
    \caption{029 Fertilization Ppo Adaptive Fixed Weather Years 5000 Seed 1 - Episode Length Vs Global Step}
\end{figure}

\begin{figure}[H]
    \centering
    \includegraphics[width=0.8\textwidth]{figures/per_run/029_fertilization_ppo_adaptive_fixed_weather_years_5000_seed_1__primary_metric_vs_global_step.png}
    \caption{029 Fertilization Ppo Adaptive Fixed Weather Years 5000 Seed 1 - Primary Metric Vs Global Step}
\end{figure}

\begin{figure}[H]
    \centering
    \includegraphics[width=0.8\textwidth]{figures/per_run/029_fertilization_ppo_adaptive_fixed_weather_years_5000_seed_1__training_reward_vs_global_step.png}
    \caption{029 Fertilization Ppo Adaptive Fixed Weather Years 5000 Seed 1 - Training Reward Vs Global Step}
\end{figure}

\clearpage
\subsection{Run: 029 P3 Ppo Adaptive Random Weather Seed0 Costw1}
\begin{figure}[H]
    \centering
    \includegraphics[width=0.8\textwidth]{figures/per_run/029_p3_ppo_adaptive_random_weather_seed0_costw1__checkpoint_eval_curves.png}
    \caption{029 P3 Ppo Adaptive Random Weather Seed0 Costw1 - Checkpoint Eval Curves}
\end{figure}

\begin{figure}[H]
    \centering
    \includegraphics[width=0.8\textwidth]{figures/per_run/029_p3_ppo_adaptive_random_weather_seed0_costw1__diagnostics_panel.png}
    \caption{029 P3 Ppo Adaptive Random Weather Seed0 Costw1 - Diagnostics Panel}
\end{figure}

\begin{figure}[H]
    \centering
    \includegraphics[width=0.8\textwidth]{figures/per_run/029_p3_ppo_adaptive_random_weather_seed0_costw1__episode_length_vs_global_step.png}
    \caption{029 P3 Ppo Adaptive Random Weather Seed0 Costw1 - Episode Length Vs Global Step}
\end{figure}

\begin{figure}[H]
    \centering
    \includegraphics[width=0.8\textwidth]{figures/per_run/029_p3_ppo_adaptive_random_weather_seed0_costw1__primary_metric_vs_global_step.png}
    \caption{029 P3 Ppo Adaptive Random Weather Seed0 Costw1 - Primary Metric Vs Global Step}
\end{figure}

\begin{figure}[H]
    \centering
    \includegraphics[width=0.8\textwidth]{figures/per_run/029_p3_ppo_adaptive_random_weather_seed0_costw1__training_reward_vs_global_step.png}
    \caption{029 P3 Ppo Adaptive Random Weather Seed0 Costw1 - Training Reward Vs Global Step}
\end{figure}

\clearpage
\subsection{Run: 030 Fertilization Ppo Nonadaptive Fixed Weather Years 5000 Seed 1}
\begin{figure}[H]
    \centering
    \includegraphics[width=0.8\textwidth]{figures/per_run/030_fertilization_ppo_nonadaptive_fixed_weather_years_5000_seed_1__checkpoint_eval_curves.png}
    \caption{030 Fertilization Ppo Nonadaptive Fixed Weather Years 5000 Seed 1 - Checkpoint Eval Curves}
\end{figure}

\begin{figure}[H]
    \centering
    \includegraphics[width=0.8\textwidth]{figures/per_run/030_fertilization_ppo_nonadaptive_fixed_weather_years_5000_seed_1__diagnostics_panel.png}
    \caption{030 Fertilization Ppo Nonadaptive Fixed Weather Years 5000 Seed 1 - Diagnostics Panel}
\end{figure}

\begin{figure}[H]
    \centering
    \includegraphics[width=0.8\textwidth]{figures/per_run/030_fertilization_ppo_nonadaptive_fixed_weather_years_5000_seed_1__episode_length_vs_global_step.png}
    \caption{030 Fertilization Ppo Nonadaptive Fixed Weather Years 5000 Seed 1 - Episode Length Vs Global Step}
\end{figure}

\begin{figure}[H]
    \centering
    \includegraphics[width=0.8\textwidth]{figures/per_run/030_fertilization_ppo_nonadaptive_fixed_weather_years_5000_seed_1__primary_metric_vs_global_step.png}
    \caption{030 Fertilization Ppo Nonadaptive Fixed Weather Years 5000 Seed 1 - Primary Metric Vs Global Step}
\end{figure}

\begin{figure}[H]
    \centering
    \includegraphics[width=0.8\textwidth]{figures/per_run/030_fertilization_ppo_nonadaptive_fixed_weather_years_5000_seed_1__training_reward_vs_global_step.png}
    \caption{030 Fertilization Ppo Nonadaptive Fixed Weather Years 5000 Seed 1 - Training Reward Vs Global Step}
\end{figure}

\clearpage
\subsection{Run: 030 P3 Ppo Adaptive Random Weather Seed0 Costw1 2}
\begin{figure}[H]
    \centering
    \includegraphics[width=0.8\textwidth]{figures/per_run/030_p3_ppo_adaptive_random_weather_seed0_costw1_2__checkpoint_eval_curves.png}
    \caption{030 P3 Ppo Adaptive Random Weather Seed0 Costw1 2 - Checkpoint Eval Curves}
\end{figure}

\begin{figure}[H]
    \centering
    \includegraphics[width=0.8\textwidth]{figures/per_run/030_p3_ppo_adaptive_random_weather_seed0_costw1_2__diagnostics_panel.png}
    \caption{030 P3 Ppo Adaptive Random Weather Seed0 Costw1 2 - Diagnostics Panel}
\end{figure}

\begin{figure}[H]
    \centering
    \includegraphics[width=0.8\textwidth]{figures/per_run/030_p3_ppo_adaptive_random_weather_seed0_costw1_2__episode_length_vs_global_step.png}
    \caption{030 P3 Ppo Adaptive Random Weather Seed0 Costw1 2 - Episode Length Vs Global Step}
\end{figure}

\begin{figure}[H]
    \centering
    \includegraphics[width=0.8\textwidth]{figures/per_run/030_p3_ppo_adaptive_random_weather_seed0_costw1_2__primary_metric_vs_global_step.png}
    \caption{030 P3 Ppo Adaptive Random Weather Seed0 Costw1 2 - Primary Metric Vs Global Step}
\end{figure}

\begin{figure}[H]
    \centering
    \includegraphics[width=0.8\textwidth]{figures/per_run/030_p3_ppo_adaptive_random_weather_seed0_costw1_2__training_reward_vs_global_step.png}
    \caption{030 P3 Ppo Adaptive Random Weather Seed0 Costw1 2 - Training Reward Vs Global Step}
\end{figure}

\clearpage
\subsection{Run: 031 Fertilization Ppo Adaptive Random Weather Years 5000 Seed 1}
\begin{figure}[H]
    \centering
    \includegraphics[width=0.8\textwidth]{figures/per_run/031_fertilization_ppo_adaptive_random_weather_years_5000_seed_1__checkpoint_eval_curves.png}
    \caption{031 Fertilization Ppo Adaptive Random Weather Years 5000 Seed 1 - Checkpoint Eval Curves}
\end{figure}

\begin{figure}[H]
    \centering
    \includegraphics[width=0.8\textwidth]{figures/per_run/031_fertilization_ppo_adaptive_random_weather_years_5000_seed_1__diagnostics_panel.png}
    \caption{031 Fertilization Ppo Adaptive Random Weather Years 5000 Seed 1 - Diagnostics Panel}
\end{figure}

\begin{figure}[H]
    \centering
    \includegraphics[width=0.8\textwidth]{figures/per_run/031_fertilization_ppo_adaptive_random_weather_years_5000_seed_1__episode_length_vs_global_step.png}
    \caption{031 Fertilization Ppo Adaptive Random Weather Years 5000 Seed 1 - Episode Length Vs Global Step}
\end{figure}

\begin{figure}[H]
    \centering
    \includegraphics[width=0.8\textwidth]{figures/per_run/031_fertilization_ppo_adaptive_random_weather_years_5000_seed_1__primary_metric_vs_global_step.png}
    \caption{031 Fertilization Ppo Adaptive Random Weather Years 5000 Seed 1 - Primary Metric Vs Global Step}
\end{figure}

\begin{figure}[H]
    \centering
    \includegraphics[width=0.8\textwidth]{figures/per_run/031_fertilization_ppo_adaptive_random_weather_years_5000_seed_1__training_reward_vs_global_step.png}
    \caption{031 Fertilization Ppo Adaptive Random Weather Years 5000 Seed 1 - Training Reward Vs Global Step}
\end{figure}

\clearpage
\subsection{Run: 031 P3 Ppo Adaptive Fixed Weather Seed1 Costw0 8}
\begin{figure}[H]
    \centering
    \includegraphics[width=0.8\textwidth]{figures/per_run/031_p3_ppo_adaptive_fixed_weather_seed1_costw0_8__checkpoint_eval_curves.png}
    \caption{031 P3 Ppo Adaptive Fixed Weather Seed1 Costw0 8 - Checkpoint Eval Curves}
\end{figure}

\begin{figure}[H]
    \centering
    \includegraphics[width=0.8\textwidth]{figures/per_run/031_p3_ppo_adaptive_fixed_weather_seed1_costw0_8__diagnostics_panel.png}
    \caption{031 P3 Ppo Adaptive Fixed Weather Seed1 Costw0 8 - Diagnostics Panel}
\end{figure}

\begin{figure}[H]
    \centering
    \includegraphics[width=0.8\textwidth]{figures/per_run/031_p3_ppo_adaptive_fixed_weather_seed1_costw0_8__episode_length_vs_global_step.png}
    \caption{031 P3 Ppo Adaptive Fixed Weather Seed1 Costw0 8 - Episode Length Vs Global Step}
\end{figure}

\begin{figure}[H]
    \centering
    \includegraphics[width=0.8\textwidth]{figures/per_run/031_p3_ppo_adaptive_fixed_weather_seed1_costw0_8__primary_metric_vs_global_step.png}
    \caption{031 P3 Ppo Adaptive Fixed Weather Seed1 Costw0 8 - Primary Metric Vs Global Step}
\end{figure}

\begin{figure}[H]
    \centering
    \includegraphics[width=0.8\textwidth]{figures/per_run/031_p3_ppo_adaptive_fixed_weather_seed1_costw0_8__training_reward_vs_global_step.png}
    \caption{031 P3 Ppo Adaptive Fixed Weather Seed1 Costw0 8 - Training Reward Vs Global Step}
\end{figure}

\clearpage
\subsection{Run: 032 Fertilization Ppo Nonadaptive Random Weather Years 5000 Seed 1}
\begin{figure}[H]
    \centering
    \includegraphics[width=0.8\textwidth]{figures/per_run/032_fertilization_ppo_nonadaptive_random_weather_years_5000_seed_1__checkpoint_eval_curves.png}
    \caption{032 Fertilization Ppo Nonadaptive Random Weather Years 5000 Seed 1 - Checkpoint Eval Curves}
\end{figure}

\begin{figure}[H]
    \centering
    \includegraphics[width=0.8\textwidth]{figures/per_run/032_fertilization_ppo_nonadaptive_random_weather_years_5000_seed_1__diagnostics_panel.png}
    \caption{032 Fertilization Ppo Nonadaptive Random Weather Years 5000 Seed 1 - Diagnostics Panel}
\end{figure}

\begin{figure}[H]
    \centering
    \includegraphics[width=0.8\textwidth]{figures/per_run/032_fertilization_ppo_nonadaptive_random_weather_years_5000_seed_1__episode_length_vs_global_step.png}
    \caption{032 Fertilization Ppo Nonadaptive Random Weather Years 5000 Seed 1 - Episode Length Vs Global Step}
\end{figure}

\begin{figure}[H]
    \centering
    \includegraphics[width=0.8\textwidth]{figures/per_run/032_fertilization_ppo_nonadaptive_random_weather_years_5000_seed_1__primary_metric_vs_global_step.png}
    \caption{032 Fertilization Ppo Nonadaptive Random Weather Years 5000 Seed 1 - Primary Metric Vs Global Step}
\end{figure}

\begin{figure}[H]
    \centering
    \includegraphics[width=0.8\textwidth]{figures/per_run/032_fertilization_ppo_nonadaptive_random_weather_years_5000_seed_1__training_reward_vs_global_step.png}
    \caption{032 Fertilization Ppo Nonadaptive Random Weather Years 5000 Seed 1 - Training Reward Vs Global Step}
\end{figure}

\clearpage
\subsection{Run: 032 P3 Ppo Adaptive Fixed Weather Seed1 Costw1}
\begin{figure}[H]
    \centering
    \includegraphics[width=0.8\textwidth]{figures/per_run/032_p3_ppo_adaptive_fixed_weather_seed1_costw1__checkpoint_eval_curves.png}
    \caption{032 P3 Ppo Adaptive Fixed Weather Seed1 Costw1 - Checkpoint Eval Curves}
\end{figure}

\begin{figure}[H]
    \centering
    \includegraphics[width=0.8\textwidth]{figures/per_run/032_p3_ppo_adaptive_fixed_weather_seed1_costw1__diagnostics_panel.png}
    \caption{032 P3 Ppo Adaptive Fixed Weather Seed1 Costw1 - Diagnostics Panel}
\end{figure}

\begin{figure}[H]
    \centering
    \includegraphics[width=0.8\textwidth]{figures/per_run/032_p3_ppo_adaptive_fixed_weather_seed1_costw1__episode_length_vs_global_step.png}
    \caption{032 P3 Ppo Adaptive Fixed Weather Seed1 Costw1 - Episode Length Vs Global Step}
\end{figure}

\begin{figure}[H]
    \centering
    \includegraphics[width=0.8\textwidth]{figures/per_run/032_p3_ppo_adaptive_fixed_weather_seed1_costw1__primary_metric_vs_global_step.png}
    \caption{032 P3 Ppo Adaptive Fixed Weather Seed1 Costw1 - Primary Metric Vs Global Step}
\end{figure}

\begin{figure}[H]
    \centering
    \includegraphics[width=0.8\textwidth]{figures/per_run/032_p3_ppo_adaptive_fixed_weather_seed1_costw1__training_reward_vs_global_step.png}
    \caption{032 P3 Ppo Adaptive Fixed Weather Seed1 Costw1 - Training Reward Vs Global Step}
\end{figure}

\clearpage
\subsection{Run: 033 Fertilization Ppo Adaptive Fixed Weather Years 5000 Seed 2}
\begin{figure}[H]
    \centering
    \includegraphics[width=0.8\textwidth]{figures/per_run/033_fertilization_ppo_adaptive_fixed_weather_years_5000_seed_2__checkpoint_eval_curves.png}
    \caption{033 Fertilization Ppo Adaptive Fixed Weather Years 5000 Seed 2 - Checkpoint Eval Curves}
\end{figure}

\begin{figure}[H]
    \centering
    \includegraphics[width=0.8\textwidth]{figures/per_run/033_fertilization_ppo_adaptive_fixed_weather_years_5000_seed_2__diagnostics_panel.png}
    \caption{033 Fertilization Ppo Adaptive Fixed Weather Years 5000 Seed 2 - Diagnostics Panel}
\end{figure}

\begin{figure}[H]
    \centering
    \includegraphics[width=0.8\textwidth]{figures/per_run/033_fertilization_ppo_adaptive_fixed_weather_years_5000_seed_2__episode_length_vs_global_step.png}
    \caption{033 Fertilization Ppo Adaptive Fixed Weather Years 5000 Seed 2 - Episode Length Vs Global Step}
\end{figure}

\begin{figure}[H]
    \centering
    \includegraphics[width=0.8\textwidth]{figures/per_run/033_fertilization_ppo_adaptive_fixed_weather_years_5000_seed_2__primary_metric_vs_global_step.png}
    \caption{033 Fertilization Ppo Adaptive Fixed Weather Years 5000 Seed 2 - Primary Metric Vs Global Step}
\end{figure}

\begin{figure}[H]
    \centering
    \includegraphics[width=0.8\textwidth]{figures/per_run/033_fertilization_ppo_adaptive_fixed_weather_years_5000_seed_2__training_reward_vs_global_step.png}
    \caption{033 Fertilization Ppo Adaptive Fixed Weather Years 5000 Seed 2 - Training Reward Vs Global Step}
\end{figure}

\clearpage
\subsection{Run: 033 P3 Ppo Adaptive Fixed Weather Seed1 Costw1 2}
\begin{figure}[H]
    \centering
    \includegraphics[width=0.8\textwidth]{figures/per_run/033_p3_ppo_adaptive_fixed_weather_seed1_costw1_2__checkpoint_eval_curves.png}
    \caption{033 P3 Ppo Adaptive Fixed Weather Seed1 Costw1 2 - Checkpoint Eval Curves}
\end{figure}

\begin{figure}[H]
    \centering
    \includegraphics[width=0.8\textwidth]{figures/per_run/033_p3_ppo_adaptive_fixed_weather_seed1_costw1_2__diagnostics_panel.png}
    \caption{033 P3 Ppo Adaptive Fixed Weather Seed1 Costw1 2 - Diagnostics Panel}
\end{figure}

\begin{figure}[H]
    \centering
    \includegraphics[width=0.8\textwidth]{figures/per_run/033_p3_ppo_adaptive_fixed_weather_seed1_costw1_2__episode_length_vs_global_step.png}
    \caption{033 P3 Ppo Adaptive Fixed Weather Seed1 Costw1 2 - Episode Length Vs Global Step}
\end{figure}

\begin{figure}[H]
    \centering
    \includegraphics[width=0.8\textwidth]{figures/per_run/033_p3_ppo_adaptive_fixed_weather_seed1_costw1_2__primary_metric_vs_global_step.png}
    \caption{033 P3 Ppo Adaptive Fixed Weather Seed1 Costw1 2 - Primary Metric Vs Global Step}
\end{figure}

\begin{figure}[H]
    \centering
    \includegraphics[width=0.8\textwidth]{figures/per_run/033_p3_ppo_adaptive_fixed_weather_seed1_costw1_2__training_reward_vs_global_step.png}
    \caption{033 P3 Ppo Adaptive Fixed Weather Seed1 Costw1 2 - Training Reward Vs Global Step}
\end{figure}

\clearpage
\subsection{Run: 034 Fertilization Ppo Nonadaptive Fixed Weather Years 5000 Seed 2}
\begin{figure}[H]
    \centering
    \includegraphics[width=0.8\textwidth]{figures/per_run/034_fertilization_ppo_nonadaptive_fixed_weather_years_5000_seed_2__checkpoint_eval_curves.png}
    \caption{034 Fertilization Ppo Nonadaptive Fixed Weather Years 5000 Seed 2 - Checkpoint Eval Curves}
\end{figure}

\begin{figure}[H]
    \centering
    \includegraphics[width=0.8\textwidth]{figures/per_run/034_fertilization_ppo_nonadaptive_fixed_weather_years_5000_seed_2__diagnostics_panel.png}
    \caption{034 Fertilization Ppo Nonadaptive Fixed Weather Years 5000 Seed 2 - Diagnostics Panel}
\end{figure}

\begin{figure}[H]
    \centering
    \includegraphics[width=0.8\textwidth]{figures/per_run/034_fertilization_ppo_nonadaptive_fixed_weather_years_5000_seed_2__episode_length_vs_global_step.png}
    \caption{034 Fertilization Ppo Nonadaptive Fixed Weather Years 5000 Seed 2 - Episode Length Vs Global Step}
\end{figure}

\begin{figure}[H]
    \centering
    \includegraphics[width=0.8\textwidth]{figures/per_run/034_fertilization_ppo_nonadaptive_fixed_weather_years_5000_seed_2__primary_metric_vs_global_step.png}
    \caption{034 Fertilization Ppo Nonadaptive Fixed Weather Years 5000 Seed 2 - Primary Metric Vs Global Step}
\end{figure}

\begin{figure}[H]
    \centering
    \includegraphics[width=0.8\textwidth]{figures/per_run/034_fertilization_ppo_nonadaptive_fixed_weather_years_5000_seed_2__training_reward_vs_global_step.png}
    \caption{034 Fertilization Ppo Nonadaptive Fixed Weather Years 5000 Seed 2 - Training Reward Vs Global Step}
\end{figure}

\clearpage
\subsection{Run: 034 P3 Ppo Adaptive Fixed Weather Seed2 Costw0 8}
\begin{figure}[H]
    \centering
    \includegraphics[width=0.8\textwidth]{figures/per_run/034_p3_ppo_adaptive_fixed_weather_seed2_costw0_8__checkpoint_eval_curves.png}
    \caption{034 P3 Ppo Adaptive Fixed Weather Seed2 Costw0 8 - Checkpoint Eval Curves}
\end{figure}

\begin{figure}[H]
    \centering
    \includegraphics[width=0.8\textwidth]{figures/per_run/034_p3_ppo_adaptive_fixed_weather_seed2_costw0_8__diagnostics_panel.png}
    \caption{034 P3 Ppo Adaptive Fixed Weather Seed2 Costw0 8 - Diagnostics Panel}
\end{figure}

\begin{figure}[H]
    \centering
    \includegraphics[width=0.8\textwidth]{figures/per_run/034_p3_ppo_adaptive_fixed_weather_seed2_costw0_8__episode_length_vs_global_step.png}
    \caption{034 P3 Ppo Adaptive Fixed Weather Seed2 Costw0 8 - Episode Length Vs Global Step}
\end{figure}

\begin{figure}[H]
    \centering
    \includegraphics[width=0.8\textwidth]{figures/per_run/034_p3_ppo_adaptive_fixed_weather_seed2_costw0_8__primary_metric_vs_global_step.png}
    \caption{034 P3 Ppo Adaptive Fixed Weather Seed2 Costw0 8 - Primary Metric Vs Global Step}
\end{figure}

\begin{figure}[H]
    \centering
    \includegraphics[width=0.8\textwidth]{figures/per_run/034_p3_ppo_adaptive_fixed_weather_seed2_costw0_8__training_reward_vs_global_step.png}
    \caption{034 P3 Ppo Adaptive Fixed Weather Seed2 Costw0 8 - Training Reward Vs Global Step}
\end{figure}

\clearpage
\subsection{Run: 035 Fertilization Ppo Adaptive Random Weather Years 5000 Seed 2}
\begin{figure}[H]
    \centering
    \includegraphics[width=0.8\textwidth]{figures/per_run/035_fertilization_ppo_adaptive_random_weather_years_5000_seed_2__checkpoint_eval_curves.png}
    \caption{035 Fertilization Ppo Adaptive Random Weather Years 5000 Seed 2 - Checkpoint Eval Curves}
\end{figure}

\begin{figure}[H]
    \centering
    \includegraphics[width=0.8\textwidth]{figures/per_run/035_fertilization_ppo_adaptive_random_weather_years_5000_seed_2__diagnostics_panel.png}
    \caption{035 Fertilization Ppo Adaptive Random Weather Years 5000 Seed 2 - Diagnostics Panel}
\end{figure}

\begin{figure}[H]
    \centering
    \includegraphics[width=0.8\textwidth]{figures/per_run/035_fertilization_ppo_adaptive_random_weather_years_5000_seed_2__episode_length_vs_global_step.png}
    \caption{035 Fertilization Ppo Adaptive Random Weather Years 5000 Seed 2 - Episode Length Vs Global Step}
\end{figure}

\begin{figure}[H]
    \centering
    \includegraphics[width=0.8\textwidth]{figures/per_run/035_fertilization_ppo_adaptive_random_weather_years_5000_seed_2__primary_metric_vs_global_step.png}
    \caption{035 Fertilization Ppo Adaptive Random Weather Years 5000 Seed 2 - Primary Metric Vs Global Step}
\end{figure}

\begin{figure}[H]
    \centering
    \includegraphics[width=0.8\textwidth]{figures/per_run/035_fertilization_ppo_adaptive_random_weather_years_5000_seed_2__training_reward_vs_global_step.png}
    \caption{035 Fertilization Ppo Adaptive Random Weather Years 5000 Seed 2 - Training Reward Vs Global Step}
\end{figure}

\clearpage
\subsection{Run: 035 P3 Ppo Adaptive Fixed Weather Seed2 Costw1}
\begin{figure}[H]
    \centering
    \includegraphics[width=0.8\textwidth]{figures/per_run/035_p3_ppo_adaptive_fixed_weather_seed2_costw1__checkpoint_eval_curves.png}
    \caption{035 P3 Ppo Adaptive Fixed Weather Seed2 Costw1 - Checkpoint Eval Curves}
\end{figure}

\begin{figure}[H]
    \centering
    \includegraphics[width=0.8\textwidth]{figures/per_run/035_p3_ppo_adaptive_fixed_weather_seed2_costw1__diagnostics_panel.png}
    \caption{035 P3 Ppo Adaptive Fixed Weather Seed2 Costw1 - Diagnostics Panel}
\end{figure}

\begin{figure}[H]
    \centering
    \includegraphics[width=0.8\textwidth]{figures/per_run/035_p3_ppo_adaptive_fixed_weather_seed2_costw1__episode_length_vs_global_step.png}
    \caption{035 P3 Ppo Adaptive Fixed Weather Seed2 Costw1 - Episode Length Vs Global Step}
\end{figure}

\begin{figure}[H]
    \centering
    \includegraphics[width=0.8\textwidth]{figures/per_run/035_p3_ppo_adaptive_fixed_weather_seed2_costw1__primary_metric_vs_global_step.png}
    \caption{035 P3 Ppo Adaptive Fixed Weather Seed2 Costw1 - Primary Metric Vs Global Step}
\end{figure}

\begin{figure}[H]
    \centering
    \includegraphics[width=0.8\textwidth]{figures/per_run/035_p3_ppo_adaptive_fixed_weather_seed2_costw1__training_reward_vs_global_step.png}
    \caption{035 P3 Ppo Adaptive Fixed Weather Seed2 Costw1 - Training Reward Vs Global Step}
\end{figure}

\clearpage
\subsection{Run: 036 Fertilization Ppo Nonadaptive Random Weather Years 5000 Seed 2}
\begin{figure}[H]
    \centering
    \includegraphics[width=0.8\textwidth]{figures/per_run/036_fertilization_ppo_nonadaptive_random_weather_years_5000_seed_2__checkpoint_eval_curves.png}
    \caption{036 Fertilization Ppo Nonadaptive Random Weather Years 5000 Seed 2 - Checkpoint Eval Curves}
\end{figure}

\begin{figure}[H]
    \centering
    \includegraphics[width=0.8\textwidth]{figures/per_run/036_fertilization_ppo_nonadaptive_random_weather_years_5000_seed_2__diagnostics_panel.png}
    \caption{036 Fertilization Ppo Nonadaptive Random Weather Years 5000 Seed 2 - Diagnostics Panel}
\end{figure}

\begin{figure}[H]
    \centering
    \includegraphics[width=0.8\textwidth]{figures/per_run/036_fertilization_ppo_nonadaptive_random_weather_years_5000_seed_2__episode_length_vs_global_step.png}
    \caption{036 Fertilization Ppo Nonadaptive Random Weather Years 5000 Seed 2 - Episode Length Vs Global Step}
\end{figure}

\begin{figure}[H]
    \centering
    \includegraphics[width=0.8\textwidth]{figures/per_run/036_fertilization_ppo_nonadaptive_random_weather_years_5000_seed_2__primary_metric_vs_global_step.png}
    \caption{036 Fertilization Ppo Nonadaptive Random Weather Years 5000 Seed 2 - Primary Metric Vs Global Step}
\end{figure}

\begin{figure}[H]
    \centering
    \includegraphics[width=0.8\textwidth]{figures/per_run/036_fertilization_ppo_nonadaptive_random_weather_years_5000_seed_2__training_reward_vs_global_step.png}
    \caption{036 Fertilization Ppo Nonadaptive Random Weather Years 5000 Seed 2 - Training Reward Vs Global Step}
\end{figure}

\clearpage
\subsection{Run: 036 P3 Ppo Adaptive Fixed Weather Seed2 Costw1 2}
\begin{figure}[H]
    \centering
    \includegraphics[width=0.8\textwidth]{figures/per_run/036_p3_ppo_adaptive_fixed_weather_seed2_costw1_2__checkpoint_eval_curves.png}
    \caption{036 P3 Ppo Adaptive Fixed Weather Seed2 Costw1 2 - Checkpoint Eval Curves}
\end{figure}

\begin{figure}[H]
    \centering
    \includegraphics[width=0.8\textwidth]{figures/per_run/036_p3_ppo_adaptive_fixed_weather_seed2_costw1_2__diagnostics_panel.png}
    \caption{036 P3 Ppo Adaptive Fixed Weather Seed2 Costw1 2 - Diagnostics Panel}
\end{figure}

\begin{figure}[H]
    \centering
    \includegraphics[width=0.8\textwidth]{figures/per_run/036_p3_ppo_adaptive_fixed_weather_seed2_costw1_2__episode_length_vs_global_step.png}
    \caption{036 P3 Ppo Adaptive Fixed Weather Seed2 Costw1 2 - Episode Length Vs Global Step}
\end{figure}

\begin{figure}[H]
    \centering
    \includegraphics[width=0.8\textwidth]{figures/per_run/036_p3_ppo_adaptive_fixed_weather_seed2_costw1_2__primary_metric_vs_global_step.png}
    \caption{036 P3 Ppo Adaptive Fixed Weather Seed2 Costw1 2 - Primary Metric Vs Global Step}
\end{figure}

\begin{figure}[H]
    \centering
    \includegraphics[width=0.8\textwidth]{figures/per_run/036_p3_ppo_adaptive_fixed_weather_seed2_costw1_2__training_reward_vs_global_step.png}
    \caption{036 P3 Ppo Adaptive Fixed Weather Seed2 Costw1 2 - Training Reward Vs Global Step}
\end{figure}

\clearpage
\subsection{Run: 037 Fertilization A2C Adaptive Fixed Weather Years 1000 Seed 0}
\begin{figure}[H]
    \centering
    \includegraphics[width=0.8\textwidth]{figures/per_run/037_fertilization_a2c_adaptive_fixed_weather_years_1000_seed_0__checkpoint_eval_curves.png}
    \caption{037 Fertilization A2C Adaptive Fixed Weather Years 1000 Seed 0 - Checkpoint Eval Curves}
\end{figure}

\begin{figure}[H]
    \centering
    \includegraphics[width=0.8\textwidth]{figures/per_run/037_fertilization_a2c_adaptive_fixed_weather_years_1000_seed_0__diagnostics_panel.png}
    \caption{037 Fertilization A2C Adaptive Fixed Weather Years 1000 Seed 0 - Diagnostics Panel}
\end{figure}

\begin{figure}[H]
    \centering
    \includegraphics[width=0.8\textwidth]{figures/per_run/037_fertilization_a2c_adaptive_fixed_weather_years_1000_seed_0__episode_length_vs_global_step.png}
    \caption{037 Fertilization A2C Adaptive Fixed Weather Years 1000 Seed 0 - Episode Length Vs Global Step}
\end{figure}

\begin{figure}[H]
    \centering
    \includegraphics[width=0.8\textwidth]{figures/per_run/037_fertilization_a2c_adaptive_fixed_weather_years_1000_seed_0__primary_metric_vs_global_step.png}
    \caption{037 Fertilization A2C Adaptive Fixed Weather Years 1000 Seed 0 - Primary Metric Vs Global Step}
\end{figure}

\begin{figure}[H]
    \centering
    \includegraphics[width=0.8\textwidth]{figures/per_run/037_fertilization_a2c_adaptive_fixed_weather_years_1000_seed_0__training_reward_vs_global_step.png}
    \caption{037 Fertilization A2C Adaptive Fixed Weather Years 1000 Seed 0 - Training Reward Vs Global Step}
\end{figure}

\clearpage
\subsection{Run: 037 P3 Ppo Adaptive Random Weather Seed1 Costw0 8}
\begin{figure}[H]
    \centering
    \includegraphics[width=0.8\textwidth]{figures/per_run/037_p3_ppo_adaptive_random_weather_seed1_costw0_8__checkpoint_eval_curves.png}
    \caption{037 P3 Ppo Adaptive Random Weather Seed1 Costw0 8 - Checkpoint Eval Curves}
\end{figure}

\begin{figure}[H]
    \centering
    \includegraphics[width=0.8\textwidth]{figures/per_run/037_p3_ppo_adaptive_random_weather_seed1_costw0_8__diagnostics_panel.png}
    \caption{037 P3 Ppo Adaptive Random Weather Seed1 Costw0 8 - Diagnostics Panel}
\end{figure}

\begin{figure}[H]
    \centering
    \includegraphics[width=0.8\textwidth]{figures/per_run/037_p3_ppo_adaptive_random_weather_seed1_costw0_8__episode_length_vs_global_step.png}
    \caption{037 P3 Ppo Adaptive Random Weather Seed1 Costw0 8 - Episode Length Vs Global Step}
\end{figure}

\begin{figure}[H]
    \centering
    \includegraphics[width=0.8\textwidth]{figures/per_run/037_p3_ppo_adaptive_random_weather_seed1_costw0_8__primary_metric_vs_global_step.png}
    \caption{037 P3 Ppo Adaptive Random Weather Seed1 Costw0 8 - Primary Metric Vs Global Step}
\end{figure}

\begin{figure}[H]
    \centering
    \includegraphics[width=0.8\textwidth]{figures/per_run/037_p3_ppo_adaptive_random_weather_seed1_costw0_8__training_reward_vs_global_step.png}
    \caption{037 P3 Ppo Adaptive Random Weather Seed1 Costw0 8 - Training Reward Vs Global Step}
\end{figure}

\clearpage
\subsection{Run: 038 Fertilization A2C Nonadaptive Fixed Weather Years 1000 Seed 0}
\begin{figure}[H]
    \centering
    \includegraphics[width=0.8\textwidth]{figures/per_run/038_fertilization_a2c_nonadaptive_fixed_weather_years_1000_seed_0__checkpoint_eval_curves.png}
    \caption{038 Fertilization A2C Nonadaptive Fixed Weather Years 1000 Seed 0 - Checkpoint Eval Curves}
\end{figure}

\begin{figure}[H]
    \centering
    \includegraphics[width=0.8\textwidth]{figures/per_run/038_fertilization_a2c_nonadaptive_fixed_weather_years_1000_seed_0__diagnostics_panel.png}
    \caption{038 Fertilization A2C Nonadaptive Fixed Weather Years 1000 Seed 0 - Diagnostics Panel}
\end{figure}

\begin{figure}[H]
    \centering
    \includegraphics[width=0.8\textwidth]{figures/per_run/038_fertilization_a2c_nonadaptive_fixed_weather_years_1000_seed_0__episode_length_vs_global_step.png}
    \caption{038 Fertilization A2C Nonadaptive Fixed Weather Years 1000 Seed 0 - Episode Length Vs Global Step}
\end{figure}

\begin{figure}[H]
    \centering
    \includegraphics[width=0.8\textwidth]{figures/per_run/038_fertilization_a2c_nonadaptive_fixed_weather_years_1000_seed_0__primary_metric_vs_global_step.png}
    \caption{038 Fertilization A2C Nonadaptive Fixed Weather Years 1000 Seed 0 - Primary Metric Vs Global Step}
\end{figure}

\begin{figure}[H]
    \centering
    \includegraphics[width=0.8\textwidth]{figures/per_run/038_fertilization_a2c_nonadaptive_fixed_weather_years_1000_seed_0__training_reward_vs_global_step.png}
    \caption{038 Fertilization A2C Nonadaptive Fixed Weather Years 1000 Seed 0 - Training Reward Vs Global Step}
\end{figure}

\clearpage
\subsection{Run: 038 P3 Ppo Adaptive Random Weather Seed1 Costw1}
\begin{figure}[H]
    \centering
    \includegraphics[width=0.8\textwidth]{figures/per_run/038_p3_ppo_adaptive_random_weather_seed1_costw1__checkpoint_eval_curves.png}
    \caption{038 P3 Ppo Adaptive Random Weather Seed1 Costw1 - Checkpoint Eval Curves}
\end{figure}

\begin{figure}[H]
    \centering
    \includegraphics[width=0.8\textwidth]{figures/per_run/038_p3_ppo_adaptive_random_weather_seed1_costw1__diagnostics_panel.png}
    \caption{038 P3 Ppo Adaptive Random Weather Seed1 Costw1 - Diagnostics Panel}
\end{figure}

\begin{figure}[H]
    \centering
    \includegraphics[width=0.8\textwidth]{figures/per_run/038_p3_ppo_adaptive_random_weather_seed1_costw1__episode_length_vs_global_step.png}
    \caption{038 P3 Ppo Adaptive Random Weather Seed1 Costw1 - Episode Length Vs Global Step}
\end{figure}

\begin{figure}[H]
    \centering
    \includegraphics[width=0.8\textwidth]{figures/per_run/038_p3_ppo_adaptive_random_weather_seed1_costw1__primary_metric_vs_global_step.png}
    \caption{038 P3 Ppo Adaptive Random Weather Seed1 Costw1 - Primary Metric Vs Global Step}
\end{figure}

\begin{figure}[H]
    \centering
    \includegraphics[width=0.8\textwidth]{figures/per_run/038_p3_ppo_adaptive_random_weather_seed1_costw1__training_reward_vs_global_step.png}
    \caption{038 P3 Ppo Adaptive Random Weather Seed1 Costw1 - Training Reward Vs Global Step}
\end{figure}

\clearpage
\subsection{Run: 039 Fertilization A2C Adaptive Random Weather Years 1000 Seed 0}
\begin{figure}[H]
    \centering
    \includegraphics[width=0.8\textwidth]{figures/per_run/039_fertilization_a2c_adaptive_random_weather_years_1000_seed_0__checkpoint_eval_curves.png}
    \caption{039 Fertilization A2C Adaptive Random Weather Years 1000 Seed 0 - Checkpoint Eval Curves}
\end{figure}

\begin{figure}[H]
    \centering
    \includegraphics[width=0.8\textwidth]{figures/per_run/039_fertilization_a2c_adaptive_random_weather_years_1000_seed_0__diagnostics_panel.png}
    \caption{039 Fertilization A2C Adaptive Random Weather Years 1000 Seed 0 - Diagnostics Panel}
\end{figure}

\begin{figure}[H]
    \centering
    \includegraphics[width=0.8\textwidth]{figures/per_run/039_fertilization_a2c_adaptive_random_weather_years_1000_seed_0__episode_length_vs_global_step.png}
    \caption{039 Fertilization A2C Adaptive Random Weather Years 1000 Seed 0 - Episode Length Vs Global Step}
\end{figure}

\begin{figure}[H]
    \centering
    \includegraphics[width=0.8\textwidth]{figures/per_run/039_fertilization_a2c_adaptive_random_weather_years_1000_seed_0__primary_metric_vs_global_step.png}
    \caption{039 Fertilization A2C Adaptive Random Weather Years 1000 Seed 0 - Primary Metric Vs Global Step}
\end{figure}

\begin{figure}[H]
    \centering
    \includegraphics[width=0.8\textwidth]{figures/per_run/039_fertilization_a2c_adaptive_random_weather_years_1000_seed_0__training_reward_vs_global_step.png}
    \caption{039 Fertilization A2C Adaptive Random Weather Years 1000 Seed 0 - Training Reward Vs Global Step}
\end{figure}

\clearpage
\subsection{Run: 039 P3 Ppo Adaptive Random Weather Seed1 Costw1 2}
\begin{figure}[H]
    \centering
    \includegraphics[width=0.8\textwidth]{figures/per_run/039_p3_ppo_adaptive_random_weather_seed1_costw1_2__checkpoint_eval_curves.png}
    \caption{039 P3 Ppo Adaptive Random Weather Seed1 Costw1 2 - Checkpoint Eval Curves}
\end{figure}

\begin{figure}[H]
    \centering
    \includegraphics[width=0.8\textwidth]{figures/per_run/039_p3_ppo_adaptive_random_weather_seed1_costw1_2__diagnostics_panel.png}
    \caption{039 P3 Ppo Adaptive Random Weather Seed1 Costw1 2 - Diagnostics Panel}
\end{figure}

\begin{figure}[H]
    \centering
    \includegraphics[width=0.8\textwidth]{figures/per_run/039_p3_ppo_adaptive_random_weather_seed1_costw1_2__episode_length_vs_global_step.png}
    \caption{039 P3 Ppo Adaptive Random Weather Seed1 Costw1 2 - Episode Length Vs Global Step}
\end{figure}

\begin{figure}[H]
    \centering
    \includegraphics[width=0.8\textwidth]{figures/per_run/039_p3_ppo_adaptive_random_weather_seed1_costw1_2__primary_metric_vs_global_step.png}
    \caption{039 P3 Ppo Adaptive Random Weather Seed1 Costw1 2 - Primary Metric Vs Global Step}
\end{figure}

\begin{figure}[H]
    \centering
    \includegraphics[width=0.8\textwidth]{figures/per_run/039_p3_ppo_adaptive_random_weather_seed1_costw1_2__training_reward_vs_global_step.png}
    \caption{039 P3 Ppo Adaptive Random Weather Seed1 Costw1 2 - Training Reward Vs Global Step}
\end{figure}

\clearpage
\subsection{Run: 040 Fertilization A2C Nonadaptive Random Weather Years 1000 Seed 0}
\begin{figure}[H]
    \centering
    \includegraphics[width=0.8\textwidth]{figures/per_run/040_fertilization_a2c_nonadaptive_random_weather_years_1000_seed_0__checkpoint_eval_curves.png}
    \caption{040 Fertilization A2C Nonadaptive Random Weather Years 1000 Seed 0 - Checkpoint Eval Curves}
\end{figure}

\begin{figure}[H]
    \centering
    \includegraphics[width=0.8\textwidth]{figures/per_run/040_fertilization_a2c_nonadaptive_random_weather_years_1000_seed_0__diagnostics_panel.png}
    \caption{040 Fertilization A2C Nonadaptive Random Weather Years 1000 Seed 0 - Diagnostics Panel}
\end{figure}

\begin{figure}[H]
    \centering
    \includegraphics[width=0.8\textwidth]{figures/per_run/040_fertilization_a2c_nonadaptive_random_weather_years_1000_seed_0__episode_length_vs_global_step.png}
    \caption{040 Fertilization A2C Nonadaptive Random Weather Years 1000 Seed 0 - Episode Length Vs Global Step}
\end{figure}

\begin{figure}[H]
    \centering
    \includegraphics[width=0.8\textwidth]{figures/per_run/040_fertilization_a2c_nonadaptive_random_weather_years_1000_seed_0__primary_metric_vs_global_step.png}
    \caption{040 Fertilization A2C Nonadaptive Random Weather Years 1000 Seed 0 - Primary Metric Vs Global Step}
\end{figure}

\begin{figure}[H]
    \centering
    \includegraphics[width=0.8\textwidth]{figures/per_run/040_fertilization_a2c_nonadaptive_random_weather_years_1000_seed_0__training_reward_vs_global_step.png}
    \caption{040 Fertilization A2C Nonadaptive Random Weather Years 1000 Seed 0 - Training Reward Vs Global Step}
\end{figure}

\clearpage
\subsection{Run: 040 P3 Ppo Adaptive Random Weather Seed2 Costw0 8}
\begin{figure}[H]
    \centering
    \includegraphics[width=0.8\textwidth]{figures/per_run/040_p3_ppo_adaptive_random_weather_seed2_costw0_8__checkpoint_eval_curves.png}
    \caption{040 P3 Ppo Adaptive Random Weather Seed2 Costw0 8 - Checkpoint Eval Curves}
\end{figure}

\begin{figure}[H]
    \centering
    \includegraphics[width=0.8\textwidth]{figures/per_run/040_p3_ppo_adaptive_random_weather_seed2_costw0_8__diagnostics_panel.png}
    \caption{040 P3 Ppo Adaptive Random Weather Seed2 Costw0 8 - Diagnostics Panel}
\end{figure}

\begin{figure}[H]
    \centering
    \includegraphics[width=0.8\textwidth]{figures/per_run/040_p3_ppo_adaptive_random_weather_seed2_costw0_8__episode_length_vs_global_step.png}
    \caption{040 P3 Ppo Adaptive Random Weather Seed2 Costw0 8 - Episode Length Vs Global Step}
\end{figure}

\begin{figure}[H]
    \centering
    \includegraphics[width=0.8\textwidth]{figures/per_run/040_p3_ppo_adaptive_random_weather_seed2_costw0_8__primary_metric_vs_global_step.png}
    \caption{040 P3 Ppo Adaptive Random Weather Seed2 Costw0 8 - Primary Metric Vs Global Step}
\end{figure}

\begin{figure}[H]
    \centering
    \includegraphics[width=0.8\textwidth]{figures/per_run/040_p3_ppo_adaptive_random_weather_seed2_costw0_8__training_reward_vs_global_step.png}
    \caption{040 P3 Ppo Adaptive Random Weather Seed2 Costw0 8 - Training Reward Vs Global Step}
\end{figure}

\clearpage
\subsection{Run: 041 Fertilization A2C Adaptive Fixed Weather Years 1000 Seed 1}
\begin{figure}[H]
    \centering
    \includegraphics[width=0.8\textwidth]{figures/per_run/041_fertilization_a2c_adaptive_fixed_weather_years_1000_seed_1__checkpoint_eval_curves.png}
    \caption{041 Fertilization A2C Adaptive Fixed Weather Years 1000 Seed 1 - Checkpoint Eval Curves}
\end{figure}

\begin{figure}[H]
    \centering
    \includegraphics[width=0.8\textwidth]{figures/per_run/041_fertilization_a2c_adaptive_fixed_weather_years_1000_seed_1__diagnostics_panel.png}
    \caption{041 Fertilization A2C Adaptive Fixed Weather Years 1000 Seed 1 - Diagnostics Panel}
\end{figure}

\begin{figure}[H]
    \centering
    \includegraphics[width=0.8\textwidth]{figures/per_run/041_fertilization_a2c_adaptive_fixed_weather_years_1000_seed_1__episode_length_vs_global_step.png}
    \caption{041 Fertilization A2C Adaptive Fixed Weather Years 1000 Seed 1 - Episode Length Vs Global Step}
\end{figure}

\begin{figure}[H]
    \centering
    \includegraphics[width=0.8\textwidth]{figures/per_run/041_fertilization_a2c_adaptive_fixed_weather_years_1000_seed_1__primary_metric_vs_global_step.png}
    \caption{041 Fertilization A2C Adaptive Fixed Weather Years 1000 Seed 1 - Primary Metric Vs Global Step}
\end{figure}

\begin{figure}[H]
    \centering
    \includegraphics[width=0.8\textwidth]{figures/per_run/041_fertilization_a2c_adaptive_fixed_weather_years_1000_seed_1__training_reward_vs_global_step.png}
    \caption{041 Fertilization A2C Adaptive Fixed Weather Years 1000 Seed 1 - Training Reward Vs Global Step}
\end{figure}

\clearpage
\subsection{Run: 041 P3 Ppo Adaptive Random Weather Seed2 Costw1}
\begin{figure}[H]
    \centering
    \includegraphics[width=0.8\textwidth]{figures/per_run/041_p3_ppo_adaptive_random_weather_seed2_costw1__checkpoint_eval_curves.png}
    \caption{041 P3 Ppo Adaptive Random Weather Seed2 Costw1 - Checkpoint Eval Curves}
\end{figure}

\begin{figure}[H]
    \centering
    \includegraphics[width=0.8\textwidth]{figures/per_run/041_p3_ppo_adaptive_random_weather_seed2_costw1__diagnostics_panel.png}
    \caption{041 P3 Ppo Adaptive Random Weather Seed2 Costw1 - Diagnostics Panel}
\end{figure}

\begin{figure}[H]
    \centering
    \includegraphics[width=0.8\textwidth]{figures/per_run/041_p3_ppo_adaptive_random_weather_seed2_costw1__episode_length_vs_global_step.png}
    \caption{041 P3 Ppo Adaptive Random Weather Seed2 Costw1 - Episode Length Vs Global Step}
\end{figure}

\begin{figure}[H]
    \centering
    \includegraphics[width=0.8\textwidth]{figures/per_run/041_p3_ppo_adaptive_random_weather_seed2_costw1__primary_metric_vs_global_step.png}
    \caption{041 P3 Ppo Adaptive Random Weather Seed2 Costw1 - Primary Metric Vs Global Step}
\end{figure}

\begin{figure}[H]
    \centering
    \includegraphics[width=0.8\textwidth]{figures/per_run/041_p3_ppo_adaptive_random_weather_seed2_costw1__training_reward_vs_global_step.png}
    \caption{041 P3 Ppo Adaptive Random Weather Seed2 Costw1 - Training Reward Vs Global Step}
\end{figure}

\clearpage
\subsection{Run: 042 Fertilization A2C Nonadaptive Fixed Weather Years 1000 Seed 1}
\begin{figure}[H]
    \centering
    \includegraphics[width=0.8\textwidth]{figures/per_run/042_fertilization_a2c_nonadaptive_fixed_weather_years_1000_seed_1__checkpoint_eval_curves.png}
    \caption{042 Fertilization A2C Nonadaptive Fixed Weather Years 1000 Seed 1 - Checkpoint Eval Curves}
\end{figure}

\begin{figure}[H]
    \centering
    \includegraphics[width=0.8\textwidth]{figures/per_run/042_fertilization_a2c_nonadaptive_fixed_weather_years_1000_seed_1__diagnostics_panel.png}
    \caption{042 Fertilization A2C Nonadaptive Fixed Weather Years 1000 Seed 1 - Diagnostics Panel}
\end{figure}

\begin{figure}[H]
    \centering
    \includegraphics[width=0.8\textwidth]{figures/per_run/042_fertilization_a2c_nonadaptive_fixed_weather_years_1000_seed_1__episode_length_vs_global_step.png}
    \caption{042 Fertilization A2C Nonadaptive Fixed Weather Years 1000 Seed 1 - Episode Length Vs Global Step}
\end{figure}

\begin{figure}[H]
    \centering
    \includegraphics[width=0.8\textwidth]{figures/per_run/042_fertilization_a2c_nonadaptive_fixed_weather_years_1000_seed_1__primary_metric_vs_global_step.png}
    \caption{042 Fertilization A2C Nonadaptive Fixed Weather Years 1000 Seed 1 - Primary Metric Vs Global Step}
\end{figure}

\begin{figure}[H]
    \centering
    \includegraphics[width=0.8\textwidth]{figures/per_run/042_fertilization_a2c_nonadaptive_fixed_weather_years_1000_seed_1__training_reward_vs_global_step.png}
    \caption{042 Fertilization A2C Nonadaptive Fixed Weather Years 1000 Seed 1 - Training Reward Vs Global Step}
\end{figure}

\clearpage
\subsection{Run: 042 P3 Ppo Adaptive Random Weather Seed2 Costw1 2}
\begin{figure}[H]
    \centering
    \includegraphics[width=0.8\textwidth]{figures/per_run/042_p3_ppo_adaptive_random_weather_seed2_costw1_2__checkpoint_eval_curves.png}
    \caption{042 P3 Ppo Adaptive Random Weather Seed2 Costw1 2 - Checkpoint Eval Curves}
\end{figure}

\begin{figure}[H]
    \centering
    \includegraphics[width=0.8\textwidth]{figures/per_run/042_p3_ppo_adaptive_random_weather_seed2_costw1_2__diagnostics_panel.png}
    \caption{042 P3 Ppo Adaptive Random Weather Seed2 Costw1 2 - Diagnostics Panel}
\end{figure}

\begin{figure}[H]
    \centering
    \includegraphics[width=0.8\textwidth]{figures/per_run/042_p3_ppo_adaptive_random_weather_seed2_costw1_2__episode_length_vs_global_step.png}
    \caption{042 P3 Ppo Adaptive Random Weather Seed2 Costw1 2 - Episode Length Vs Global Step}
\end{figure}

\begin{figure}[H]
    \centering
    \includegraphics[width=0.8\textwidth]{figures/per_run/042_p3_ppo_adaptive_random_weather_seed2_costw1_2__primary_metric_vs_global_step.png}
    \caption{042 P3 Ppo Adaptive Random Weather Seed2 Costw1 2 - Primary Metric Vs Global Step}
\end{figure}

\begin{figure}[H]
    \centering
    \includegraphics[width=0.8\textwidth]{figures/per_run/042_p3_ppo_adaptive_random_weather_seed2_costw1_2__training_reward_vs_global_step.png}
    \caption{042 P3 Ppo Adaptive Random Weather Seed2 Costw1 2 - Training Reward Vs Global Step}
\end{figure}

\clearpage
\subsection{Run: 043 Fertilization A2C Adaptive Random Weather Years 1000 Seed 1}
\begin{figure}[H]
    \centering
    \includegraphics[width=0.8\textwidth]{figures/per_run/043_fertilization_a2c_adaptive_random_weather_years_1000_seed_1__checkpoint_eval_curves.png}
    \caption{043 Fertilization A2C Adaptive Random Weather Years 1000 Seed 1 - Checkpoint Eval Curves}
\end{figure}

\begin{figure}[H]
    \centering
    \includegraphics[width=0.8\textwidth]{figures/per_run/043_fertilization_a2c_adaptive_random_weather_years_1000_seed_1__diagnostics_panel.png}
    \caption{043 Fertilization A2C Adaptive Random Weather Years 1000 Seed 1 - Diagnostics Panel}
\end{figure}

\begin{figure}[H]
    \centering
    \includegraphics[width=0.8\textwidth]{figures/per_run/043_fertilization_a2c_adaptive_random_weather_years_1000_seed_1__episode_length_vs_global_step.png}
    \caption{043 Fertilization A2C Adaptive Random Weather Years 1000 Seed 1 - Episode Length Vs Global Step}
\end{figure}

\begin{figure}[H]
    \centering
    \includegraphics[width=0.8\textwidth]{figures/per_run/043_fertilization_a2c_adaptive_random_weather_years_1000_seed_1__primary_metric_vs_global_step.png}
    \caption{043 Fertilization A2C Adaptive Random Weather Years 1000 Seed 1 - Primary Metric Vs Global Step}
\end{figure}

\begin{figure}[H]
    \centering
    \includegraphics[width=0.8\textwidth]{figures/per_run/043_fertilization_a2c_adaptive_random_weather_years_1000_seed_1__training_reward_vs_global_step.png}
    \caption{043 Fertilization A2C Adaptive Random Weather Years 1000 Seed 1 - Training Reward Vs Global Step}
\end{figure}

\clearpage
\subsection{Run: 044 Fertilization A2C Nonadaptive Random Weather Years 1000 Seed 1}
\begin{figure}[H]
    \centering
    \includegraphics[width=0.8\textwidth]{figures/per_run/044_fertilization_a2c_nonadaptive_random_weather_years_1000_seed_1__checkpoint_eval_curves.png}
    \caption{044 Fertilization A2C Nonadaptive Random Weather Years 1000 Seed 1 - Checkpoint Eval Curves}
\end{figure}

\begin{figure}[H]
    \centering
    \includegraphics[width=0.8\textwidth]{figures/per_run/044_fertilization_a2c_nonadaptive_random_weather_years_1000_seed_1__diagnostics_panel.png}
    \caption{044 Fertilization A2C Nonadaptive Random Weather Years 1000 Seed 1 - Diagnostics Panel}
\end{figure}

\begin{figure}[H]
    \centering
    \includegraphics[width=0.8\textwidth]{figures/per_run/044_fertilization_a2c_nonadaptive_random_weather_years_1000_seed_1__episode_length_vs_global_step.png}
    \caption{044 Fertilization A2C Nonadaptive Random Weather Years 1000 Seed 1 - Episode Length Vs Global Step}
\end{figure}

\begin{figure}[H]
    \centering
    \includegraphics[width=0.8\textwidth]{figures/per_run/044_fertilization_a2c_nonadaptive_random_weather_years_1000_seed_1__primary_metric_vs_global_step.png}
    \caption{044 Fertilization A2C Nonadaptive Random Weather Years 1000 Seed 1 - Primary Metric Vs Global Step}
\end{figure}

\begin{figure}[H]
    \centering
    \includegraphics[width=0.8\textwidth]{figures/per_run/044_fertilization_a2c_nonadaptive_random_weather_years_1000_seed_1__training_reward_vs_global_step.png}
    \caption{044 Fertilization A2C Nonadaptive Random Weather Years 1000 Seed 1 - Training Reward Vs Global Step}
\end{figure}

\clearpage
\subsection{Run: 045 Fertilization A2C Adaptive Fixed Weather Years 1000 Seed 2}
\begin{figure}[H]
    \centering
    \includegraphics[width=0.8\textwidth]{figures/per_run/045_fertilization_a2c_adaptive_fixed_weather_years_1000_seed_2__checkpoint_eval_curves.png}
    \caption{045 Fertilization A2C Adaptive Fixed Weather Years 1000 Seed 2 - Checkpoint Eval Curves}
\end{figure}

\begin{figure}[H]
    \centering
    \includegraphics[width=0.8\textwidth]{figures/per_run/045_fertilization_a2c_adaptive_fixed_weather_years_1000_seed_2__diagnostics_panel.png}
    \caption{045 Fertilization A2C Adaptive Fixed Weather Years 1000 Seed 2 - Diagnostics Panel}
\end{figure}

\begin{figure}[H]
    \centering
    \includegraphics[width=0.8\textwidth]{figures/per_run/045_fertilization_a2c_adaptive_fixed_weather_years_1000_seed_2__episode_length_vs_global_step.png}
    \caption{045 Fertilization A2C Adaptive Fixed Weather Years 1000 Seed 2 - Episode Length Vs Global Step}
\end{figure}

\begin{figure}[H]
    \centering
    \includegraphics[width=0.8\textwidth]{figures/per_run/045_fertilization_a2c_adaptive_fixed_weather_years_1000_seed_2__primary_metric_vs_global_step.png}
    \caption{045 Fertilization A2C Adaptive Fixed Weather Years 1000 Seed 2 - Primary Metric Vs Global Step}
\end{figure}

\begin{figure}[H]
    \centering
    \includegraphics[width=0.8\textwidth]{figures/per_run/045_fertilization_a2c_adaptive_fixed_weather_years_1000_seed_2__training_reward_vs_global_step.png}
    \caption{045 Fertilization A2C Adaptive Fixed Weather Years 1000 Seed 2 - Training Reward Vs Global Step}
\end{figure}

\clearpage
\subsection{Run: 046 Fertilization A2C Nonadaptive Fixed Weather Years 1000 Seed 2}
\begin{figure}[H]
    \centering
    \includegraphics[width=0.8\textwidth]{figures/per_run/046_fertilization_a2c_nonadaptive_fixed_weather_years_1000_seed_2__checkpoint_eval_curves.png}
    \caption{046 Fertilization A2C Nonadaptive Fixed Weather Years 1000 Seed 2 - Checkpoint Eval Curves}
\end{figure}

\begin{figure}[H]
    \centering
    \includegraphics[width=0.8\textwidth]{figures/per_run/046_fertilization_a2c_nonadaptive_fixed_weather_years_1000_seed_2__diagnostics_panel.png}
    \caption{046 Fertilization A2C Nonadaptive Fixed Weather Years 1000 Seed 2 - Diagnostics Panel}
\end{figure}

\begin{figure}[H]
    \centering
    \includegraphics[width=0.8\textwidth]{figures/per_run/046_fertilization_a2c_nonadaptive_fixed_weather_years_1000_seed_2__episode_length_vs_global_step.png}
    \caption{046 Fertilization A2C Nonadaptive Fixed Weather Years 1000 Seed 2 - Episode Length Vs Global Step}
\end{figure}

\begin{figure}[H]
    \centering
    \includegraphics[width=0.8\textwidth]{figures/per_run/046_fertilization_a2c_nonadaptive_fixed_weather_years_1000_seed_2__primary_metric_vs_global_step.png}
    \caption{046 Fertilization A2C Nonadaptive Fixed Weather Years 1000 Seed 2 - Primary Metric Vs Global Step}
\end{figure}

\begin{figure}[H]
    \centering
    \includegraphics[width=0.8\textwidth]{figures/per_run/046_fertilization_a2c_nonadaptive_fixed_weather_years_1000_seed_2__training_reward_vs_global_step.png}
    \caption{046 Fertilization A2C Nonadaptive Fixed Weather Years 1000 Seed 2 - Training Reward Vs Global Step}
\end{figure}

\clearpage
\subsection{Run: 047 Fertilization A2C Adaptive Random Weather Years 1000 Seed 2}
\begin{figure}[H]
    \centering
    \includegraphics[width=0.8\textwidth]{figures/per_run/047_fertilization_a2c_adaptive_random_weather_years_1000_seed_2__checkpoint_eval_curves.png}
    \caption{047 Fertilization A2C Adaptive Random Weather Years 1000 Seed 2 - Checkpoint Eval Curves}
\end{figure}

\begin{figure}[H]
    \centering
    \includegraphics[width=0.8\textwidth]{figures/per_run/047_fertilization_a2c_adaptive_random_weather_years_1000_seed_2__diagnostics_panel.png}
    \caption{047 Fertilization A2C Adaptive Random Weather Years 1000 Seed 2 - Diagnostics Panel}
\end{figure}

\begin{figure}[H]
    \centering
    \includegraphics[width=0.8\textwidth]{figures/per_run/047_fertilization_a2c_adaptive_random_weather_years_1000_seed_2__episode_length_vs_global_step.png}
    \caption{047 Fertilization A2C Adaptive Random Weather Years 1000 Seed 2 - Episode Length Vs Global Step}
\end{figure}

\begin{figure}[H]
    \centering
    \includegraphics[width=0.8\textwidth]{figures/per_run/047_fertilization_a2c_adaptive_random_weather_years_1000_seed_2__primary_metric_vs_global_step.png}
    \caption{047 Fertilization A2C Adaptive Random Weather Years 1000 Seed 2 - Primary Metric Vs Global Step}
\end{figure}

\begin{figure}[H]
    \centering
    \includegraphics[width=0.8\textwidth]{figures/per_run/047_fertilization_a2c_adaptive_random_weather_years_1000_seed_2__training_reward_vs_global_step.png}
    \caption{047 Fertilization A2C Adaptive Random Weather Years 1000 Seed 2 - Training Reward Vs Global Step}
\end{figure}

\clearpage
\subsection{Run: 048 Fertilization A2C Nonadaptive Random Weather Years 1000 Seed 2}
\begin{figure}[H]
    \centering
    \includegraphics[width=0.8\textwidth]{figures/per_run/048_fertilization_a2c_nonadaptive_random_weather_years_1000_seed_2__checkpoint_eval_curves.png}
    \caption{048 Fertilization A2C Nonadaptive Random Weather Years 1000 Seed 2 - Checkpoint Eval Curves}
\end{figure}

\begin{figure}[H]
    \centering
    \includegraphics[width=0.8\textwidth]{figures/per_run/048_fertilization_a2c_nonadaptive_random_weather_years_1000_seed_2__diagnostics_panel.png}
    \caption{048 Fertilization A2C Nonadaptive Random Weather Years 1000 Seed 2 - Diagnostics Panel}
\end{figure}

\begin{figure}[H]
    \centering
    \includegraphics[width=0.8\textwidth]{figures/per_run/048_fertilization_a2c_nonadaptive_random_weather_years_1000_seed_2__episode_length_vs_global_step.png}
    \caption{048 Fertilization A2C Nonadaptive Random Weather Years 1000 Seed 2 - Episode Length Vs Global Step}
\end{figure}

\begin{figure}[H]
    \centering
    \includegraphics[width=0.8\textwidth]{figures/per_run/048_fertilization_a2c_nonadaptive_random_weather_years_1000_seed_2__primary_metric_vs_global_step.png}
    \caption{048 Fertilization A2C Nonadaptive Random Weather Years 1000 Seed 2 - Primary Metric Vs Global Step}
\end{figure}

\begin{figure}[H]
    \centering
    \includegraphics[width=0.8\textwidth]{figures/per_run/048_fertilization_a2c_nonadaptive_random_weather_years_1000_seed_2__training_reward_vs_global_step.png}
    \caption{048 Fertilization A2C Nonadaptive Random Weather Years 1000 Seed 2 - Training Reward Vs Global Step}
\end{figure}

\clearpage
\subsection{Run: 049 Fertilization A2C Adaptive Fixed Weather Years 3000 Seed 0}
\begin{figure}[H]
    \centering
    \includegraphics[width=0.8\textwidth]{figures/per_run/049_fertilization_a2c_adaptive_fixed_weather_years_3000_seed_0__checkpoint_eval_curves.png}
    \caption{049 Fertilization A2C Adaptive Fixed Weather Years 3000 Seed 0 - Checkpoint Eval Curves}
\end{figure}

\begin{figure}[H]
    \centering
    \includegraphics[width=0.8\textwidth]{figures/per_run/049_fertilization_a2c_adaptive_fixed_weather_years_3000_seed_0__diagnostics_panel.png}
    \caption{049 Fertilization A2C Adaptive Fixed Weather Years 3000 Seed 0 - Diagnostics Panel}
\end{figure}

\begin{figure}[H]
    \centering
    \includegraphics[width=0.8\textwidth]{figures/per_run/049_fertilization_a2c_adaptive_fixed_weather_years_3000_seed_0__episode_length_vs_global_step.png}
    \caption{049 Fertilization A2C Adaptive Fixed Weather Years 3000 Seed 0 - Episode Length Vs Global Step}
\end{figure}

\begin{figure}[H]
    \centering
    \includegraphics[width=0.8\textwidth]{figures/per_run/049_fertilization_a2c_adaptive_fixed_weather_years_3000_seed_0__primary_metric_vs_global_step.png}
    \caption{049 Fertilization A2C Adaptive Fixed Weather Years 3000 Seed 0 - Primary Metric Vs Global Step}
\end{figure}

\begin{figure}[H]
    \centering
    \includegraphics[width=0.8\textwidth]{figures/per_run/049_fertilization_a2c_adaptive_fixed_weather_years_3000_seed_0__training_reward_vs_global_step.png}
    \caption{049 Fertilization A2C Adaptive Fixed Weather Years 3000 Seed 0 - Training Reward Vs Global Step}
\end{figure}

\clearpage
\subsection{Run: 050 Fertilization A2C Nonadaptive Fixed Weather Years 3000 Seed 0}
\begin{figure}[H]
    \centering
    \includegraphics[width=0.8\textwidth]{figures/per_run/050_fertilization_a2c_nonadaptive_fixed_weather_years_3000_seed_0__checkpoint_eval_curves.png}
    \caption{050 Fertilization A2C Nonadaptive Fixed Weather Years 3000 Seed 0 - Checkpoint Eval Curves}
\end{figure}

\begin{figure}[H]
    \centering
    \includegraphics[width=0.8\textwidth]{figures/per_run/050_fertilization_a2c_nonadaptive_fixed_weather_years_3000_seed_0__diagnostics_panel.png}
    \caption{050 Fertilization A2C Nonadaptive Fixed Weather Years 3000 Seed 0 - Diagnostics Panel}
\end{figure}

\begin{figure}[H]
    \centering
    \includegraphics[width=0.8\textwidth]{figures/per_run/050_fertilization_a2c_nonadaptive_fixed_weather_years_3000_seed_0__episode_length_vs_global_step.png}
    \caption{050 Fertilization A2C Nonadaptive Fixed Weather Years 3000 Seed 0 - Episode Length Vs Global Step}
\end{figure}

\begin{figure}[H]
    \centering
    \includegraphics[width=0.8\textwidth]{figures/per_run/050_fertilization_a2c_nonadaptive_fixed_weather_years_3000_seed_0__primary_metric_vs_global_step.png}
    \caption{050 Fertilization A2C Nonadaptive Fixed Weather Years 3000 Seed 0 - Primary Metric Vs Global Step}
\end{figure}

\begin{figure}[H]
    \centering
    \includegraphics[width=0.8\textwidth]{figures/per_run/050_fertilization_a2c_nonadaptive_fixed_weather_years_3000_seed_0__training_reward_vs_global_step.png}
    \caption{050 Fertilization A2C Nonadaptive Fixed Weather Years 3000 Seed 0 - Training Reward Vs Global Step}
\end{figure}

\clearpage
\subsection{Run: 051 Fertilization A2C Adaptive Random Weather Years 3000 Seed 0}
\begin{figure}[H]
    \centering
    \includegraphics[width=0.8\textwidth]{figures/per_run/051_fertilization_a2c_adaptive_random_weather_years_3000_seed_0__checkpoint_eval_curves.png}
    \caption{051 Fertilization A2C Adaptive Random Weather Years 3000 Seed 0 - Checkpoint Eval Curves}
\end{figure}

\begin{figure}[H]
    \centering
    \includegraphics[width=0.8\textwidth]{figures/per_run/051_fertilization_a2c_adaptive_random_weather_years_3000_seed_0__diagnostics_panel.png}
    \caption{051 Fertilization A2C Adaptive Random Weather Years 3000 Seed 0 - Diagnostics Panel}
\end{figure}

\begin{figure}[H]
    \centering
    \includegraphics[width=0.8\textwidth]{figures/per_run/051_fertilization_a2c_adaptive_random_weather_years_3000_seed_0__episode_length_vs_global_step.png}
    \caption{051 Fertilization A2C Adaptive Random Weather Years 3000 Seed 0 - Episode Length Vs Global Step}
\end{figure}

\begin{figure}[H]
    \centering
    \includegraphics[width=0.8\textwidth]{figures/per_run/051_fertilization_a2c_adaptive_random_weather_years_3000_seed_0__primary_metric_vs_global_step.png}
    \caption{051 Fertilization A2C Adaptive Random Weather Years 3000 Seed 0 - Primary Metric Vs Global Step}
\end{figure}

\begin{figure}[H]
    \centering
    \includegraphics[width=0.8\textwidth]{figures/per_run/051_fertilization_a2c_adaptive_random_weather_years_3000_seed_0__training_reward_vs_global_step.png}
    \caption{051 Fertilization A2C Adaptive Random Weather Years 3000 Seed 0 - Training Reward Vs Global Step}
\end{figure}

\clearpage
\subsection{Run: 052 Fertilization A2C Nonadaptive Random Weather Years 3000 Seed 0}
\begin{figure}[H]
    \centering
    \includegraphics[width=0.8\textwidth]{figures/per_run/052_fertilization_a2c_nonadaptive_random_weather_years_3000_seed_0__checkpoint_eval_curves.png}
    \caption{052 Fertilization A2C Nonadaptive Random Weather Years 3000 Seed 0 - Checkpoint Eval Curves}
\end{figure}

\begin{figure}[H]
    \centering
    \includegraphics[width=0.8\textwidth]{figures/per_run/052_fertilization_a2c_nonadaptive_random_weather_years_3000_seed_0__diagnostics_panel.png}
    \caption{052 Fertilization A2C Nonadaptive Random Weather Years 3000 Seed 0 - Diagnostics Panel}
\end{figure}

\begin{figure}[H]
    \centering
    \includegraphics[width=0.8\textwidth]{figures/per_run/052_fertilization_a2c_nonadaptive_random_weather_years_3000_seed_0__episode_length_vs_global_step.png}
    \caption{052 Fertilization A2C Nonadaptive Random Weather Years 3000 Seed 0 - Episode Length Vs Global Step}
\end{figure}

\begin{figure}[H]
    \centering
    \includegraphics[width=0.8\textwidth]{figures/per_run/052_fertilization_a2c_nonadaptive_random_weather_years_3000_seed_0__primary_metric_vs_global_step.png}
    \caption{052 Fertilization A2C Nonadaptive Random Weather Years 3000 Seed 0 - Primary Metric Vs Global Step}
\end{figure}

\begin{figure}[H]
    \centering
    \includegraphics[width=0.8\textwidth]{figures/per_run/052_fertilization_a2c_nonadaptive_random_weather_years_3000_seed_0__training_reward_vs_global_step.png}
    \caption{052 Fertilization A2C Nonadaptive Random Weather Years 3000 Seed 0 - Training Reward Vs Global Step}
\end{figure}

\clearpage
\subsection{Run: 053 Fertilization A2C Adaptive Fixed Weather Years 3000 Seed 1}
\begin{figure}[H]
    \centering
    \includegraphics[width=0.8\textwidth]{figures/per_run/053_fertilization_a2c_adaptive_fixed_weather_years_3000_seed_1__checkpoint_eval_curves.png}
    \caption{053 Fertilization A2C Adaptive Fixed Weather Years 3000 Seed 1 - Checkpoint Eval Curves}
\end{figure}

\begin{figure}[H]
    \centering
    \includegraphics[width=0.8\textwidth]{figures/per_run/053_fertilization_a2c_adaptive_fixed_weather_years_3000_seed_1__diagnostics_panel.png}
    \caption{053 Fertilization A2C Adaptive Fixed Weather Years 3000 Seed 1 - Diagnostics Panel}
\end{figure}

\begin{figure}[H]
    \centering
    \includegraphics[width=0.8\textwidth]{figures/per_run/053_fertilization_a2c_adaptive_fixed_weather_years_3000_seed_1__episode_length_vs_global_step.png}
    \caption{053 Fertilization A2C Adaptive Fixed Weather Years 3000 Seed 1 - Episode Length Vs Global Step}
\end{figure}

\begin{figure}[H]
    \centering
    \includegraphics[width=0.8\textwidth]{figures/per_run/053_fertilization_a2c_adaptive_fixed_weather_years_3000_seed_1__primary_metric_vs_global_step.png}
    \caption{053 Fertilization A2C Adaptive Fixed Weather Years 3000 Seed 1 - Primary Metric Vs Global Step}
\end{figure}

\begin{figure}[H]
    \centering
    \includegraphics[width=0.8\textwidth]{figures/per_run/053_fertilization_a2c_adaptive_fixed_weather_years_3000_seed_1__training_reward_vs_global_step.png}
    \caption{053 Fertilization A2C Adaptive Fixed Weather Years 3000 Seed 1 - Training Reward Vs Global Step}
\end{figure}

\clearpage
\subsection{Run: 054 Fertilization A2C Nonadaptive Fixed Weather Years 3000 Seed 1}
\begin{figure}[H]
    \centering
    \includegraphics[width=0.8\textwidth]{figures/per_run/054_fertilization_a2c_nonadaptive_fixed_weather_years_3000_seed_1__checkpoint_eval_curves.png}
    \caption{054 Fertilization A2C Nonadaptive Fixed Weather Years 3000 Seed 1 - Checkpoint Eval Curves}
\end{figure}

\begin{figure}[H]
    \centering
    \includegraphics[width=0.8\textwidth]{figures/per_run/054_fertilization_a2c_nonadaptive_fixed_weather_years_3000_seed_1__diagnostics_panel.png}
    \caption{054 Fertilization A2C Nonadaptive Fixed Weather Years 3000 Seed 1 - Diagnostics Panel}
\end{figure}

\begin{figure}[H]
    \centering
    \includegraphics[width=0.8\textwidth]{figures/per_run/054_fertilization_a2c_nonadaptive_fixed_weather_years_3000_seed_1__episode_length_vs_global_step.png}
    \caption{054 Fertilization A2C Nonadaptive Fixed Weather Years 3000 Seed 1 - Episode Length Vs Global Step}
\end{figure}

\begin{figure}[H]
    \centering
    \includegraphics[width=0.8\textwidth]{figures/per_run/054_fertilization_a2c_nonadaptive_fixed_weather_years_3000_seed_1__primary_metric_vs_global_step.png}
    \caption{054 Fertilization A2C Nonadaptive Fixed Weather Years 3000 Seed 1 - Primary Metric Vs Global Step}
\end{figure}

\begin{figure}[H]
    \centering
    \includegraphics[width=0.8\textwidth]{figures/per_run/054_fertilization_a2c_nonadaptive_fixed_weather_years_3000_seed_1__training_reward_vs_global_step.png}
    \caption{054 Fertilization A2C Nonadaptive Fixed Weather Years 3000 Seed 1 - Training Reward Vs Global Step}
\end{figure}

\clearpage
\subsection{Run: 055 Fertilization A2C Adaptive Random Weather Years 3000 Seed 1}
\begin{figure}[H]
    \centering
    \includegraphics[width=0.8\textwidth]{figures/per_run/055_fertilization_a2c_adaptive_random_weather_years_3000_seed_1__checkpoint_eval_curves.png}
    \caption{055 Fertilization A2C Adaptive Random Weather Years 3000 Seed 1 - Checkpoint Eval Curves}
\end{figure}

\begin{figure}[H]
    \centering
    \includegraphics[width=0.8\textwidth]{figures/per_run/055_fertilization_a2c_adaptive_random_weather_years_3000_seed_1__diagnostics_panel.png}
    \caption{055 Fertilization A2C Adaptive Random Weather Years 3000 Seed 1 - Diagnostics Panel}
\end{figure}

\begin{figure}[H]
    \centering
    \includegraphics[width=0.8\textwidth]{figures/per_run/055_fertilization_a2c_adaptive_random_weather_years_3000_seed_1__episode_length_vs_global_step.png}
    \caption{055 Fertilization A2C Adaptive Random Weather Years 3000 Seed 1 - Episode Length Vs Global Step}
\end{figure}

\begin{figure}[H]
    \centering
    \includegraphics[width=0.8\textwidth]{figures/per_run/055_fertilization_a2c_adaptive_random_weather_years_3000_seed_1__primary_metric_vs_global_step.png}
    \caption{055 Fertilization A2C Adaptive Random Weather Years 3000 Seed 1 - Primary Metric Vs Global Step}
\end{figure}

\begin{figure}[H]
    \centering
    \includegraphics[width=0.8\textwidth]{figures/per_run/055_fertilization_a2c_adaptive_random_weather_years_3000_seed_1__training_reward_vs_global_step.png}
    \caption{055 Fertilization A2C Adaptive Random Weather Years 3000 Seed 1 - Training Reward Vs Global Step}
\end{figure}

\clearpage
\subsection{Run: 056 Fertilization A2C Nonadaptive Random Weather Years 3000 Seed 1}
\begin{figure}[H]
    \centering
    \includegraphics[width=0.8\textwidth]{figures/per_run/056_fertilization_a2c_nonadaptive_random_weather_years_3000_seed_1__checkpoint_eval_curves.png}
    \caption{056 Fertilization A2C Nonadaptive Random Weather Years 3000 Seed 1 - Checkpoint Eval Curves}
\end{figure}

\begin{figure}[H]
    \centering
    \includegraphics[width=0.8\textwidth]{figures/per_run/056_fertilization_a2c_nonadaptive_random_weather_years_3000_seed_1__diagnostics_panel.png}
    \caption{056 Fertilization A2C Nonadaptive Random Weather Years 3000 Seed 1 - Diagnostics Panel}
\end{figure}

\begin{figure}[H]
    \centering
    \includegraphics[width=0.8\textwidth]{figures/per_run/056_fertilization_a2c_nonadaptive_random_weather_years_3000_seed_1__episode_length_vs_global_step.png}
    \caption{056 Fertilization A2C Nonadaptive Random Weather Years 3000 Seed 1 - Episode Length Vs Global Step}
\end{figure}

\begin{figure}[H]
    \centering
    \includegraphics[width=0.8\textwidth]{figures/per_run/056_fertilization_a2c_nonadaptive_random_weather_years_3000_seed_1__primary_metric_vs_global_step.png}
    \caption{056 Fertilization A2C Nonadaptive Random Weather Years 3000 Seed 1 - Primary Metric Vs Global Step}
\end{figure}

\begin{figure}[H]
    \centering
    \includegraphics[width=0.8\textwidth]{figures/per_run/056_fertilization_a2c_nonadaptive_random_weather_years_3000_seed_1__training_reward_vs_global_step.png}
    \caption{056 Fertilization A2C Nonadaptive Random Weather Years 3000 Seed 1 - Training Reward Vs Global Step}
\end{figure}

\clearpage
\subsection{Run: 057 Fertilization A2C Adaptive Fixed Weather Years 3000 Seed 2}
\begin{figure}[H]
    \centering
    \includegraphics[width=0.8\textwidth]{figures/per_run/057_fertilization_a2c_adaptive_fixed_weather_years_3000_seed_2__checkpoint_eval_curves.png}
    \caption{057 Fertilization A2C Adaptive Fixed Weather Years 3000 Seed 2 - Checkpoint Eval Curves}
\end{figure}

\begin{figure}[H]
    \centering
    \includegraphics[width=0.8\textwidth]{figures/per_run/057_fertilization_a2c_adaptive_fixed_weather_years_3000_seed_2__diagnostics_panel.png}
    \caption{057 Fertilization A2C Adaptive Fixed Weather Years 3000 Seed 2 - Diagnostics Panel}
\end{figure}

\begin{figure}[H]
    \centering
    \includegraphics[width=0.8\textwidth]{figures/per_run/057_fertilization_a2c_adaptive_fixed_weather_years_3000_seed_2__episode_length_vs_global_step.png}
    \caption{057 Fertilization A2C Adaptive Fixed Weather Years 3000 Seed 2 - Episode Length Vs Global Step}
\end{figure}

\begin{figure}[H]
    \centering
    \includegraphics[width=0.8\textwidth]{figures/per_run/057_fertilization_a2c_adaptive_fixed_weather_years_3000_seed_2__primary_metric_vs_global_step.png}
    \caption{057 Fertilization A2C Adaptive Fixed Weather Years 3000 Seed 2 - Primary Metric Vs Global Step}
\end{figure}

\begin{figure}[H]
    \centering
    \includegraphics[width=0.8\textwidth]{figures/per_run/057_fertilization_a2c_adaptive_fixed_weather_years_3000_seed_2__training_reward_vs_global_step.png}
    \caption{057 Fertilization A2C Adaptive Fixed Weather Years 3000 Seed 2 - Training Reward Vs Global Step}
\end{figure}

\clearpage
\subsection{Run: 058 Fertilization A2C Nonadaptive Fixed Weather Years 3000 Seed 2}
\begin{figure}[H]
    \centering
    \includegraphics[width=0.8\textwidth]{figures/per_run/058_fertilization_a2c_nonadaptive_fixed_weather_years_3000_seed_2__checkpoint_eval_curves.png}
    \caption{058 Fertilization A2C Nonadaptive Fixed Weather Years 3000 Seed 2 - Checkpoint Eval Curves}
\end{figure}

\begin{figure}[H]
    \centering
    \includegraphics[width=0.8\textwidth]{figures/per_run/058_fertilization_a2c_nonadaptive_fixed_weather_years_3000_seed_2__diagnostics_panel.png}
    \caption{058 Fertilization A2C Nonadaptive Fixed Weather Years 3000 Seed 2 - Diagnostics Panel}
\end{figure}

\begin{figure}[H]
    \centering
    \includegraphics[width=0.8\textwidth]{figures/per_run/058_fertilization_a2c_nonadaptive_fixed_weather_years_3000_seed_2__episode_length_vs_global_step.png}
    \caption{058 Fertilization A2C Nonadaptive Fixed Weather Years 3000 Seed 2 - Episode Length Vs Global Step}
\end{figure}

\begin{figure}[H]
    \centering
    \includegraphics[width=0.8\textwidth]{figures/per_run/058_fertilization_a2c_nonadaptive_fixed_weather_years_3000_seed_2__primary_metric_vs_global_step.png}
    \caption{058 Fertilization A2C Nonadaptive Fixed Weather Years 3000 Seed 2 - Primary Metric Vs Global Step}
\end{figure}

\begin{figure}[H]
    \centering
    \includegraphics[width=0.8\textwidth]{figures/per_run/058_fertilization_a2c_nonadaptive_fixed_weather_years_3000_seed_2__training_reward_vs_global_step.png}
    \caption{058 Fertilization A2C Nonadaptive Fixed Weather Years 3000 Seed 2 - Training Reward Vs Global Step}
\end{figure}

\clearpage
\subsection{Run: 059 Fertilization A2C Adaptive Random Weather Years 3000 Seed 2}
\begin{figure}[H]
    \centering
    \includegraphics[width=0.8\textwidth]{figures/per_run/059_fertilization_a2c_adaptive_random_weather_years_3000_seed_2__checkpoint_eval_curves.png}
    \caption{059 Fertilization A2C Adaptive Random Weather Years 3000 Seed 2 - Checkpoint Eval Curves}
\end{figure}

\begin{figure}[H]
    \centering
    \includegraphics[width=0.8\textwidth]{figures/per_run/059_fertilization_a2c_adaptive_random_weather_years_3000_seed_2__diagnostics_panel.png}
    \caption{059 Fertilization A2C Adaptive Random Weather Years 3000 Seed 2 - Diagnostics Panel}
\end{figure}

\begin{figure}[H]
    \centering
    \includegraphics[width=0.8\textwidth]{figures/per_run/059_fertilization_a2c_adaptive_random_weather_years_3000_seed_2__episode_length_vs_global_step.png}
    \caption{059 Fertilization A2C Adaptive Random Weather Years 3000 Seed 2 - Episode Length Vs Global Step}
\end{figure}

\begin{figure}[H]
    \centering
    \includegraphics[width=0.8\textwidth]{figures/per_run/059_fertilization_a2c_adaptive_random_weather_years_3000_seed_2__primary_metric_vs_global_step.png}
    \caption{059 Fertilization A2C Adaptive Random Weather Years 3000 Seed 2 - Primary Metric Vs Global Step}
\end{figure}

\begin{figure}[H]
    \centering
    \includegraphics[width=0.8\textwidth]{figures/per_run/059_fertilization_a2c_adaptive_random_weather_years_3000_seed_2__training_reward_vs_global_step.png}
    \caption{059 Fertilization A2C Adaptive Random Weather Years 3000 Seed 2 - Training Reward Vs Global Step}
\end{figure}

\clearpage
\subsection{Run: 060 Fertilization A2C Nonadaptive Random Weather Years 3000 Seed 2}
\begin{figure}[H]
    \centering
    \includegraphics[width=0.8\textwidth]{figures/per_run/060_fertilization_a2c_nonadaptive_random_weather_years_3000_seed_2__checkpoint_eval_curves.png}
    \caption{060 Fertilization A2C Nonadaptive Random Weather Years 3000 Seed 2 - Checkpoint Eval Curves}
\end{figure}

\begin{figure}[H]
    \centering
    \includegraphics[width=0.8\textwidth]{figures/per_run/060_fertilization_a2c_nonadaptive_random_weather_years_3000_seed_2__diagnostics_panel.png}
    \caption{060 Fertilization A2C Nonadaptive Random Weather Years 3000 Seed 2 - Diagnostics Panel}
\end{figure}

\begin{figure}[H]
    \centering
    \includegraphics[width=0.8\textwidth]{figures/per_run/060_fertilization_a2c_nonadaptive_random_weather_years_3000_seed_2__episode_length_vs_global_step.png}
    \caption{060 Fertilization A2C Nonadaptive Random Weather Years 3000 Seed 2 - Episode Length Vs Global Step}
\end{figure}

\begin{figure}[H]
    \centering
    \includegraphics[width=0.8\textwidth]{figures/per_run/060_fertilization_a2c_nonadaptive_random_weather_years_3000_seed_2__primary_metric_vs_global_step.png}
    \caption{060 Fertilization A2C Nonadaptive Random Weather Years 3000 Seed 2 - Primary Metric Vs Global Step}
\end{figure}

\begin{figure}[H]
    \centering
    \includegraphics[width=0.8\textwidth]{figures/per_run/060_fertilization_a2c_nonadaptive_random_weather_years_3000_seed_2__training_reward_vs_global_step.png}
    \caption{060 Fertilization A2C Nonadaptive Random Weather Years 3000 Seed 2 - Training Reward Vs Global Step}
\end{figure}

\clearpage
\subsection{Run: 061 Fertilization A2C Adaptive Fixed Weather Years 5000 Seed 0}
\begin{figure}[H]
    \centering
    \includegraphics[width=0.8\textwidth]{figures/per_run/061_fertilization_a2c_adaptive_fixed_weather_years_5000_seed_0__checkpoint_eval_curves.png}
    \caption{061 Fertilization A2C Adaptive Fixed Weather Years 5000 Seed 0 - Checkpoint Eval Curves}
\end{figure}

\begin{figure}[H]
    \centering
    \includegraphics[width=0.8\textwidth]{figures/per_run/061_fertilization_a2c_adaptive_fixed_weather_years_5000_seed_0__diagnostics_panel.png}
    \caption{061 Fertilization A2C Adaptive Fixed Weather Years 5000 Seed 0 - Diagnostics Panel}
\end{figure}

\begin{figure}[H]
    \centering
    \includegraphics[width=0.8\textwidth]{figures/per_run/061_fertilization_a2c_adaptive_fixed_weather_years_5000_seed_0__episode_length_vs_global_step.png}
    \caption{061 Fertilization A2C Adaptive Fixed Weather Years 5000 Seed 0 - Episode Length Vs Global Step}
\end{figure}

\begin{figure}[H]
    \centering
    \includegraphics[width=0.8\textwidth]{figures/per_run/061_fertilization_a2c_adaptive_fixed_weather_years_5000_seed_0__primary_metric_vs_global_step.png}
    \caption{061 Fertilization A2C Adaptive Fixed Weather Years 5000 Seed 0 - Primary Metric Vs Global Step}
\end{figure}

\begin{figure}[H]
    \centering
    \includegraphics[width=0.8\textwidth]{figures/per_run/061_fertilization_a2c_adaptive_fixed_weather_years_5000_seed_0__training_reward_vs_global_step.png}
    \caption{061 Fertilization A2C Adaptive Fixed Weather Years 5000 Seed 0 - Training Reward Vs Global Step}
\end{figure}

\clearpage
\subsection{Run: 062 Fertilization A2C Nonadaptive Fixed Weather Years 5000 Seed 0}
\begin{figure}[H]
    \centering
    \includegraphics[width=0.8\textwidth]{figures/per_run/062_fertilization_a2c_nonadaptive_fixed_weather_years_5000_seed_0__checkpoint_eval_curves.png}
    \caption{062 Fertilization A2C Nonadaptive Fixed Weather Years 5000 Seed 0 - Checkpoint Eval Curves}
\end{figure}

\begin{figure}[H]
    \centering
    \includegraphics[width=0.8\textwidth]{figures/per_run/062_fertilization_a2c_nonadaptive_fixed_weather_years_5000_seed_0__diagnostics_panel.png}
    \caption{062 Fertilization A2C Nonadaptive Fixed Weather Years 5000 Seed 0 - Diagnostics Panel}
\end{figure}

\begin{figure}[H]
    \centering
    \includegraphics[width=0.8\textwidth]{figures/per_run/062_fertilization_a2c_nonadaptive_fixed_weather_years_5000_seed_0__episode_length_vs_global_step.png}
    \caption{062 Fertilization A2C Nonadaptive Fixed Weather Years 5000 Seed 0 - Episode Length Vs Global Step}
\end{figure}

\begin{figure}[H]
    \centering
    \includegraphics[width=0.8\textwidth]{figures/per_run/062_fertilization_a2c_nonadaptive_fixed_weather_years_5000_seed_0__primary_metric_vs_global_step.png}
    \caption{062 Fertilization A2C Nonadaptive Fixed Weather Years 5000 Seed 0 - Primary Metric Vs Global Step}
\end{figure}

\begin{figure}[H]
    \centering
    \includegraphics[width=0.8\textwidth]{figures/per_run/062_fertilization_a2c_nonadaptive_fixed_weather_years_5000_seed_0__training_reward_vs_global_step.png}
    \caption{062 Fertilization A2C Nonadaptive Fixed Weather Years 5000 Seed 0 - Training Reward Vs Global Step}
\end{figure}

\clearpage
\subsection{Run: 063 Fertilization A2C Adaptive Random Weather Years 5000 Seed 0}
\begin{figure}[H]
    \centering
    \includegraphics[width=0.8\textwidth]{figures/per_run/063_fertilization_a2c_adaptive_random_weather_years_5000_seed_0__checkpoint_eval_curves.png}
    \caption{063 Fertilization A2C Adaptive Random Weather Years 5000 Seed 0 - Checkpoint Eval Curves}
\end{figure}

\begin{figure}[H]
    \centering
    \includegraphics[width=0.8\textwidth]{figures/per_run/063_fertilization_a2c_adaptive_random_weather_years_5000_seed_0__diagnostics_panel.png}
    \caption{063 Fertilization A2C Adaptive Random Weather Years 5000 Seed 0 - Diagnostics Panel}
\end{figure}

\begin{figure}[H]
    \centering
    \includegraphics[width=0.8\textwidth]{figures/per_run/063_fertilization_a2c_adaptive_random_weather_years_5000_seed_0__episode_length_vs_global_step.png}
    \caption{063 Fertilization A2C Adaptive Random Weather Years 5000 Seed 0 - Episode Length Vs Global Step}
\end{figure}

\begin{figure}[H]
    \centering
    \includegraphics[width=0.8\textwidth]{figures/per_run/063_fertilization_a2c_adaptive_random_weather_years_5000_seed_0__primary_metric_vs_global_step.png}
    \caption{063 Fertilization A2C Adaptive Random Weather Years 5000 Seed 0 - Primary Metric Vs Global Step}
\end{figure}

\begin{figure}[H]
    \centering
    \includegraphics[width=0.8\textwidth]{figures/per_run/063_fertilization_a2c_adaptive_random_weather_years_5000_seed_0__training_reward_vs_global_step.png}
    \caption{063 Fertilization A2C Adaptive Random Weather Years 5000 Seed 0 - Training Reward Vs Global Step}
\end{figure}

\clearpage
\subsection{Run: 064 Fertilization A2C Nonadaptive Random Weather Years 5000 Seed 0}
\begin{figure}[H]
    \centering
    \includegraphics[width=0.8\textwidth]{figures/per_run/064_fertilization_a2c_nonadaptive_random_weather_years_5000_seed_0__checkpoint_eval_curves.png}
    \caption{064 Fertilization A2C Nonadaptive Random Weather Years 5000 Seed 0 - Checkpoint Eval Curves}
\end{figure}

\begin{figure}[H]
    \centering
    \includegraphics[width=0.8\textwidth]{figures/per_run/064_fertilization_a2c_nonadaptive_random_weather_years_5000_seed_0__diagnostics_panel.png}
    \caption{064 Fertilization A2C Nonadaptive Random Weather Years 5000 Seed 0 - Diagnostics Panel}
\end{figure}

\begin{figure}[H]
    \centering
    \includegraphics[width=0.8\textwidth]{figures/per_run/064_fertilization_a2c_nonadaptive_random_weather_years_5000_seed_0__episode_length_vs_global_step.png}
    \caption{064 Fertilization A2C Nonadaptive Random Weather Years 5000 Seed 0 - Episode Length Vs Global Step}
\end{figure}

\begin{figure}[H]
    \centering
    \includegraphics[width=0.8\textwidth]{figures/per_run/064_fertilization_a2c_nonadaptive_random_weather_years_5000_seed_0__primary_metric_vs_global_step.png}
    \caption{064 Fertilization A2C Nonadaptive Random Weather Years 5000 Seed 0 - Primary Metric Vs Global Step}
\end{figure}

\begin{figure}[H]
    \centering
    \includegraphics[width=0.8\textwidth]{figures/per_run/064_fertilization_a2c_nonadaptive_random_weather_years_5000_seed_0__training_reward_vs_global_step.png}
    \caption{064 Fertilization A2C Nonadaptive Random Weather Years 5000 Seed 0 - Training Reward Vs Global Step}
\end{figure}

\clearpage
\subsection{Run: 065 Fertilization A2C Adaptive Fixed Weather Years 5000 Seed 1}
\begin{figure}[H]
    \centering
    \includegraphics[width=0.8\textwidth]{figures/per_run/065_fertilization_a2c_adaptive_fixed_weather_years_5000_seed_1__checkpoint_eval_curves.png}
    \caption{065 Fertilization A2C Adaptive Fixed Weather Years 5000 Seed 1 - Checkpoint Eval Curves}
\end{figure}

\begin{figure}[H]
    \centering
    \includegraphics[width=0.8\textwidth]{figures/per_run/065_fertilization_a2c_adaptive_fixed_weather_years_5000_seed_1__diagnostics_panel.png}
    \caption{065 Fertilization A2C Adaptive Fixed Weather Years 5000 Seed 1 - Diagnostics Panel}
\end{figure}

\begin{figure}[H]
    \centering
    \includegraphics[width=0.8\textwidth]{figures/per_run/065_fertilization_a2c_adaptive_fixed_weather_years_5000_seed_1__episode_length_vs_global_step.png}
    \caption{065 Fertilization A2C Adaptive Fixed Weather Years 5000 Seed 1 - Episode Length Vs Global Step}
\end{figure}

\begin{figure}[H]
    \centering
    \includegraphics[width=0.8\textwidth]{figures/per_run/065_fertilization_a2c_adaptive_fixed_weather_years_5000_seed_1__primary_metric_vs_global_step.png}
    \caption{065 Fertilization A2C Adaptive Fixed Weather Years 5000 Seed 1 - Primary Metric Vs Global Step}
\end{figure}

\begin{figure}[H]
    \centering
    \includegraphics[width=0.8\textwidth]{figures/per_run/065_fertilization_a2c_adaptive_fixed_weather_years_5000_seed_1__training_reward_vs_global_step.png}
    \caption{065 Fertilization A2C Adaptive Fixed Weather Years 5000 Seed 1 - Training Reward Vs Global Step}
\end{figure}

\clearpage
\subsection{Run: 066 Fertilization A2C Nonadaptive Fixed Weather Years 5000 Seed 1}
\begin{figure}[H]
    \centering
    \includegraphics[width=0.8\textwidth]{figures/per_run/066_fertilization_a2c_nonadaptive_fixed_weather_years_5000_seed_1__checkpoint_eval_curves.png}
    \caption{066 Fertilization A2C Nonadaptive Fixed Weather Years 5000 Seed 1 - Checkpoint Eval Curves}
\end{figure}

\begin{figure}[H]
    \centering
    \includegraphics[width=0.8\textwidth]{figures/per_run/066_fertilization_a2c_nonadaptive_fixed_weather_years_5000_seed_1__diagnostics_panel.png}
    \caption{066 Fertilization A2C Nonadaptive Fixed Weather Years 5000 Seed 1 - Diagnostics Panel}
\end{figure}

\begin{figure}[H]
    \centering
    \includegraphics[width=0.8\textwidth]{figures/per_run/066_fertilization_a2c_nonadaptive_fixed_weather_years_5000_seed_1__episode_length_vs_global_step.png}
    \caption{066 Fertilization A2C Nonadaptive Fixed Weather Years 5000 Seed 1 - Episode Length Vs Global Step}
\end{figure}

\begin{figure}[H]
    \centering
    \includegraphics[width=0.8\textwidth]{figures/per_run/066_fertilization_a2c_nonadaptive_fixed_weather_years_5000_seed_1__primary_metric_vs_global_step.png}
    \caption{066 Fertilization A2C Nonadaptive Fixed Weather Years 5000 Seed 1 - Primary Metric Vs Global Step}
\end{figure}

\begin{figure}[H]
    \centering
    \includegraphics[width=0.8\textwidth]{figures/per_run/066_fertilization_a2c_nonadaptive_fixed_weather_years_5000_seed_1__training_reward_vs_global_step.png}
    \caption{066 Fertilization A2C Nonadaptive Fixed Weather Years 5000 Seed 1 - Training Reward Vs Global Step}
\end{figure}

\clearpage
\subsection{Run: 067 Fertilization A2C Adaptive Random Weather Years 5000 Seed 1}
\begin{figure}[H]
    \centering
    \includegraphics[width=0.8\textwidth]{figures/per_run/067_fertilization_a2c_adaptive_random_weather_years_5000_seed_1__checkpoint_eval_curves.png}
    \caption{067 Fertilization A2C Adaptive Random Weather Years 5000 Seed 1 - Checkpoint Eval Curves}
\end{figure}

\begin{figure}[H]
    \centering
    \includegraphics[width=0.8\textwidth]{figures/per_run/067_fertilization_a2c_adaptive_random_weather_years_5000_seed_1__diagnostics_panel.png}
    \caption{067 Fertilization A2C Adaptive Random Weather Years 5000 Seed 1 - Diagnostics Panel}
\end{figure}

\begin{figure}[H]
    \centering
    \includegraphics[width=0.8\textwidth]{figures/per_run/067_fertilization_a2c_adaptive_random_weather_years_5000_seed_1__episode_length_vs_global_step.png}
    \caption{067 Fertilization A2C Adaptive Random Weather Years 5000 Seed 1 - Episode Length Vs Global Step}
\end{figure}

\begin{figure}[H]
    \centering
    \includegraphics[width=0.8\textwidth]{figures/per_run/067_fertilization_a2c_adaptive_random_weather_years_5000_seed_1__primary_metric_vs_global_step.png}
    \caption{067 Fertilization A2C Adaptive Random Weather Years 5000 Seed 1 - Primary Metric Vs Global Step}
\end{figure}

\begin{figure}[H]
    \centering
    \includegraphics[width=0.8\textwidth]{figures/per_run/067_fertilization_a2c_adaptive_random_weather_years_5000_seed_1__training_reward_vs_global_step.png}
    \caption{067 Fertilization A2C Adaptive Random Weather Years 5000 Seed 1 - Training Reward Vs Global Step}
\end{figure}

\clearpage
\subsection{Run: 068 Fertilization A2C Nonadaptive Random Weather Years 5000 Seed 1}
\begin{figure}[H]
    \centering
    \includegraphics[width=0.8\textwidth]{figures/per_run/068_fertilization_a2c_nonadaptive_random_weather_years_5000_seed_1__checkpoint_eval_curves.png}
    \caption{068 Fertilization A2C Nonadaptive Random Weather Years 5000 Seed 1 - Checkpoint Eval Curves}
\end{figure}

\begin{figure}[H]
    \centering
    \includegraphics[width=0.8\textwidth]{figures/per_run/068_fertilization_a2c_nonadaptive_random_weather_years_5000_seed_1__diagnostics_panel.png}
    \caption{068 Fertilization A2C Nonadaptive Random Weather Years 5000 Seed 1 - Diagnostics Panel}
\end{figure}

\begin{figure}[H]
    \centering
    \includegraphics[width=0.8\textwidth]{figures/per_run/068_fertilization_a2c_nonadaptive_random_weather_years_5000_seed_1__episode_length_vs_global_step.png}
    \caption{068 Fertilization A2C Nonadaptive Random Weather Years 5000 Seed 1 - Episode Length Vs Global Step}
\end{figure}

\begin{figure}[H]
    \centering
    \includegraphics[width=0.8\textwidth]{figures/per_run/068_fertilization_a2c_nonadaptive_random_weather_years_5000_seed_1__primary_metric_vs_global_step.png}
    \caption{068 Fertilization A2C Nonadaptive Random Weather Years 5000 Seed 1 - Primary Metric Vs Global Step}
\end{figure}

\begin{figure}[H]
    \centering
    \includegraphics[width=0.8\textwidth]{figures/per_run/068_fertilization_a2c_nonadaptive_random_weather_years_5000_seed_1__training_reward_vs_global_step.png}
    \caption{068 Fertilization A2C Nonadaptive Random Weather Years 5000 Seed 1 - Training Reward Vs Global Step}
\end{figure}

\clearpage
\subsection{Run: 069 Fertilization A2C Adaptive Fixed Weather Years 5000 Seed 2}
\begin{figure}[H]
    \centering
    \includegraphics[width=0.8\textwidth]{figures/per_run/069_fertilization_a2c_adaptive_fixed_weather_years_5000_seed_2__checkpoint_eval_curves.png}
    \caption{069 Fertilization A2C Adaptive Fixed Weather Years 5000 Seed 2 - Checkpoint Eval Curves}
\end{figure}

\begin{figure}[H]
    \centering
    \includegraphics[width=0.8\textwidth]{figures/per_run/069_fertilization_a2c_adaptive_fixed_weather_years_5000_seed_2__diagnostics_panel.png}
    \caption{069 Fertilization A2C Adaptive Fixed Weather Years 5000 Seed 2 - Diagnostics Panel}
\end{figure}

\begin{figure}[H]
    \centering
    \includegraphics[width=0.8\textwidth]{figures/per_run/069_fertilization_a2c_adaptive_fixed_weather_years_5000_seed_2__episode_length_vs_global_step.png}
    \caption{069 Fertilization A2C Adaptive Fixed Weather Years 5000 Seed 2 - Episode Length Vs Global Step}
\end{figure}

\begin{figure}[H]
    \centering
    \includegraphics[width=0.8\textwidth]{figures/per_run/069_fertilization_a2c_adaptive_fixed_weather_years_5000_seed_2__primary_metric_vs_global_step.png}
    \caption{069 Fertilization A2C Adaptive Fixed Weather Years 5000 Seed 2 - Primary Metric Vs Global Step}
\end{figure}

\begin{figure}[H]
    \centering
    \includegraphics[width=0.8\textwidth]{figures/per_run/069_fertilization_a2c_adaptive_fixed_weather_years_5000_seed_2__training_reward_vs_global_step.png}
    \caption{069 Fertilization A2C Adaptive Fixed Weather Years 5000 Seed 2 - Training Reward Vs Global Step}
\end{figure}

\clearpage
\subsection{Run: 070 Fertilization A2C Nonadaptive Fixed Weather Years 5000 Seed 2}
\begin{figure}[H]
    \centering
    \includegraphics[width=0.8\textwidth]{figures/per_run/070_fertilization_a2c_nonadaptive_fixed_weather_years_5000_seed_2__checkpoint_eval_curves.png}
    \caption{070 Fertilization A2C Nonadaptive Fixed Weather Years 5000 Seed 2 - Checkpoint Eval Curves}
\end{figure}

\begin{figure}[H]
    \centering
    \includegraphics[width=0.8\textwidth]{figures/per_run/070_fertilization_a2c_nonadaptive_fixed_weather_years_5000_seed_2__diagnostics_panel.png}
    \caption{070 Fertilization A2C Nonadaptive Fixed Weather Years 5000 Seed 2 - Diagnostics Panel}
\end{figure}

\begin{figure}[H]
    \centering
    \includegraphics[width=0.8\textwidth]{figures/per_run/070_fertilization_a2c_nonadaptive_fixed_weather_years_5000_seed_2__episode_length_vs_global_step.png}
    \caption{070 Fertilization A2C Nonadaptive Fixed Weather Years 5000 Seed 2 - Episode Length Vs Global Step}
\end{figure}

\begin{figure}[H]
    \centering
    \includegraphics[width=0.8\textwidth]{figures/per_run/070_fertilization_a2c_nonadaptive_fixed_weather_years_5000_seed_2__primary_metric_vs_global_step.png}
    \caption{070 Fertilization A2C Nonadaptive Fixed Weather Years 5000 Seed 2 - Primary Metric Vs Global Step}
\end{figure}

\begin{figure}[H]
    \centering
    \includegraphics[width=0.8\textwidth]{figures/per_run/070_fertilization_a2c_nonadaptive_fixed_weather_years_5000_seed_2__training_reward_vs_global_step.png}
    \caption{070 Fertilization A2C Nonadaptive Fixed Weather Years 5000 Seed 2 - Training Reward Vs Global Step}
\end{figure}

\clearpage
\subsection{Run: 071 Fertilization A2C Adaptive Random Weather Years 5000 Seed 2}
\begin{figure}[H]
    \centering
    \includegraphics[width=0.8\textwidth]{figures/per_run/071_fertilization_a2c_adaptive_random_weather_years_5000_seed_2__checkpoint_eval_curves.png}
    \caption{071 Fertilization A2C Adaptive Random Weather Years 5000 Seed 2 - Checkpoint Eval Curves}
\end{figure}

\begin{figure}[H]
    \centering
    \includegraphics[width=0.8\textwidth]{figures/per_run/071_fertilization_a2c_adaptive_random_weather_years_5000_seed_2__diagnostics_panel.png}
    \caption{071 Fertilization A2C Adaptive Random Weather Years 5000 Seed 2 - Diagnostics Panel}
\end{figure}

\begin{figure}[H]
    \centering
    \includegraphics[width=0.8\textwidth]{figures/per_run/071_fertilization_a2c_adaptive_random_weather_years_5000_seed_2__episode_length_vs_global_step.png}
    \caption{071 Fertilization A2C Adaptive Random Weather Years 5000 Seed 2 - Episode Length Vs Global Step}
\end{figure}

\begin{figure}[H]
    \centering
    \includegraphics[width=0.8\textwidth]{figures/per_run/071_fertilization_a2c_adaptive_random_weather_years_5000_seed_2__primary_metric_vs_global_step.png}
    \caption{071 Fertilization A2C Adaptive Random Weather Years 5000 Seed 2 - Primary Metric Vs Global Step}
\end{figure}

\begin{figure}[H]
    \centering
    \includegraphics[width=0.8\textwidth]{figures/per_run/071_fertilization_a2c_adaptive_random_weather_years_5000_seed_2__training_reward_vs_global_step.png}
    \caption{071 Fertilization A2C Adaptive Random Weather Years 5000 Seed 2 - Training Reward Vs Global Step}
\end{figure}

\clearpage
\subsection{Run: 072 Fertilization A2C Nonadaptive Random Weather Years 5000 Seed 2}
\begin{figure}[H]
    \centering
    \includegraphics[width=0.8\textwidth]{figures/per_run/072_fertilization_a2c_nonadaptive_random_weather_years_5000_seed_2__checkpoint_eval_curves.png}
    \caption{072 Fertilization A2C Nonadaptive Random Weather Years 5000 Seed 2 - Checkpoint Eval Curves}
\end{figure}

\begin{figure}[H]
    \centering
    \includegraphics[width=0.8\textwidth]{figures/per_run/072_fertilization_a2c_nonadaptive_random_weather_years_5000_seed_2__diagnostics_panel.png}
    \caption{072 Fertilization A2C Nonadaptive Random Weather Years 5000 Seed 2 - Diagnostics Panel}
\end{figure}

\begin{figure}[H]
    \centering
    \includegraphics[width=0.8\textwidth]{figures/per_run/072_fertilization_a2c_nonadaptive_random_weather_years_5000_seed_2__episode_length_vs_global_step.png}
    \caption{072 Fertilization A2C Nonadaptive Random Weather Years 5000 Seed 2 - Episode Length Vs Global Step}
\end{figure}

\begin{figure}[H]
    \centering
    \includegraphics[width=0.8\textwidth]{figures/per_run/072_fertilization_a2c_nonadaptive_random_weather_years_5000_seed_2__primary_metric_vs_global_step.png}
    \caption{072 Fertilization A2C Nonadaptive Random Weather Years 5000 Seed 2 - Primary Metric Vs Global Step}
\end{figure}

\begin{figure}[H]
    \centering
    \includegraphics[width=0.8\textwidth]{figures/per_run/072_fertilization_a2c_nonadaptive_random_weather_years_5000_seed_2__training_reward_vs_global_step.png}
    \caption{072 Fertilization A2C Nonadaptive Random Weather Years 5000 Seed 2 - Training Reward Vs Global Step}
\end{figure}

\clearpage
\subsection{Run: 073 Crop Planning Ppo Adaptive Fixed Weather Seed 0}
\begin{figure}[H]
    \centering
    \includegraphics[width=0.8\textwidth]{figures/per_run/073_crop_planning_ppo_adaptive_fixed_weather_seed_0__checkpoint_eval_curves.png}
    \caption{073 Crop Planning Ppo Adaptive Fixed Weather Seed 0 - Checkpoint Eval Curves}
\end{figure}

\begin{figure}[H]
    \centering
    \includegraphics[width=0.8\textwidth]{figures/per_run/073_crop_planning_ppo_adaptive_fixed_weather_seed_0__diagnostics_panel.png}
    \caption{073 Crop Planning Ppo Adaptive Fixed Weather Seed 0 - Diagnostics Panel}
\end{figure}

\begin{figure}[H]
    \centering
    \includegraphics[width=0.8\textwidth]{figures/per_run/073_crop_planning_ppo_adaptive_fixed_weather_seed_0__episode_length_vs_global_step.png}
    \caption{073 Crop Planning Ppo Adaptive Fixed Weather Seed 0 - Episode Length Vs Global Step}
\end{figure}

\begin{figure}[H]
    \centering
    \includegraphics[width=0.8\textwidth]{figures/per_run/073_crop_planning_ppo_adaptive_fixed_weather_seed_0__primary_metric_vs_global_step.png}
    \caption{073 Crop Planning Ppo Adaptive Fixed Weather Seed 0 - Primary Metric Vs Global Step}
\end{figure}

\begin{figure}[H]
    \centering
    \includegraphics[width=0.8\textwidth]{figures/per_run/073_crop_planning_ppo_adaptive_fixed_weather_seed_0__training_reward_vs_global_step.png}
    \caption{073 Crop Planning Ppo Adaptive Fixed Weather Seed 0 - Training Reward Vs Global Step}
\end{figure}

\clearpage
\subsection{Run: 074 Crop Planning Ppo Nonadaptive Fixed Weather Seed 0}
\begin{figure}[H]
    \centering
    \includegraphics[width=0.8\textwidth]{figures/per_run/074_crop_planning_ppo_nonadaptive_fixed_weather_seed_0__checkpoint_eval_curves.png}
    \caption{074 Crop Planning Ppo Nonadaptive Fixed Weather Seed 0 - Checkpoint Eval Curves}
\end{figure}

\begin{figure}[H]
    \centering
    \includegraphics[width=0.8\textwidth]{figures/per_run/074_crop_planning_ppo_nonadaptive_fixed_weather_seed_0__diagnostics_panel.png}
    \caption{074 Crop Planning Ppo Nonadaptive Fixed Weather Seed 0 - Diagnostics Panel}
\end{figure}

\begin{figure}[H]
    \centering
    \includegraphics[width=0.8\textwidth]{figures/per_run/074_crop_planning_ppo_nonadaptive_fixed_weather_seed_0__episode_length_vs_global_step.png}
    \caption{074 Crop Planning Ppo Nonadaptive Fixed Weather Seed 0 - Episode Length Vs Global Step}
\end{figure}

\begin{figure}[H]
    \centering
    \includegraphics[width=0.8\textwidth]{figures/per_run/074_crop_planning_ppo_nonadaptive_fixed_weather_seed_0__primary_metric_vs_global_step.png}
    \caption{074 Crop Planning Ppo Nonadaptive Fixed Weather Seed 0 - Primary Metric Vs Global Step}
\end{figure}

\begin{figure}[H]
    \centering
    \includegraphics[width=0.8\textwidth]{figures/per_run/074_crop_planning_ppo_nonadaptive_fixed_weather_seed_0__training_reward_vs_global_step.png}
    \caption{074 Crop Planning Ppo Nonadaptive Fixed Weather Seed 0 - Training Reward Vs Global Step}
\end{figure}

\clearpage
\subsection{Run: 075 Crop Planning Hierarchical Ppo Fixed Weather Seed 0 Profile Pakistan Baseline}
\begin{figure}[H]
    \centering
    \includegraphics[width=0.8\textwidth]{figures/per_run/075_crop_planning_hierarchical_ppo_fixed_weather_seed_0_profile_pakistan_baseline__diagnostics_panel.png}
    \caption{075 Crop Planning Hierarchical Ppo Fixed Weather Seed 0 Profile Pakistan Baseline - Diagnostics Panel}
\end{figure}

\begin{figure}[H]
    \centering
    \includegraphics[width=0.8\textwidth]{figures/per_run/075_crop_planning_hierarchical_ppo_fixed_weather_seed_0_profile_pakistan_baseline__primary_metric_vs_global_step.png}
    \caption{075 Crop Planning Hierarchical Ppo Fixed Weather Seed 0 Profile Pakistan Baseline - Primary Metric Vs Global Step}
\end{figure}

\clearpage
\subsection{Run: 076 Crop Planning Ppo Adaptive Random Weather Seed 0}
\begin{figure}[H]
    \centering
    \includegraphics[width=0.8\textwidth]{figures/per_run/076_crop_planning_ppo_adaptive_random_weather_seed_0__checkpoint_eval_curves.png}
    \caption{076 Crop Planning Ppo Adaptive Random Weather Seed 0 - Checkpoint Eval Curves}
\end{figure}

\begin{figure}[H]
    \centering
    \includegraphics[width=0.8\textwidth]{figures/per_run/076_crop_planning_ppo_adaptive_random_weather_seed_0__diagnostics_panel.png}
    \caption{076 Crop Planning Ppo Adaptive Random Weather Seed 0 - Diagnostics Panel}
\end{figure}

\begin{figure}[H]
    \centering
    \includegraphics[width=0.8\textwidth]{figures/per_run/076_crop_planning_ppo_adaptive_random_weather_seed_0__episode_length_vs_global_step.png}
    \caption{076 Crop Planning Ppo Adaptive Random Weather Seed 0 - Episode Length Vs Global Step}
\end{figure}

\begin{figure}[H]
    \centering
    \includegraphics[width=0.8\textwidth]{figures/per_run/076_crop_planning_ppo_adaptive_random_weather_seed_0__primary_metric_vs_global_step.png}
    \caption{076 Crop Planning Ppo Adaptive Random Weather Seed 0 - Primary Metric Vs Global Step}
\end{figure}

\begin{figure}[H]
    \centering
    \includegraphics[width=0.8\textwidth]{figures/per_run/076_crop_planning_ppo_adaptive_random_weather_seed_0__training_reward_vs_global_step.png}
    \caption{076 Crop Planning Ppo Adaptive Random Weather Seed 0 - Training Reward Vs Global Step}
\end{figure}

\clearpage
\subsection{Run: 077 Crop Planning Ppo Nonadaptive Random Weather Seed 0}
\begin{figure}[H]
    \centering
    \includegraphics[width=0.8\textwidth]{figures/per_run/077_crop_planning_ppo_nonadaptive_random_weather_seed_0__checkpoint_eval_curves.png}
    \caption{077 Crop Planning Ppo Nonadaptive Random Weather Seed 0 - Checkpoint Eval Curves}
\end{figure}

\begin{figure}[H]
    \centering
    \includegraphics[width=0.8\textwidth]{figures/per_run/077_crop_planning_ppo_nonadaptive_random_weather_seed_0__diagnostics_panel.png}
    \caption{077 Crop Planning Ppo Nonadaptive Random Weather Seed 0 - Diagnostics Panel}
\end{figure}

\begin{figure}[H]
    \centering
    \includegraphics[width=0.8\textwidth]{figures/per_run/077_crop_planning_ppo_nonadaptive_random_weather_seed_0__episode_length_vs_global_step.png}
    \caption{077 Crop Planning Ppo Nonadaptive Random Weather Seed 0 - Episode Length Vs Global Step}
\end{figure}

\begin{figure}[H]
    \centering
    \includegraphics[width=0.8\textwidth]{figures/per_run/077_crop_planning_ppo_nonadaptive_random_weather_seed_0__primary_metric_vs_global_step.png}
    \caption{077 Crop Planning Ppo Nonadaptive Random Weather Seed 0 - Primary Metric Vs Global Step}
\end{figure}

\begin{figure}[H]
    \centering
    \includegraphics[width=0.8\textwidth]{figures/per_run/077_crop_planning_ppo_nonadaptive_random_weather_seed_0__training_reward_vs_global_step.png}
    \caption{077 Crop Planning Ppo Nonadaptive Random Weather Seed 0 - Training Reward Vs Global Step}
\end{figure}

\clearpage
\subsection{Run: 078 Crop Planning Hierarchical Ppo Random Weather Seed 0 Profile Pakistan Baseline}
\begin{figure}[H]
    \centering
    \includegraphics[width=0.8\textwidth]{figures/per_run/078_crop_planning_hierarchical_ppo_random_weather_seed_0_profile_pakistan_baseline__diagnostics_panel.png}
    \caption{078 Crop Planning Hierarchical Ppo Random Weather Seed 0 Profile Pakistan Baseline - Diagnostics Panel}
\end{figure}

\begin{figure}[H]
    \centering
    \includegraphics[width=0.8\textwidth]{figures/per_run/078_crop_planning_hierarchical_ppo_random_weather_seed_0_profile_pakistan_baseline__primary_metric_vs_global_step.png}
    \caption{078 Crop Planning Hierarchical Ppo Random Weather Seed 0 Profile Pakistan Baseline - Primary Metric Vs Global Step}
\end{figure}

\clearpage
\subsection{Run: 079 Crop Planning Ppo Adaptive Fixed Weather Seed 1}
\begin{figure}[H]
    \centering
    \includegraphics[width=0.8\textwidth]{figures/per_run/079_crop_planning_ppo_adaptive_fixed_weather_seed_1__checkpoint_eval_curves.png}
    \caption{079 Crop Planning Ppo Adaptive Fixed Weather Seed 1 - Checkpoint Eval Curves}
\end{figure}

\begin{figure}[H]
    \centering
    \includegraphics[width=0.8\textwidth]{figures/per_run/079_crop_planning_ppo_adaptive_fixed_weather_seed_1__diagnostics_panel.png}
    \caption{079 Crop Planning Ppo Adaptive Fixed Weather Seed 1 - Diagnostics Panel}
\end{figure}

\begin{figure}[H]
    \centering
    \includegraphics[width=0.8\textwidth]{figures/per_run/079_crop_planning_ppo_adaptive_fixed_weather_seed_1__episode_length_vs_global_step.png}
    \caption{079 Crop Planning Ppo Adaptive Fixed Weather Seed 1 - Episode Length Vs Global Step}
\end{figure}

\begin{figure}[H]
    \centering
    \includegraphics[width=0.8\textwidth]{figures/per_run/079_crop_planning_ppo_adaptive_fixed_weather_seed_1__primary_metric_vs_global_step.png}
    \caption{079 Crop Planning Ppo Adaptive Fixed Weather Seed 1 - Primary Metric Vs Global Step}
\end{figure}

\begin{figure}[H]
    \centering
    \includegraphics[width=0.8\textwidth]{figures/per_run/079_crop_planning_ppo_adaptive_fixed_weather_seed_1__training_reward_vs_global_step.png}
    \caption{079 Crop Planning Ppo Adaptive Fixed Weather Seed 1 - Training Reward Vs Global Step}
\end{figure}

\clearpage
\subsection{Run: 080 Crop Planning Ppo Nonadaptive Fixed Weather Seed 1}
\begin{figure}[H]
    \centering
    \includegraphics[width=0.8\textwidth]{figures/per_run/080_crop_planning_ppo_nonadaptive_fixed_weather_seed_1__checkpoint_eval_curves.png}
    \caption{080 Crop Planning Ppo Nonadaptive Fixed Weather Seed 1 - Checkpoint Eval Curves}
\end{figure}

\begin{figure}[H]
    \centering
    \includegraphics[width=0.8\textwidth]{figures/per_run/080_crop_planning_ppo_nonadaptive_fixed_weather_seed_1__diagnostics_panel.png}
    \caption{080 Crop Planning Ppo Nonadaptive Fixed Weather Seed 1 - Diagnostics Panel}
\end{figure}

\begin{figure}[H]
    \centering
    \includegraphics[width=0.8\textwidth]{figures/per_run/080_crop_planning_ppo_nonadaptive_fixed_weather_seed_1__episode_length_vs_global_step.png}
    \caption{080 Crop Planning Ppo Nonadaptive Fixed Weather Seed 1 - Episode Length Vs Global Step}
\end{figure}

\begin{figure}[H]
    \centering
    \includegraphics[width=0.8\textwidth]{figures/per_run/080_crop_planning_ppo_nonadaptive_fixed_weather_seed_1__primary_metric_vs_global_step.png}
    \caption{080 Crop Planning Ppo Nonadaptive Fixed Weather Seed 1 - Primary Metric Vs Global Step}
\end{figure}

\begin{figure}[H]
    \centering
    \includegraphics[width=0.8\textwidth]{figures/per_run/080_crop_planning_ppo_nonadaptive_fixed_weather_seed_1__training_reward_vs_global_step.png}
    \caption{080 Crop Planning Ppo Nonadaptive Fixed Weather Seed 1 - Training Reward Vs Global Step}
\end{figure}

\clearpage
\subsection{Run: 081 Crop Planning Hierarchical Ppo Fixed Weather Seed 1 Profile Pakistan Baseline}
\begin{figure}[H]
    \centering
    \includegraphics[width=0.8\textwidth]{figures/per_run/081_crop_planning_hierarchical_ppo_fixed_weather_seed_1_profile_pakistan_baseline__diagnostics_panel.png}
    \caption{081 Crop Planning Hierarchical Ppo Fixed Weather Seed 1 Profile Pakistan Baseline - Diagnostics Panel}
\end{figure}

\begin{figure}[H]
    \centering
    \includegraphics[width=0.8\textwidth]{figures/per_run/081_crop_planning_hierarchical_ppo_fixed_weather_seed_1_profile_pakistan_baseline__primary_metric_vs_global_step.png}
    \caption{081 Crop Planning Hierarchical Ppo Fixed Weather Seed 1 Profile Pakistan Baseline - Primary Metric Vs Global Step}
\end{figure}

\clearpage
\subsection{Run: 082 Crop Planning Ppo Adaptive Random Weather Seed 1}
\begin{figure}[H]
    \centering
    \includegraphics[width=0.8\textwidth]{figures/per_run/082_crop_planning_ppo_adaptive_random_weather_seed_1__checkpoint_eval_curves.png}
    \caption{082 Crop Planning Ppo Adaptive Random Weather Seed 1 - Checkpoint Eval Curves}
\end{figure}

\begin{figure}[H]
    \centering
    \includegraphics[width=0.8\textwidth]{figures/per_run/082_crop_planning_ppo_adaptive_random_weather_seed_1__diagnostics_panel.png}
    \caption{082 Crop Planning Ppo Adaptive Random Weather Seed 1 - Diagnostics Panel}
\end{figure}

\begin{figure}[H]
    \centering
    \includegraphics[width=0.8\textwidth]{figures/per_run/082_crop_planning_ppo_adaptive_random_weather_seed_1__episode_length_vs_global_step.png}
    \caption{082 Crop Planning Ppo Adaptive Random Weather Seed 1 - Episode Length Vs Global Step}
\end{figure}

\begin{figure}[H]
    \centering
    \includegraphics[width=0.8\textwidth]{figures/per_run/082_crop_planning_ppo_adaptive_random_weather_seed_1__primary_metric_vs_global_step.png}
    \caption{082 Crop Planning Ppo Adaptive Random Weather Seed 1 - Primary Metric Vs Global Step}
\end{figure}

\begin{figure}[H]
    \centering
    \includegraphics[width=0.8\textwidth]{figures/per_run/082_crop_planning_ppo_adaptive_random_weather_seed_1__training_reward_vs_global_step.png}
    \caption{082 Crop Planning Ppo Adaptive Random Weather Seed 1 - Training Reward Vs Global Step}
\end{figure}

\clearpage
\subsection{Run: 083 Crop Planning Ppo Nonadaptive Random Weather Seed 1}
\begin{figure}[H]
    \centering
    \includegraphics[width=0.8\textwidth]{figures/per_run/083_crop_planning_ppo_nonadaptive_random_weather_seed_1__checkpoint_eval_curves.png}
    \caption{083 Crop Planning Ppo Nonadaptive Random Weather Seed 1 - Checkpoint Eval Curves}
\end{figure}

\begin{figure}[H]
    \centering
    \includegraphics[width=0.8\textwidth]{figures/per_run/083_crop_planning_ppo_nonadaptive_random_weather_seed_1__diagnostics_panel.png}
    \caption{083 Crop Planning Ppo Nonadaptive Random Weather Seed 1 - Diagnostics Panel}
\end{figure}

\begin{figure}[H]
    \centering
    \includegraphics[width=0.8\textwidth]{figures/per_run/083_crop_planning_ppo_nonadaptive_random_weather_seed_1__episode_length_vs_global_step.png}
    \caption{083 Crop Planning Ppo Nonadaptive Random Weather Seed 1 - Episode Length Vs Global Step}
\end{figure}

\begin{figure}[H]
    \centering
    \includegraphics[width=0.8\textwidth]{figures/per_run/083_crop_planning_ppo_nonadaptive_random_weather_seed_1__primary_metric_vs_global_step.png}
    \caption{083 Crop Planning Ppo Nonadaptive Random Weather Seed 1 - Primary Metric Vs Global Step}
\end{figure}

\begin{figure}[H]
    \centering
    \includegraphics[width=0.8\textwidth]{figures/per_run/083_crop_planning_ppo_nonadaptive_random_weather_seed_1__training_reward_vs_global_step.png}
    \caption{083 Crop Planning Ppo Nonadaptive Random Weather Seed 1 - Training Reward Vs Global Step}
\end{figure}

\clearpage
\subsection{Run: 084 Crop Planning Hierarchical Ppo Random Weather Seed 1 Profile Pakistan Baseline}
\begin{figure}[H]
    \centering
    \includegraphics[width=0.8\textwidth]{figures/per_run/084_crop_planning_hierarchical_ppo_random_weather_seed_1_profile_pakistan_baseline__diagnostics_panel.png}
    \caption{084 Crop Planning Hierarchical Ppo Random Weather Seed 1 Profile Pakistan Baseline - Diagnostics Panel}
\end{figure}

\begin{figure}[H]
    \centering
    \includegraphics[width=0.8\textwidth]{figures/per_run/084_crop_planning_hierarchical_ppo_random_weather_seed_1_profile_pakistan_baseline__primary_metric_vs_global_step.png}
    \caption{084 Crop Planning Hierarchical Ppo Random Weather Seed 1 Profile Pakistan Baseline - Primary Metric Vs Global Step}
\end{figure}

\clearpage
\subsection{Run: 085 Crop Planning Ppo Adaptive Fixed Weather Seed 2}
\begin{figure}[H]
    \centering
    \includegraphics[width=0.8\textwidth]{figures/per_run/085_crop_planning_ppo_adaptive_fixed_weather_seed_2__checkpoint_eval_curves.png}
    \caption{085 Crop Planning Ppo Adaptive Fixed Weather Seed 2 - Checkpoint Eval Curves}
\end{figure}

\begin{figure}[H]
    \centering
    \includegraphics[width=0.8\textwidth]{figures/per_run/085_crop_planning_ppo_adaptive_fixed_weather_seed_2__diagnostics_panel.png}
    \caption{085 Crop Planning Ppo Adaptive Fixed Weather Seed 2 - Diagnostics Panel}
\end{figure}

\begin{figure}[H]
    \centering
    \includegraphics[width=0.8\textwidth]{figures/per_run/085_crop_planning_ppo_adaptive_fixed_weather_seed_2__episode_length_vs_global_step.png}
    \caption{085 Crop Planning Ppo Adaptive Fixed Weather Seed 2 - Episode Length Vs Global Step}
\end{figure}

\begin{figure}[H]
    \centering
    \includegraphics[width=0.8\textwidth]{figures/per_run/085_crop_planning_ppo_adaptive_fixed_weather_seed_2__primary_metric_vs_global_step.png}
    \caption{085 Crop Planning Ppo Adaptive Fixed Weather Seed 2 - Primary Metric Vs Global Step}
\end{figure}

\begin{figure}[H]
    \centering
    \includegraphics[width=0.8\textwidth]{figures/per_run/085_crop_planning_ppo_adaptive_fixed_weather_seed_2__training_reward_vs_global_step.png}
    \caption{085 Crop Planning Ppo Adaptive Fixed Weather Seed 2 - Training Reward Vs Global Step}
\end{figure}

\clearpage
\subsection{Run: 086 Crop Planning Ppo Nonadaptive Fixed Weather Seed 2}
\begin{figure}[H]
    \centering
    \includegraphics[width=0.8\textwidth]{figures/per_run/086_crop_planning_ppo_nonadaptive_fixed_weather_seed_2__checkpoint_eval_curves.png}
    \caption{086 Crop Planning Ppo Nonadaptive Fixed Weather Seed 2 - Checkpoint Eval Curves}
\end{figure}

\begin{figure}[H]
    \centering
    \includegraphics[width=0.8\textwidth]{figures/per_run/086_crop_planning_ppo_nonadaptive_fixed_weather_seed_2__diagnostics_panel.png}
    \caption{086 Crop Planning Ppo Nonadaptive Fixed Weather Seed 2 - Diagnostics Panel}
\end{figure}

\begin{figure}[H]
    \centering
    \includegraphics[width=0.8\textwidth]{figures/per_run/086_crop_planning_ppo_nonadaptive_fixed_weather_seed_2__episode_length_vs_global_step.png}
    \caption{086 Crop Planning Ppo Nonadaptive Fixed Weather Seed 2 - Episode Length Vs Global Step}
\end{figure}

\begin{figure}[H]
    \centering
    \includegraphics[width=0.8\textwidth]{figures/per_run/086_crop_planning_ppo_nonadaptive_fixed_weather_seed_2__primary_metric_vs_global_step.png}
    \caption{086 Crop Planning Ppo Nonadaptive Fixed Weather Seed 2 - Primary Metric Vs Global Step}
\end{figure}

\begin{figure}[H]
    \centering
    \includegraphics[width=0.8\textwidth]{figures/per_run/086_crop_planning_ppo_nonadaptive_fixed_weather_seed_2__training_reward_vs_global_step.png}
    \caption{086 Crop Planning Ppo Nonadaptive Fixed Weather Seed 2 - Training Reward Vs Global Step}
\end{figure}

\clearpage
\subsection{Run: 087 Crop Planning Hierarchical Ppo Fixed Weather Seed 2 Profile Pakistan Baseline}
\begin{figure}[H]
    \centering
    \includegraphics[width=0.8\textwidth]{figures/per_run/087_crop_planning_hierarchical_ppo_fixed_weather_seed_2_profile_pakistan_baseline__diagnostics_panel.png}
    \caption{087 Crop Planning Hierarchical Ppo Fixed Weather Seed 2 Profile Pakistan Baseline - Diagnostics Panel}
\end{figure}

\begin{figure}[H]
    \centering
    \includegraphics[width=0.8\textwidth]{figures/per_run/087_crop_planning_hierarchical_ppo_fixed_weather_seed_2_profile_pakistan_baseline__primary_metric_vs_global_step.png}
    \caption{087 Crop Planning Hierarchical Ppo Fixed Weather Seed 2 Profile Pakistan Baseline - Primary Metric Vs Global Step}
\end{figure}

\clearpage
\subsection{Run: 088 Crop Planning Ppo Adaptive Random Weather Seed 2}
\begin{figure}[H]
    \centering
    \includegraphics[width=0.8\textwidth]{figures/per_run/088_crop_planning_ppo_adaptive_random_weather_seed_2__checkpoint_eval_curves.png}
    \caption{088 Crop Planning Ppo Adaptive Random Weather Seed 2 - Checkpoint Eval Curves}
\end{figure}

\begin{figure}[H]
    \centering
    \includegraphics[width=0.8\textwidth]{figures/per_run/088_crop_planning_ppo_adaptive_random_weather_seed_2__diagnostics_panel.png}
    \caption{088 Crop Planning Ppo Adaptive Random Weather Seed 2 - Diagnostics Panel}
\end{figure}

\begin{figure}[H]
    \centering
    \includegraphics[width=0.8\textwidth]{figures/per_run/088_crop_planning_ppo_adaptive_random_weather_seed_2__episode_length_vs_global_step.png}
    \caption{088 Crop Planning Ppo Adaptive Random Weather Seed 2 - Episode Length Vs Global Step}
\end{figure}

\begin{figure}[H]
    \centering
    \includegraphics[width=0.8\textwidth]{figures/per_run/088_crop_planning_ppo_adaptive_random_weather_seed_2__primary_metric_vs_global_step.png}
    \caption{088 Crop Planning Ppo Adaptive Random Weather Seed 2 - Primary Metric Vs Global Step}
\end{figure}

\begin{figure}[H]
    \centering
    \includegraphics[width=0.8\textwidth]{figures/per_run/088_crop_planning_ppo_adaptive_random_weather_seed_2__training_reward_vs_global_step.png}
    \caption{088 Crop Planning Ppo Adaptive Random Weather Seed 2 - Training Reward Vs Global Step}
\end{figure}

\clearpage
\subsection{Run: 089 Crop Planning Ppo Nonadaptive Random Weather Seed 2}
\begin{figure}[H]
    \centering
    \includegraphics[width=0.8\textwidth]{figures/per_run/089_crop_planning_ppo_nonadaptive_random_weather_seed_2__checkpoint_eval_curves.png}
    \caption{089 Crop Planning Ppo Nonadaptive Random Weather Seed 2 - Checkpoint Eval Curves}
\end{figure}

\begin{figure}[H]
    \centering
    \includegraphics[width=0.8\textwidth]{figures/per_run/089_crop_planning_ppo_nonadaptive_random_weather_seed_2__diagnostics_panel.png}
    \caption{089 Crop Planning Ppo Nonadaptive Random Weather Seed 2 - Diagnostics Panel}
\end{figure}

\begin{figure}[H]
    \centering
    \includegraphics[width=0.8\textwidth]{figures/per_run/089_crop_planning_ppo_nonadaptive_random_weather_seed_2__episode_length_vs_global_step.png}
    \caption{089 Crop Planning Ppo Nonadaptive Random Weather Seed 2 - Episode Length Vs Global Step}
\end{figure}

\begin{figure}[H]
    \centering
    \includegraphics[width=0.8\textwidth]{figures/per_run/089_crop_planning_ppo_nonadaptive_random_weather_seed_2__primary_metric_vs_global_step.png}
    \caption{089 Crop Planning Ppo Nonadaptive Random Weather Seed 2 - Primary Metric Vs Global Step}
\end{figure}

\begin{figure}[H]
    \centering
    \includegraphics[width=0.8\textwidth]{figures/per_run/089_crop_planning_ppo_nonadaptive_random_weather_seed_2__training_reward_vs_global_step.png}
    \caption{089 Crop Planning Ppo Nonadaptive Random Weather Seed 2 - Training Reward Vs Global Step}
\end{figure}

\clearpage
\subsection{Run: 090 Crop Planning Hierarchical Ppo Random Weather Seed 2 Profile Pakistan Baseline}
\begin{figure}[H]
    \centering
    \includegraphics[width=0.8\textwidth]{figures/per_run/090_crop_planning_hierarchical_ppo_random_weather_seed_2_profile_pakistan_baseline__diagnostics_panel.png}
    \caption{090 Crop Planning Hierarchical Ppo Random Weather Seed 2 Profile Pakistan Baseline - Diagnostics Panel}
\end{figure}

\begin{figure}[H]
    \centering
    \includegraphics[width=0.8\textwidth]{figures/per_run/090_crop_planning_hierarchical_ppo_random_weather_seed_2_profile_pakistan_baseline__primary_metric_vs_global_step.png}
    \caption{090 Crop Planning Hierarchical Ppo Random Weather Seed 2 Profile Pakistan Baseline - Primary Metric Vs Global Step}
\end{figure}

\clearpage
\subsection{Run: 091 Crop Planning A2C Adaptive Fixed Weather Seed 0}
\begin{figure}[H]
    \centering
    \includegraphics[width=0.8\textwidth]{figures/per_run/091_crop_planning_a2c_adaptive_fixed_weather_seed_0__checkpoint_eval_curves.png}
    \caption{091 Crop Planning A2C Adaptive Fixed Weather Seed 0 - Checkpoint Eval Curves}
\end{figure}

\begin{figure}[H]
    \centering
    \includegraphics[width=0.8\textwidth]{figures/per_run/091_crop_planning_a2c_adaptive_fixed_weather_seed_0__diagnostics_panel.png}
    \caption{091 Crop Planning A2C Adaptive Fixed Weather Seed 0 - Diagnostics Panel}
\end{figure}

\begin{figure}[H]
    \centering
    \includegraphics[width=0.8\textwidth]{figures/per_run/091_crop_planning_a2c_adaptive_fixed_weather_seed_0__episode_length_vs_global_step.png}
    \caption{091 Crop Planning A2C Adaptive Fixed Weather Seed 0 - Episode Length Vs Global Step}
\end{figure}

\begin{figure}[H]
    \centering
    \includegraphics[width=0.8\textwidth]{figures/per_run/091_crop_planning_a2c_adaptive_fixed_weather_seed_0__primary_metric_vs_global_step.png}
    \caption{091 Crop Planning A2C Adaptive Fixed Weather Seed 0 - Primary Metric Vs Global Step}
\end{figure}

\begin{figure}[H]
    \centering
    \includegraphics[width=0.8\textwidth]{figures/per_run/091_crop_planning_a2c_adaptive_fixed_weather_seed_0__training_reward_vs_global_step.png}
    \caption{091 Crop Planning A2C Adaptive Fixed Weather Seed 0 - Training Reward Vs Global Step}
\end{figure}

\clearpage
\subsection{Run: 092 Crop Planning A2C Nonadaptive Fixed Weather Seed 0}
\begin{figure}[H]
    \centering
    \includegraphics[width=0.8\textwidth]{figures/per_run/092_crop_planning_a2c_nonadaptive_fixed_weather_seed_0__checkpoint_eval_curves.png}
    \caption{092 Crop Planning A2C Nonadaptive Fixed Weather Seed 0 - Checkpoint Eval Curves}
\end{figure}

\begin{figure}[H]
    \centering
    \includegraphics[width=0.8\textwidth]{figures/per_run/092_crop_planning_a2c_nonadaptive_fixed_weather_seed_0__diagnostics_panel.png}
    \caption{092 Crop Planning A2C Nonadaptive Fixed Weather Seed 0 - Diagnostics Panel}
\end{figure}

\begin{figure}[H]
    \centering
    \includegraphics[width=0.8\textwidth]{figures/per_run/092_crop_planning_a2c_nonadaptive_fixed_weather_seed_0__episode_length_vs_global_step.png}
    \caption{092 Crop Planning A2C Nonadaptive Fixed Weather Seed 0 - Episode Length Vs Global Step}
\end{figure}

\begin{figure}[H]
    \centering
    \includegraphics[width=0.8\textwidth]{figures/per_run/092_crop_planning_a2c_nonadaptive_fixed_weather_seed_0__primary_metric_vs_global_step.png}
    \caption{092 Crop Planning A2C Nonadaptive Fixed Weather Seed 0 - Primary Metric Vs Global Step}
\end{figure}

\begin{figure}[H]
    \centering
    \includegraphics[width=0.8\textwidth]{figures/per_run/092_crop_planning_a2c_nonadaptive_fixed_weather_seed_0__training_reward_vs_global_step.png}
    \caption{092 Crop Planning A2C Nonadaptive Fixed Weather Seed 0 - Training Reward Vs Global Step}
\end{figure}

\clearpage
\subsection{Run: 093 Crop Planning Hierarchical A2C Fixed Weather Seed 0 Profile Pakistan Baseline}
\begin{figure}[H]
    \centering
    \includegraphics[width=0.8\textwidth]{figures/per_run/093_crop_planning_hierarchical_a2c_fixed_weather_seed_0_profile_pakistan_baseline__diagnostics_panel.png}
    \caption{093 Crop Planning Hierarchical A2C Fixed Weather Seed 0 Profile Pakistan Baseline - Diagnostics Panel}
\end{figure}

\begin{figure}[H]
    \centering
    \includegraphics[width=0.8\textwidth]{figures/per_run/093_crop_planning_hierarchical_a2c_fixed_weather_seed_0_profile_pakistan_baseline__primary_metric_vs_global_step.png}
    \caption{093 Crop Planning Hierarchical A2C Fixed Weather Seed 0 Profile Pakistan Baseline - Primary Metric Vs Global Step}
\end{figure}

\clearpage
\subsection{Run: 094 Crop Planning A2C Adaptive Random Weather Seed 0}
\begin{figure}[H]
    \centering
    \includegraphics[width=0.8\textwidth]{figures/per_run/094_crop_planning_a2c_adaptive_random_weather_seed_0__checkpoint_eval_curves.png}
    \caption{094 Crop Planning A2C Adaptive Random Weather Seed 0 - Checkpoint Eval Curves}
\end{figure}

\begin{figure}[H]
    \centering
    \includegraphics[width=0.8\textwidth]{figures/per_run/094_crop_planning_a2c_adaptive_random_weather_seed_0__diagnostics_panel.png}
    \caption{094 Crop Planning A2C Adaptive Random Weather Seed 0 - Diagnostics Panel}
\end{figure}

\begin{figure}[H]
    \centering
    \includegraphics[width=0.8\textwidth]{figures/per_run/094_crop_planning_a2c_adaptive_random_weather_seed_0__episode_length_vs_global_step.png}
    \caption{094 Crop Planning A2C Adaptive Random Weather Seed 0 - Episode Length Vs Global Step}
\end{figure}

\begin{figure}[H]
    \centering
    \includegraphics[width=0.8\textwidth]{figures/per_run/094_crop_planning_a2c_adaptive_random_weather_seed_0__primary_metric_vs_global_step.png}
    \caption{094 Crop Planning A2C Adaptive Random Weather Seed 0 - Primary Metric Vs Global Step}
\end{figure}

\begin{figure}[H]
    \centering
    \includegraphics[width=0.8\textwidth]{figures/per_run/094_crop_planning_a2c_adaptive_random_weather_seed_0__training_reward_vs_global_step.png}
    \caption{094 Crop Planning A2C Adaptive Random Weather Seed 0 - Training Reward Vs Global Step}
\end{figure}

\clearpage
\subsection{Run: 095 Crop Planning A2C Nonadaptive Random Weather Seed 0}
\begin{figure}[H]
    \centering
    \includegraphics[width=0.8\textwidth]{figures/per_run/095_crop_planning_a2c_nonadaptive_random_weather_seed_0__checkpoint_eval_curves.png}
    \caption{095 Crop Planning A2C Nonadaptive Random Weather Seed 0 - Checkpoint Eval Curves}
\end{figure}

\begin{figure}[H]
    \centering
    \includegraphics[width=0.8\textwidth]{figures/per_run/095_crop_planning_a2c_nonadaptive_random_weather_seed_0__diagnostics_panel.png}
    \caption{095 Crop Planning A2C Nonadaptive Random Weather Seed 0 - Diagnostics Panel}
\end{figure}

\begin{figure}[H]
    \centering
    \includegraphics[width=0.8\textwidth]{figures/per_run/095_crop_planning_a2c_nonadaptive_random_weather_seed_0__episode_length_vs_global_step.png}
    \caption{095 Crop Planning A2C Nonadaptive Random Weather Seed 0 - Episode Length Vs Global Step}
\end{figure}

\begin{figure}[H]
    \centering
    \includegraphics[width=0.8\textwidth]{figures/per_run/095_crop_planning_a2c_nonadaptive_random_weather_seed_0__primary_metric_vs_global_step.png}
    \caption{095 Crop Planning A2C Nonadaptive Random Weather Seed 0 - Primary Metric Vs Global Step}
\end{figure}

\begin{figure}[H]
    \centering
    \includegraphics[width=0.8\textwidth]{figures/per_run/095_crop_planning_a2c_nonadaptive_random_weather_seed_0__training_reward_vs_global_step.png}
    \caption{095 Crop Planning A2C Nonadaptive Random Weather Seed 0 - Training Reward Vs Global Step}
\end{figure}

\clearpage
\subsection{Run: 096 Crop Planning Hierarchical A2C Random Weather Seed 0 Profile Pakistan Baseline}
\begin{figure}[H]
    \centering
    \includegraphics[width=0.8\textwidth]{figures/per_run/096_crop_planning_hierarchical_a2c_random_weather_seed_0_profile_pakistan_baseline__diagnostics_panel.png}
    \caption{096 Crop Planning Hierarchical A2C Random Weather Seed 0 Profile Pakistan Baseline - Diagnostics Panel}
\end{figure}

\begin{figure}[H]
    \centering
    \includegraphics[width=0.8\textwidth]{figures/per_run/096_crop_planning_hierarchical_a2c_random_weather_seed_0_profile_pakistan_baseline__primary_metric_vs_global_step.png}
    \caption{096 Crop Planning Hierarchical A2C Random Weather Seed 0 Profile Pakistan Baseline - Primary Metric Vs Global Step}
\end{figure}

\clearpage
\subsection{Run: 097 Crop Planning A2C Adaptive Fixed Weather Seed 1}
\begin{figure}[H]
    \centering
    \includegraphics[width=0.8\textwidth]{figures/per_run/097_crop_planning_a2c_adaptive_fixed_weather_seed_1__checkpoint_eval_curves.png}
    \caption{097 Crop Planning A2C Adaptive Fixed Weather Seed 1 - Checkpoint Eval Curves}
\end{figure}

\begin{figure}[H]
    \centering
    \includegraphics[width=0.8\textwidth]{figures/per_run/097_crop_planning_a2c_adaptive_fixed_weather_seed_1__diagnostics_panel.png}
    \caption{097 Crop Planning A2C Adaptive Fixed Weather Seed 1 - Diagnostics Panel}
\end{figure}

\begin{figure}[H]
    \centering
    \includegraphics[width=0.8\textwidth]{figures/per_run/097_crop_planning_a2c_adaptive_fixed_weather_seed_1__episode_length_vs_global_step.png}
    \caption{097 Crop Planning A2C Adaptive Fixed Weather Seed 1 - Episode Length Vs Global Step}
\end{figure}

\begin{figure}[H]
    \centering
    \includegraphics[width=0.8\textwidth]{figures/per_run/097_crop_planning_a2c_adaptive_fixed_weather_seed_1__primary_metric_vs_global_step.png}
    \caption{097 Crop Planning A2C Adaptive Fixed Weather Seed 1 - Primary Metric Vs Global Step}
\end{figure}

\begin{figure}[H]
    \centering
    \includegraphics[width=0.8\textwidth]{figures/per_run/097_crop_planning_a2c_adaptive_fixed_weather_seed_1__training_reward_vs_global_step.png}
    \caption{097 Crop Planning A2C Adaptive Fixed Weather Seed 1 - Training Reward Vs Global Step}
\end{figure}

\clearpage
\subsection{Run: 098 Crop Planning A2C Nonadaptive Fixed Weather Seed 1}
\begin{figure}[H]
    \centering
    \includegraphics[width=0.8\textwidth]{figures/per_run/098_crop_planning_a2c_nonadaptive_fixed_weather_seed_1__checkpoint_eval_curves.png}
    \caption{098 Crop Planning A2C Nonadaptive Fixed Weather Seed 1 - Checkpoint Eval Curves}
\end{figure}

\begin{figure}[H]
    \centering
    \includegraphics[width=0.8\textwidth]{figures/per_run/098_crop_planning_a2c_nonadaptive_fixed_weather_seed_1__diagnostics_panel.png}
    \caption{098 Crop Planning A2C Nonadaptive Fixed Weather Seed 1 - Diagnostics Panel}
\end{figure}

\begin{figure}[H]
    \centering
    \includegraphics[width=0.8\textwidth]{figures/per_run/098_crop_planning_a2c_nonadaptive_fixed_weather_seed_1__episode_length_vs_global_step.png}
    \caption{098 Crop Planning A2C Nonadaptive Fixed Weather Seed 1 - Episode Length Vs Global Step}
\end{figure}

\begin{figure}[H]
    \centering
    \includegraphics[width=0.8\textwidth]{figures/per_run/098_crop_planning_a2c_nonadaptive_fixed_weather_seed_1__primary_metric_vs_global_step.png}
    \caption{098 Crop Planning A2C Nonadaptive Fixed Weather Seed 1 - Primary Metric Vs Global Step}
\end{figure}

\begin{figure}[H]
    \centering
    \includegraphics[width=0.8\textwidth]{figures/per_run/098_crop_planning_a2c_nonadaptive_fixed_weather_seed_1__training_reward_vs_global_step.png}
    \caption{098 Crop Planning A2C Nonadaptive Fixed Weather Seed 1 - Training Reward Vs Global Step}
\end{figure}

\clearpage
\subsection{Run: 099 Crop Planning Hierarchical A2C Fixed Weather Seed 1 Profile Pakistan Baseline}
\begin{figure}[H]
    \centering
    \includegraphics[width=0.8\textwidth]{figures/per_run/099_crop_planning_hierarchical_a2c_fixed_weather_seed_1_profile_pakistan_baseline__diagnostics_panel.png}
    \caption{099 Crop Planning Hierarchical A2C Fixed Weather Seed 1 Profile Pakistan Baseline - Diagnostics Panel}
\end{figure}

\begin{figure}[H]
    \centering
    \includegraphics[width=0.8\textwidth]{figures/per_run/099_crop_planning_hierarchical_a2c_fixed_weather_seed_1_profile_pakistan_baseline__primary_metric_vs_global_step.png}
    \caption{099 Crop Planning Hierarchical A2C Fixed Weather Seed 1 Profile Pakistan Baseline - Primary Metric Vs Global Step}
\end{figure}

\clearpage
\subsection{Run: 100 Crop Planning A2C Adaptive Random Weather Seed 1}
\begin{figure}[H]
    \centering
    \includegraphics[width=0.8\textwidth]{figures/per_run/100_crop_planning_a2c_adaptive_random_weather_seed_1__checkpoint_eval_curves.png}
    \caption{100 Crop Planning A2C Adaptive Random Weather Seed 1 - Checkpoint Eval Curves}
\end{figure}

\begin{figure}[H]
    \centering
    \includegraphics[width=0.8\textwidth]{figures/per_run/100_crop_planning_a2c_adaptive_random_weather_seed_1__diagnostics_panel.png}
    \caption{100 Crop Planning A2C Adaptive Random Weather Seed 1 - Diagnostics Panel}
\end{figure}

\begin{figure}[H]
    \centering
    \includegraphics[width=0.8\textwidth]{figures/per_run/100_crop_planning_a2c_adaptive_random_weather_seed_1__episode_length_vs_global_step.png}
    \caption{100 Crop Planning A2C Adaptive Random Weather Seed 1 - Episode Length Vs Global Step}
\end{figure}

\begin{figure}[H]
    \centering
    \includegraphics[width=0.8\textwidth]{figures/per_run/100_crop_planning_a2c_adaptive_random_weather_seed_1__primary_metric_vs_global_step.png}
    \caption{100 Crop Planning A2C Adaptive Random Weather Seed 1 - Primary Metric Vs Global Step}
\end{figure}

\begin{figure}[H]
    \centering
    \includegraphics[width=0.8\textwidth]{figures/per_run/100_crop_planning_a2c_adaptive_random_weather_seed_1__training_reward_vs_global_step.png}
    \caption{100 Crop Planning A2C Adaptive Random Weather Seed 1 - Training Reward Vs Global Step}
\end{figure}

\clearpage
\subsection{Run: 101 Crop Planning A2C Nonadaptive Random Weather Seed 1}
\begin{figure}[H]
    \centering
    \includegraphics[width=0.8\textwidth]{figures/per_run/101_crop_planning_a2c_nonadaptive_random_weather_seed_1__checkpoint_eval_curves.png}
    \caption{101 Crop Planning A2C Nonadaptive Random Weather Seed 1 - Checkpoint Eval Curves}
\end{figure}

\begin{figure}[H]
    \centering
    \includegraphics[width=0.8\textwidth]{figures/per_run/101_crop_planning_a2c_nonadaptive_random_weather_seed_1__diagnostics_panel.png}
    \caption{101 Crop Planning A2C Nonadaptive Random Weather Seed 1 - Diagnostics Panel}
\end{figure}

\begin{figure}[H]
    \centering
    \includegraphics[width=0.8\textwidth]{figures/per_run/101_crop_planning_a2c_nonadaptive_random_weather_seed_1__episode_length_vs_global_step.png}
    \caption{101 Crop Planning A2C Nonadaptive Random Weather Seed 1 - Episode Length Vs Global Step}
\end{figure}

\begin{figure}[H]
    \centering
    \includegraphics[width=0.8\textwidth]{figures/per_run/101_crop_planning_a2c_nonadaptive_random_weather_seed_1__primary_metric_vs_global_step.png}
    \caption{101 Crop Planning A2C Nonadaptive Random Weather Seed 1 - Primary Metric Vs Global Step}
\end{figure}

\begin{figure}[H]
    \centering
    \includegraphics[width=0.8\textwidth]{figures/per_run/101_crop_planning_a2c_nonadaptive_random_weather_seed_1__training_reward_vs_global_step.png}
    \caption{101 Crop Planning A2C Nonadaptive Random Weather Seed 1 - Training Reward Vs Global Step}
\end{figure}

\clearpage
\subsection{Run: 102 Crop Planning Hierarchical A2C Random Weather Seed 1 Profile Pakistan Baseline}
\begin{figure}[H]
    \centering
    \includegraphics[width=0.8\textwidth]{figures/per_run/102_crop_planning_hierarchical_a2c_random_weather_seed_1_profile_pakistan_baseline__diagnostics_panel.png}
    \caption{102 Crop Planning Hierarchical A2C Random Weather Seed 1 Profile Pakistan Baseline - Diagnostics Panel}
\end{figure}

\begin{figure}[H]
    \centering
    \includegraphics[width=0.8\textwidth]{figures/per_run/102_crop_planning_hierarchical_a2c_random_weather_seed_1_profile_pakistan_baseline__primary_metric_vs_global_step.png}
    \caption{102 Crop Planning Hierarchical A2C Random Weather Seed 1 Profile Pakistan Baseline - Primary Metric Vs Global Step}
\end{figure}

\clearpage
\subsection{Run: 103 Crop Planning A2C Adaptive Fixed Weather Seed 2}
\begin{figure}[H]
    \centering
    \includegraphics[width=0.8\textwidth]{figures/per_run/103_crop_planning_a2c_adaptive_fixed_weather_seed_2__checkpoint_eval_curves.png}
    \caption{103 Crop Planning A2C Adaptive Fixed Weather Seed 2 - Checkpoint Eval Curves}
\end{figure}

\begin{figure}[H]
    \centering
    \includegraphics[width=0.8\textwidth]{figures/per_run/103_crop_planning_a2c_adaptive_fixed_weather_seed_2__diagnostics_panel.png}
    \caption{103 Crop Planning A2C Adaptive Fixed Weather Seed 2 - Diagnostics Panel}
\end{figure}

\begin{figure}[H]
    \centering
    \includegraphics[width=0.8\textwidth]{figures/per_run/103_crop_planning_a2c_adaptive_fixed_weather_seed_2__episode_length_vs_global_step.png}
    \caption{103 Crop Planning A2C Adaptive Fixed Weather Seed 2 - Episode Length Vs Global Step}
\end{figure}

\begin{figure}[H]
    \centering
    \includegraphics[width=0.8\textwidth]{figures/per_run/103_crop_planning_a2c_adaptive_fixed_weather_seed_2__primary_metric_vs_global_step.png}
    \caption{103 Crop Planning A2C Adaptive Fixed Weather Seed 2 - Primary Metric Vs Global Step}
\end{figure}

\begin{figure}[H]
    \centering
    \includegraphics[width=0.8\textwidth]{figures/per_run/103_crop_planning_a2c_adaptive_fixed_weather_seed_2__training_reward_vs_global_step.png}
    \caption{103 Crop Planning A2C Adaptive Fixed Weather Seed 2 - Training Reward Vs Global Step}
\end{figure}

\clearpage
\subsection{Run: 104 Crop Planning A2C Nonadaptive Fixed Weather Seed 2}
\begin{figure}[H]
    \centering
    \includegraphics[width=0.8\textwidth]{figures/per_run/104_crop_planning_a2c_nonadaptive_fixed_weather_seed_2__checkpoint_eval_curves.png}
    \caption{104 Crop Planning A2C Nonadaptive Fixed Weather Seed 2 - Checkpoint Eval Curves}
\end{figure}

\begin{figure}[H]
    \centering
    \includegraphics[width=0.8\textwidth]{figures/per_run/104_crop_planning_a2c_nonadaptive_fixed_weather_seed_2__diagnostics_panel.png}
    \caption{104 Crop Planning A2C Nonadaptive Fixed Weather Seed 2 - Diagnostics Panel}
\end{figure}

\begin{figure}[H]
    \centering
    \includegraphics[width=0.8\textwidth]{figures/per_run/104_crop_planning_a2c_nonadaptive_fixed_weather_seed_2__episode_length_vs_global_step.png}
    \caption{104 Crop Planning A2C Nonadaptive Fixed Weather Seed 2 - Episode Length Vs Global Step}
\end{figure}

\begin{figure}[H]
    \centering
    \includegraphics[width=0.8\textwidth]{figures/per_run/104_crop_planning_a2c_nonadaptive_fixed_weather_seed_2__primary_metric_vs_global_step.png}
    \caption{104 Crop Planning A2C Nonadaptive Fixed Weather Seed 2 - Primary Metric Vs Global Step}
\end{figure}

\begin{figure}[H]
    \centering
    \includegraphics[width=0.8\textwidth]{figures/per_run/104_crop_planning_a2c_nonadaptive_fixed_weather_seed_2__training_reward_vs_global_step.png}
    \caption{104 Crop Planning A2C Nonadaptive Fixed Weather Seed 2 - Training Reward Vs Global Step}
\end{figure}

\clearpage
\subsection{Run: 105 Crop Planning Hierarchical A2C Fixed Weather Seed 2 Profile Pakistan Baseline}
\begin{figure}[H]
    \centering
    \includegraphics[width=0.8\textwidth]{figures/per_run/105_crop_planning_hierarchical_a2c_fixed_weather_seed_2_profile_pakistan_baseline__diagnostics_panel.png}
    \caption{105 Crop Planning Hierarchical A2C Fixed Weather Seed 2 Profile Pakistan Baseline - Diagnostics Panel}
\end{figure}

\begin{figure}[H]
    \centering
    \includegraphics[width=0.8\textwidth]{figures/per_run/105_crop_planning_hierarchical_a2c_fixed_weather_seed_2_profile_pakistan_baseline__primary_metric_vs_global_step.png}
    \caption{105 Crop Planning Hierarchical A2C Fixed Weather Seed 2 Profile Pakistan Baseline - Primary Metric Vs Global Step}
\end{figure}

\clearpage
\subsection{Run: 106 Crop Planning A2C Adaptive Random Weather Seed 2}
\begin{figure}[H]
    \centering
    \includegraphics[width=0.8\textwidth]{figures/per_run/106_crop_planning_a2c_adaptive_random_weather_seed_2__checkpoint_eval_curves.png}
    \caption{106 Crop Planning A2C Adaptive Random Weather Seed 2 - Checkpoint Eval Curves}
\end{figure}

\begin{figure}[H]
    \centering
    \includegraphics[width=0.8\textwidth]{figures/per_run/106_crop_planning_a2c_adaptive_random_weather_seed_2__diagnostics_panel.png}
    \caption{106 Crop Planning A2C Adaptive Random Weather Seed 2 - Diagnostics Panel}
\end{figure}

\begin{figure}[H]
    \centering
    \includegraphics[width=0.8\textwidth]{figures/per_run/106_crop_planning_a2c_adaptive_random_weather_seed_2__episode_length_vs_global_step.png}
    \caption{106 Crop Planning A2C Adaptive Random Weather Seed 2 - Episode Length Vs Global Step}
\end{figure}

\begin{figure}[H]
    \centering
    \includegraphics[width=0.8\textwidth]{figures/per_run/106_crop_planning_a2c_adaptive_random_weather_seed_2__primary_metric_vs_global_step.png}
    \caption{106 Crop Planning A2C Adaptive Random Weather Seed 2 - Primary Metric Vs Global Step}
\end{figure}

\begin{figure}[H]
    \centering
    \includegraphics[width=0.8\textwidth]{figures/per_run/106_crop_planning_a2c_adaptive_random_weather_seed_2__training_reward_vs_global_step.png}
    \caption{106 Crop Planning A2C Adaptive Random Weather Seed 2 - Training Reward Vs Global Step}
\end{figure}

\clearpage
\subsection{Run: 107 Crop Planning A2C Nonadaptive Random Weather Seed 2}
\begin{figure}[H]
    \centering
    \includegraphics[width=0.8\textwidth]{figures/per_run/107_crop_planning_a2c_nonadaptive_random_weather_seed_2__checkpoint_eval_curves.png}
    \caption{107 Crop Planning A2C Nonadaptive Random Weather Seed 2 - Checkpoint Eval Curves}
\end{figure}

\begin{figure}[H]
    \centering
    \includegraphics[width=0.8\textwidth]{figures/per_run/107_crop_planning_a2c_nonadaptive_random_weather_seed_2__diagnostics_panel.png}
    \caption{107 Crop Planning A2C Nonadaptive Random Weather Seed 2 - Diagnostics Panel}
\end{figure}

\begin{figure}[H]
    \centering
    \includegraphics[width=0.8\textwidth]{figures/per_run/107_crop_planning_a2c_nonadaptive_random_weather_seed_2__episode_length_vs_global_step.png}
    \caption{107 Crop Planning A2C Nonadaptive Random Weather Seed 2 - Episode Length Vs Global Step}
\end{figure}

\begin{figure}[H]
    \centering
    \includegraphics[width=0.8\textwidth]{figures/per_run/107_crop_planning_a2c_nonadaptive_random_weather_seed_2__primary_metric_vs_global_step.png}
    \caption{107 Crop Planning A2C Nonadaptive Random Weather Seed 2 - Primary Metric Vs Global Step}
\end{figure}

\begin{figure}[H]
    \centering
    \includegraphics[width=0.8\textwidth]{figures/per_run/107_crop_planning_a2c_nonadaptive_random_weather_seed_2__training_reward_vs_global_step.png}
    \caption{107 Crop Planning A2C Nonadaptive Random Weather Seed 2 - Training Reward Vs Global Step}
\end{figure}

\clearpage
\subsection{Run: 108 Crop Planning Hierarchical A2C Random Weather Seed 2 Profile Pakistan Baseline}
\begin{figure}[H]
    \centering
    \includegraphics[width=0.8\textwidth]{figures/per_run/108_crop_planning_hierarchical_a2c_random_weather_seed_2_profile_pakistan_baseline__diagnostics_panel.png}
    \caption{108 Crop Planning Hierarchical A2C Random Weather Seed 2 Profile Pakistan Baseline - Diagnostics Panel}
\end{figure}

\begin{figure}[H]
    \centering
    \includegraphics[width=0.8\textwidth]{figures/per_run/108_crop_planning_hierarchical_a2c_random_weather_seed_2_profile_pakistan_baseline__primary_metric_vs_global_step.png}
    \caption{108 Crop Planning Hierarchical A2C Random Weather Seed 2 Profile Pakistan Baseline - Primary Metric Vs Global Step}
\end{figure}

\clearpage
\subsection{Run: 109 Fertilization Dqn Ablation Adaptive Fixed Weather Years 5000 Seed 0}
\begin{figure}[H]
    \centering
    \includegraphics[width=0.8\textwidth]{figures/per_run/109_fertilization_dqn_ablation_adaptive_fixed_weather_years_5000_seed_0__diagnostics_panel.png}
    \caption{109 Fertilization Dqn Ablation Adaptive Fixed Weather Years 5000 Seed 0 - Diagnostics Panel}
\end{figure}

\begin{figure}[H]
    \centering
    \includegraphics[width=0.8\textwidth]{figures/per_run/109_fertilization_dqn_ablation_adaptive_fixed_weather_years_5000_seed_0__episode_length_vs_global_step.png}
    \caption{109 Fertilization Dqn Ablation Adaptive Fixed Weather Years 5000 Seed 0 - Episode Length Vs Global Step}
\end{figure}

\begin{figure}[H]
    \centering
    \includegraphics[width=0.8\textwidth]{figures/per_run/109_fertilization_dqn_ablation_adaptive_fixed_weather_years_5000_seed_0__primary_metric_vs_global_step.png}
    \caption{109 Fertilization Dqn Ablation Adaptive Fixed Weather Years 5000 Seed 0 - Primary Metric Vs Global Step}
\end{figure}

\begin{figure}[H]
    \centering
    \includegraphics[width=0.8\textwidth]{figures/per_run/109_fertilization_dqn_ablation_adaptive_fixed_weather_years_5000_seed_0__training_reward_vs_global_step.png}
    \caption{109 Fertilization Dqn Ablation Adaptive Fixed Weather Years 5000 Seed 0 - Training Reward Vs Global Step}
\end{figure}

\clearpage
\subsection{Run: 110 Fertilization Dqn Ablation Adaptive Random Weather Years 5000 Seed 0}
\begin{figure}[H]
    \centering
    \includegraphics[width=0.8\textwidth]{figures/per_run/110_fertilization_dqn_ablation_adaptive_random_weather_years_5000_seed_0__diagnostics_panel.png}
    \caption{110 Fertilization Dqn Ablation Adaptive Random Weather Years 5000 Seed 0 - Diagnostics Panel}
\end{figure}

\begin{figure}[H]
    \centering
    \includegraphics[width=0.8\textwidth]{figures/per_run/110_fertilization_dqn_ablation_adaptive_random_weather_years_5000_seed_0__episode_length_vs_global_step.png}
    \caption{110 Fertilization Dqn Ablation Adaptive Random Weather Years 5000 Seed 0 - Episode Length Vs Global Step}
\end{figure}

\begin{figure}[H]
    \centering
    \includegraphics[width=0.8\textwidth]{figures/per_run/110_fertilization_dqn_ablation_adaptive_random_weather_years_5000_seed_0__primary_metric_vs_global_step.png}
    \caption{110 Fertilization Dqn Ablation Adaptive Random Weather Years 5000 Seed 0 - Primary Metric Vs Global Step}
\end{figure}

\begin{figure}[H]
    \centering
    \includegraphics[width=0.8\textwidth]{figures/per_run/110_fertilization_dqn_ablation_adaptive_random_weather_years_5000_seed_0__training_reward_vs_global_step.png}
    \caption{110 Fertilization Dqn Ablation Adaptive Random Weather Years 5000 Seed 0 - Training Reward Vs Global Step}
\end{figure}

\clearpage
\subsection{Run: 111 Crop Planning Dqn Ablation Nonadaptive Fixed Weather Seed 0}
\begin{figure}[H]
    \centering
    \includegraphics[width=0.8\textwidth]{figures/per_run/111_crop_planning_dqn_ablation_nonadaptive_fixed_weather_seed_0__diagnostics_panel.png}
    \caption{111 Crop Planning Dqn Ablation Nonadaptive Fixed Weather Seed 0 - Diagnostics Panel}
\end{figure}

\begin{figure}[H]
    \centering
    \includegraphics[width=0.8\textwidth]{figures/per_run/111_crop_planning_dqn_ablation_nonadaptive_fixed_weather_seed_0__episode_length_vs_global_step.png}
    \caption{111 Crop Planning Dqn Ablation Nonadaptive Fixed Weather Seed 0 - Episode Length Vs Global Step}
\end{figure}

\begin{figure}[H]
    \centering
    \includegraphics[width=0.8\textwidth]{figures/per_run/111_crop_planning_dqn_ablation_nonadaptive_fixed_weather_seed_0__primary_metric_vs_global_step.png}
    \caption{111 Crop Planning Dqn Ablation Nonadaptive Fixed Weather Seed 0 - Primary Metric Vs Global Step}
\end{figure}

\begin{figure}[H]
    \centering
    \includegraphics[width=0.8\textwidth]{figures/per_run/111_crop_planning_dqn_ablation_nonadaptive_fixed_weather_seed_0__training_reward_vs_global_step.png}
    \caption{111 Crop Planning Dqn Ablation Nonadaptive Fixed Weather Seed 0 - Training Reward Vs Global Step}
\end{figure}

\clearpage
\subsection{Run: 112 Crop Planning Dqn Ablation Nonadaptive Random Weather Seed 0}
\begin{figure}[H]
    \centering
    \includegraphics[width=0.8\textwidth]{figures/per_run/112_crop_planning_dqn_ablation_nonadaptive_random_weather_seed_0__diagnostics_panel.png}
    \caption{112 Crop Planning Dqn Ablation Nonadaptive Random Weather Seed 0 - Diagnostics Panel}
\end{figure}

\begin{figure}[H]
    \centering
    \includegraphics[width=0.8\textwidth]{figures/per_run/112_crop_planning_dqn_ablation_nonadaptive_random_weather_seed_0__episode_length_vs_global_step.png}
    \caption{112 Crop Planning Dqn Ablation Nonadaptive Random Weather Seed 0 - Episode Length Vs Global Step}
\end{figure}

\begin{figure}[H]
    \centering
    \includegraphics[width=0.8\textwidth]{figures/per_run/112_crop_planning_dqn_ablation_nonadaptive_random_weather_seed_0__primary_metric_vs_global_step.png}
    \caption{112 Crop Planning Dqn Ablation Nonadaptive Random Weather Seed 0 - Primary Metric Vs Global Step}
\end{figure}

\begin{figure}[H]
    \centering
    \includegraphics[width=0.8\textwidth]{figures/per_run/112_crop_planning_dqn_ablation_nonadaptive_random_weather_seed_0__training_reward_vs_global_step.png}
    \caption{112 Crop Planning Dqn Ablation Nonadaptive Random Weather Seed 0 - Training Reward Vs Global Step}
\end{figure}

\clearpage
